\chapter{Primary Statutory Sources (United States Code)}
\label{apx:primary-usc}

\USCSourceNote

The same \texttt{.tex} bodies are also pulled inline where the narrative cites each provision; this appendix collects them for uninterrupted reading. Git diffs used in Section~\ref{sec:statute_evolution_usc} are stored as \texttt{.diff} files under \texttt{Paper/usc\_snippets/} and rebuilt by \texttt{make usc-diffs} (which runs \texttt{tools/usc\_diff.sh} with no arguments).

\section*{Title 18 --- Firearms and related}
\addcontentsline{toc}{section}{Title 18 --- Firearms and related}

\subsection*{§922(g) (excerpt)}
\label{apx:usc:18:922g}
\input{usc_snippets/usc_18_922g}

\subsection*{§922(o) (excerpt)}
\label{apx:usc:18:922o}
\input{usc_snippets/usc_18_922o}

\subsection*{§926B (LEOSA)}
\label{apx:usc:18:926b}
\input{usc_snippets/usc_18_926b}

\subsection*{§521 (criminal street gangs)}
\label{apx:usc:18:521}
\input{usc_snippets/usc_18_521}

\section*{Title 18 --- Peonage, slavery, and trafficking in persons}
\addcontentsline{toc}{section}{Title 18 --- Peonage, slavery, and trafficking in persons}

\subsection*{§1581 (Peonage)}
\label{apx:usc:18:1581}
\input{usc_snippets/usc_18_1581}

\subsection*{§1589 (Forced labor)}
\label{apx:usc:18:1589}
\input{usc_snippets/usc_18_1589}

\section*{Title 18 --- Civil rights (conspiracy and color of law)}
\addcontentsline{toc}{section}{Title 18 --- Civil rights (conspiracy and color of law)}

\subsection*{§241}
\label{apx:usc:18:241}
\input{usc_snippets/usc_18_241}

\subsection*{§242}
\label{apx:usc:18:242}
\input{usc_snippets/usc_18_242}

\section*{Title 18 --- Forfeiture}
\addcontentsline{toc}{section}{Title 18 --- Forfeiture}

\subsection*{§981(a) (Civil forfeiture)}
\label{apx:usc:18:981a}
\input{usc_snippets/usc_18_981a}

\section*{Title 18 --- Prison-made goods}
\addcontentsline{toc}{section}{Title 18 --- Prison-made goods}

\subsection*{§1761 (Transportation or importation)}
\label{apx:usc:18:1761}
\input{usc_snippets/usc_18_1761}

\section*{Title 18 --- Sentencing}
\addcontentsline{toc}{section}{Title 18 --- Sentencing}

\subsection*{§3559(c) (Three-strikes mandatory life imprisonment)}
\label{apx:usc:18:3559c}
\input{usc_snippets/usc_18_3559c}

\section*{Title 18 --- Wiretap and electronic surveillance}
\addcontentsline{toc}{section}{Title 18 --- Wiretap and electronic surveillance}

\subsection*{§2511 (Interception prohibition)}
\label{apx:usc:18:2511}
\input{usc_snippets/usc_18_2511}

\section*{Title 8 --- Aliens and nationality}
\addcontentsline{toc}{section}{Title 8 --- Aliens and nationality}

\subsection*{§§262--297 (Exclusion of Chinese --- Repealed)}
\label{apx:usc:8:ch7}
\input{usc_snippets/usc_8_ch7}

\section*{Title 21 --- Drug abuse prevention and control}
\addcontentsline{toc}{section}{Title 21 --- Drug abuse prevention and control}

\subsection*{§812(b) (Controlled Substances scheduling criteria)}
\label{apx:usc:21:812b}
\input{usc_snippets/usc_21_812b}

\subsection*{§841(a) (Prohibited acts)}
\label{apx:usc:21:841a}
\input{usc_snippets/usc_21_841a}

\section*{Title 42 --- Public health and welfare}
\addcontentsline{toc}{section}{Title 42 --- Public health and welfare}

\subsection*{§2000d (Title VI operative text)}
\label{apx:usc:42:2000d}
\input{usc_snippets/usc_42_2000d}

\section*{Title 50 --- War and national defense}
\addcontentsline{toc}{section}{Title 50 --- War and national defense}

\subsection*{§1801 (FISA definitions, excerpt)}
\label{apx:usc:50:1801abc}
\input{usc_snippets/usc_50_1801abc}

\section*{Title 52 --- Voting and elections}
\addcontentsline{toc}{section}{Title 52 --- Voting and elections}

\subsection*{§10301 (VRA enforcement)}
\label{apx:usc:52:10301}
\input{usc_snippets/usc_52_10301}

\chapter{Equation Registry and Era-Level Calibration}
\label{sec:equation_registry}

\section{Canonical Equation: The Predatory Min-Max Function}

The complete formal definition of the governing optimization algorithm is reproduced here as the canonical reference. All component equations in the manuscript derive from or refer back to this canonical equation.

\begin{definition}[The Predatory Min-Max Function --- Canonical Entry]
\textbf{Full Set-Theoretic Hierarchy:}\footnote{Registry restatement of eq.~\eqref{eq:10.4-predatory-minmax-full-definition} (Chapter~10).} $E \;\subset\; P_{\text{uppet}} \;\subset\; I_{\text{buffer}} \;\subset\; I$, $O_{\text{racialized}} \cap I = \emptyset$.\\
\textbf{Kernel Objective:}\footnote{Registry restatement of eq.~\eqref{eq:1.5-kernel-optimization} (Chapter~1) and eq.~\eqref{eq:10.5-kernel-objective-canonical} (Chapter~10). Empirical validation: see the Antebellum Kernel Optimization case study (p.~\pageref{cs:kernel_antebellum}) and the Piketty-Saez-Zucman top-decile wealth concentration analysis. Tier~3 (structural); see the Empirical Methodology chapter (p.~\pageref{ch:empirical_methodology}).}
\begin{equation}\label{eq:19.1-kernel-objective-registry-entry}
\max_{\{P_i\}}\; \mathcal{E}(t) \quad \text{subject to} \quad M_{\text{eff}}(t) = M(t) - \lambda\,\Phi_{\text{load}}(t) < \tau
\end{equation}
where the effective resistance variable is:\footnote{Registry restatement of eq.~\eqref{eq:10.6-effective-resistance-variable} (Chapter~10).}
\begin{equation}\label{eq:19.1a-effective-resistance-registry-entry}
M_{\text{eff}}(t) = M(t) - \lambda\,\Phi_{\text{load}}(t)
\end{equation}
\textbf{Suppression Envelope:}\footnote{Registry restatement of eq.~\eqref{eq:1.8-suppression-envelope} (Chapter~1) and eq.~\eqref{eq:10.7-suppression-envelope-canonical} (Chapter~10). Empirical validation: see the Post-1965 Backlash Wave case study (p.~\pageref{cs:backlash_wave}) and the ANES cross-racial solidarity proxy analysis. Tier~3 (structural); see the Empirical Methodology chapter (p.~\pageref{ch:empirical_methodology}).}
\begin{equation}\label{eq:19.2-suppression-envelope-registry-entry}
\Sigma_{\text{sup}}(t) = \psi_s(t) + \psi_m(t) + R(t) + \Phi_{\text{load}}(t)
\end{equation}
\textbf{Crash Condition:}\footnote{Registry restatement of eq.~\eqref{eq:1.9-crash-condition-derivative} (Chapter~1) and eq.~\eqref{eq:10.8-crash-condition-canonical} (Chapter~10). Empirical validation: Bacon's Rebellion (1676) and the Haitian Revolution (1791--1804) are the two documented cases where $M_{\text{eff}}(t) > \tau$ produced systemic kernel failure. Tier~3 (structural); see the Empirical Methodology chapter (p.~\pageref{ch:empirical_methodology}).}
\begin{equation}\label{eq:19.3-crash-condition-registry-entry}
\frac{dM}{dt} > \frac{d\Sigma_{\text{sup}}}{dt} \quad \Longleftrightarrow \quad M_{\text{eff}}(t) > \tau
\end{equation}
\textbf{Tier Benefit Ordering:}\footnote{Registry restatement of eq.~\eqref{eq:10.9-tier-benefit-ordering-canonical} (Chapter~10). Empirical validation: Piketty-Saez-Zucman top-decile wealth concentration, Federal Reserve SCF median household wealth by quintile, and BJS incarceration data confirming $E \cap O_{\text{final}} = \emptyset$. Tier~3 (ordinal); see the Empirical Methodology chapter (p.~\pageref{ch:empirical_methodology}).}
\begin{equation}\label{eq:19.4-tier-ordering-registry-entry}
\text{Benefit}(E) \gg \text{Benefit}(P_{\text{uppet}}) > \text{Benefit}(F_{\text{enforce}}) > \text{Benefit}(I_{\text{buffer}}) > \text{Benefit}(O_{\text{racialized}})
\end{equation}
\textbf{Concession Invariant:} $\forall\, R_i \in \mathcal{R}$ (non-kinetic reforms in the 1450--2026 dataset): $\Delta\max(R_i) = 0$\\
\textbf{Suppression Allocation:} $\psi_m(t) > 0$ only when $M(t) \to \tau$; sourced from increased $O_{\text{racialized}}$ extraction.\\
\textbf{Boundary Dynamics:} The reclassification operator $\mathcal{R}(x_i)$ defined in eq.~\eqref{eq:10.13-reclassification-operator} executes whenever $K(x_i) > K_{\text{tolerated}} \vee \mathrm{comply}(x_i) = 0$, moving $x_i$ from $I_{\text{buffer}}$ to $O_{\text{final}}$ with $\psi_s(x_i) \to 0$ instantaneously; empirical instantiations include Christiana~1851, Hamburg~1876, Mulford~1967, Ruby~Ridge~1992, Waco~1993, and Perdi~2026 ($\S$\ref{sec:ruby_ridge_waco}).\\
\textbf{Full definition and variable legend}: Section~\ref{sec:full_reveal}.
\end{definition}

\section{Electrodynamic Framework Registry}

The following equations formalize the Unified Electrodynamic Framework (UEF)
introduced in Chapter~\ref{ch:system_init}. They provide the hardware-layer
translation of the set-theoretic and wave-physics formalisms deployed
throughout the book. Their license to apply across substrates---the
reason the same field equations govern both electrons in copper and
humans in institutions---is the substrate-independence theorem
established in Section~\ref{subsec:substrate_independence}: humans were
executing these dynamics millennia before electromagnetism was
formalized, and the electromagnetic vocabulary was borrowed from the
prior social-power vocabulary. The registry
below is therefore the formal apparatus the physicists left behind,
applied to the system it was originally about.

\begin{definition}[Unified Electrodynamic Framework --- Canonical Entry]
\textbf{Root Equation --- The Unified Lorentz Force:}
\begin{equation}\label{eq:19.5-unified-lorentz-force-registry}
\vec{F}_{\text{total}} = \mathbf{Q}\,\vec{\mathcal{E}}_{\text{mat}} + \mathbf{Q}\left(\vec{v} \times \sum_{k=1}^{N} \rho_k \vec{B}_k\right)
\end{equation}
where $\mathbf{Q}$ is the intersectional charge multivector,
$\vec{\mathcal{E}}_{\text{mat}}$ is the material-wage electric field,
and $\sum \rho_k \vec{B}_k$ is the superposition of $N$ cultural magnetic
fields. The cross product $\vec{v} \times \vec{B}$ is the mathematical
signature of the psychological wage: always orthogonal to velocity, so it
does zero work.\footnote{Registry restatement of eq.~\eqref{eq:0.1-unified-lorentz-force} (Chapter~0). Tier~3 (ordinal/structural); see the Empirical Methodology chapter (p.~\pageref{ch:empirical_methodology}).}

\textbf{Complex Suppression Allocation:}
\begin{equation}\label{eq:19.6-complex-wage-registry}
W = \psi_m + j\psi_s
\end{equation}
where $\psi_m \in \mathbb{R}$ is the material wage (real power) and
$j\psi_s \in j\mathbb{R}$ is the psychological wage (reactive power). In
polar form: $W = R\,e^{j\theta}$ with $R = |W| = \sqrt{\psi_m^2 + \psi_s^2}$.
The phase angle $\theta$ diagnoses the operating mode:
$\theta = 90°$ is the default (purely psychological);
$\theta > 90°$ is the fascism threshold (Quadrant~II).\footnote{Registry restatement of eq.~\eqref{eq:2.2b-complex-wage-def} (Chapter~1). Tier~3 (ordinal/structural).}

\textbf{Solidarity as Imaginary Cancellation:}
\begin{equation}\label{eq:19.7-solidarity-cancellation-registry}
W + W^* = 2\psi_m
\end{equation}
When the Buffer Class adopts the complex conjugate $W^* = \psi_m - j\psi_s$
and aligns with the Out-group, the imaginary status wage cancels entirely
and real material power doubles.\footnote{Registry restatement of eq.~\eqref{eq:2.2e-solidarity-addition} (Chapter~1). Tier~3 (ordinal/structural).}

\textbf{Systemic Resistance (Drude Model):}
\begin{equation}\label{eq:19.8-systemic-resistance-registry}
R = \frac{m \cdot L}{n \cdot q^2 \cdot \tau \cdot A}
\end{equation}
The Elite engineers high resistance for the Out-group by minimizing $\tau$
(mean free time between collisions). Stop-and-frisk, welfare audits, and
algorithmic flags force individuals to collide with the state before building
drift velocity.\footnote{Registry restatement of eq.~\eqref{eq:0.4-systemic-resistance} (Chapter~0). Tier~3 (ordinal/structural).}

\textbf{Enclosure Capacitance:}
\begin{equation}\label{eq:19.9-enclosure-capacitance-registry}
C = \frac{\epsilon A}{d}
\end{equation}
where $\epsilon$ is the legal dielectric (property rights, inheritance law,
FHA lending), $A$ is plate area (geographic scope), and $d$ is distance
(redlining, segregation).\footnote{Registry restatement of eq.~\eqref{eq:0.5-enclosure-capacitance} (Chapter~0). Tier~3 (ordinal/structural).}

\textbf{Inductive Kickback:}
\begin{equation}\label{eq:19.10-inductive-kickback-registry}
V = -L\,\frac{di}{dt}
\end{equation}
When reform attempts to break the oppressive current ($dt \to 0$), the
inductor of cultural inertia discharges stored energy as fascist backlash.\footnote{Registry restatement of eq.~\eqref{eq:0.6-inductive-kickback} (Chapter~0). Tier~3 (ordinal/structural).}

\textbf{AC Complex Power and Real/Reactive Decomposition:}
\begin{equation}\label{eq:19.11-ac-complex-power-registry}
S = V \cdot I^{*} = |V|\,|I|\,e^{j\theta} = P_{\text{real}} + jQ_{\text{reactive}}
\end{equation}
where $P_{\text{real}} = |V|\,|I|\cos\theta = \psi_m$ is the material
wage (real power) and $Q_{\text{reactive}} = |V|\,|I|\sin\theta = \psi_s$
is the psychological wage (reactive power). DC mode is the
$\theta = 0$ limit; fascism threshold is $\theta > 90°$.\footnote{Registry restatement of eqs.~\eqref{eq:0.8-ac-complex-power}--\eqref{eq:0.9-ac-real-reactive-decomp} (Chapter~0). Tier~3 (ordinal/structural).}

\textbf{N-Weighted Phase Current at the Elite Node:}
\begin{equation}\label{eq:19.12-n-weighted-phase-registry}
I_E(t) = \sum_{k=1}^{N} \rho_k\,A_k\,\cos\!\bigl(\omega t + \phi_k\bigr)
\end{equation}
The intersectional generalization of three-phase AC: $N$ extraction
axes weighted by $\rho_k$ (population/coupling) with phases $\phi_k$
tuned by the interference engine. Kirchhoff's law at the Elite node
enforces extraction-current conservation under axis dropout: suppression
of one $A_k$ is rebalanced by the others weighted by
$\rho_k / (\sum_{m \ne k} \rho_m)$.\footnote{Registry restatement of eqs.~\eqref{eq:0.10-n-weighted-phase-current}--\eqref{eq:0.11-elite-node-kirchhoff} (Chapter~0). Tier~3 (ordinal/structural).}

\textbf{Buffer-Class Work Theorem:}
\begin{equation}\label{eq:19.13-buffer-work-theorem-registry}
W_{I_{\text{buffer}}}^{\text{psych}}(t_0, t_1) = \int_{t_0}^{t_1} \mathbf{Q}\,(\vec{v} \times \vec{B}) \cdot \vec{v}\, d\tau \;\equiv\; 0
\end{equation}
The systemic work performed on the Buffer Class by the psychological
wage is identically zero over any interval, because $\vec{v} \times \vec{B}$
is orthogonal to $\vec{v}$ by construction. Formal proof of Du Bois's
``public and psychological wage'' \cite{dubois}.\footnote{Registry restatement of eq.~\eqref{eq:0.15-buffer-class-work-theorem} (Chapter~0). Tier~3 (structural/derivational).}

\textbf{Reparations Integral:}
\begin{equation}\label{eq:19.14-reparations-integral-registry}
\mathcal{R}(t_0, t_{\text{now}}) = \int_{t_0}^{t_{\text{now}}} \operatorname{Re}\!\Bigl[V_{\text{state}}(\tau)\, I_{O}^{*}(\tau)\Bigr]\, d\tau
\end{equation}
The time integral of real power flowing from the Out-group node to the
Elite node over the historical operating interval of the extraction
system. Lower limits are fixed by documented system-initialization
dates ($t_0 = 1444$ Portuguese, $t_0 = 1492$ Hispaniola, $t_0 = 1619$
US). Three independent peer-reviewed calibrations exist: Craemer's
compounded-labor estimate yields \$14--17 trillion \cite{craemer2015slavery};
Darity \& Mullen's wealth-gap-closure estimate yields \$10--12 trillion
\cite{darity_mullen}; Coates' housing/education/credit
survey provides the public articulation \cite{coates_case_for_reparations}.
Tier 2 (multiple peer-reviewed calibrations of the same definite integral).\footnote{Registry restatement of eq.~\eqref{eq:0.16-reparations-integral} (Chapter~0). Tier~2 (multiple independent peer-reviewed calibrations).}

\textbf{Augmented Lagrangian Root Equation:}
\begin{equation}\label{eq:lag-root-registry}
\mathcal{L}^*(q,\dot{q},t) = T(\dot{q},t) - V(q,t) - \mathcal{D}(q,\dot{q},t) + \lambda\bigl(\tau - M_{\mathrm{eff}}(t)\bigr)
\end{equation}
where $T$ is the mobilized social capacity (kinetic energy of the Out-group and Buffer class), $V$ is the engineered constraint landscape (systemic diodes, redlining boundaries, Lyapunov ceilings), $\mathcal{D}$ is the Rayleigh dissipation function (bureaucratic friction, policing drag, administrative heat loss), and $\lambda$ is the marginal cost of suppression (the Complex Wage $W = \psi_m + j\psi_s$). The constraint $M_{\mathrm{eff}}(t) < \tau$ enforces the rebellion threshold.\footnote{Registry restatement of eq.~\eqref{eq:lag-augmented} (Chapter~\ref{ch:lagrangian}). Tier~3 (ordinal/structural); the augmented Lagrangian formalizes the directional claim that the Elite optimize extraction subject to a rebellion-threshold constraint. See the Empirical Methodology chapter (p.~\pageref{ch:empirical_methodology}).}
\end{definition}

\section{Canonical Symbol Registry}

The manuscript uses a dual representation: kernel semantics (\texttt{Max/Min}) and dynamic state variables. The registry below is canonical.

\begin{longtable}{>{\raggedright}p{2.7cm} >{\raggedright}p{3.2cm} >{\raggedright}p{5.3cm} >{\raggedright\arraybackslash}p{3.3cm}}
\toprule
\textbf{Symbol} & \textbf{Type} & \textbf{Definition} & \textbf{Primary Mapping} \\
\midrule
\endfirsthead
\toprule
\textbf{Symbol} & \textbf{Type} & \textbf{Definition} & \textbf{Primary Mapping} \\
\midrule
\endhead
$\mathcal{E}(t)$ & Dynamic objective & Time-indexed extraction output routed toward $E$. & Dynamic analogue of $\max$ \\
$M(t)$ & Dynamic state & Time-indexed class-coherence threat to kernel stability. & Dynamic analogue of $\min$ stress \\
$\tau$ & Threshold & Domestic collapse threshold for coherent class solidarity. & Failure boundary for $M(t)$ / $|S_{\text{total}}|$ \\
$\psi(t) = \psi_s(t)+\psi_m(t)$ & Control variable & Suppression allocation to $I_{\text{buffer}}$: $\psi_s$ (status wage, always active) $+$ $\psi_m$ (material wage, $\geq 0$; deployed when $M(t)\to\tau$; sourced from increased $O_{\text{racialized}}$ extraction). & Suppression envelope component \\
$R(t)$ & Control variable & Kinetic repression capacity via $F_{\text{enforce}}$. & Suppression envelope component \\
$\phi_{k,j}$ & Phase variable & Axis-level phase shift for axis $k$ and subgroup $j$. & Intersectional interference component \\
$\Phi_j$ & Phase aggregate & Compound subgroup phase shift: $\Phi_j=\sum_k \phi_{k,j}$. & Subgroup wave displacement \\
$\Phi_{\text{load}}(t)$ & Dispersion index & Aggregate phase-dispersion load across active subgroups. & Suppression envelope component \\
$S_{\text{total}}(t)$ & Wave state & Net class-solidarity signal under intersectional phase loading. & Collapse test against $\tau$ \\
$\rho_{\tau}$ & Diagnostic ratio & Proximity index $\rho_{\tau}=M(t)/\tau$. & Runtime-log risk indicator \\
$K(x_i)$ & Individual state & Kinetic capacity of individual $x_i$: armament $\times$ demonstrated willingness to deploy against $F_{\text{enforce}}$. & Reclassification input; $\S$\ref{sec:ruby_ridge_waco} \\
$K_{\text{tolerated}}$ & Threshold & $E$'s threshold for permitted Buffer-Class kinetic capacity under the current runtime; era-calibrated. & Reclassification threshold; $\S$\ref{sec:ruby_ridge_waco} \\
$\mathcal{R}(x_i)$ & Boundary operator & Reclassification operator: conditional assignment of $x_i$ to $I_{\text{buffer}}$ or $O_{\text{final}}$ based on $(K(x_i), K_{\text{tolerated}}, \mathrm{comply}(x_i))$; see eq.~\eqref{eq:10.13-reclassification-operator}. & Dynamic boundary function on eq.~\eqref{eq:10.12-o-final-construction} \\
% --- New variables: Temporal Enclosures, Strauss-Howe, Interface Swap (Eq. 10.15–10.20, 13.15–13.17) ---
$X_{\text{temporal}}$ & Mechanism set & Temporal Enclosure: the set of all sovereign debt instruments $D_{\text{sovereign}}$ whose issuance securitizes $L_{\text{future}}(O \cup I_{\text{buffer}})$ and transfers the discounted future labor surplus to $E$. See eq.~\eqref{eq:10.15-temporal-enclosure-definition}. & Post-kinetic extraction vector; replaces spatial/kinetic enclosures as primary mechanism in mature architectures. \\
$L_{\text{future}}$ & Resource variable & Aggregate future labor capacity of $O \cup I_{\text{buffer}}$: the productive output not yet performed, whose present value is securitized by $D_{\text{sovereign}}$. & Extraction target of $X_{\text{temporal}}$. \\
$D_{\text{sovereign}}$ & Stock variable & Sovereign debt outstanding: the total accumulated issuance of $X_{\text{temporal}}$ instruments; triggers eq.~\eqref{eq:13.15-interface-swap-trigger} when $D_{\text{sovereign}} > \bar{D}$. & Primary instrument of $X_{\text{temporal}}$; accumulation clock for Interface Swap. \\
$\delta_E(t)$ & Extraction rate & Annual extraction increment from financial repression: $\delta_E = \pi(t) - i(t) > 0$ when $r_{\text{real}} < 0$; the real purchasing power transferred per unit time from $I_{\text{buffer}}$/$O$ to $E$. See eq.~\eqref{eq:10.16-financial-repression-condition}. & Stealth extraction rate; operates without kinetic enforcement. \\
$\mathcal{E}_{X_{\text{temporal}}}(t)$ & Dynamic objective & Instantaneous temporal extraction rate: $\frac{dD_{\text{sovereign}}}{dt} \cdot \delta_E(t) \cdot \mathbf{1}[r < g]$. See eq.~\eqref{eq:10.17-temporal-extraction-rate}. & Component of $\mathcal{E}(t)$ active under $r < g$ condition; scales with sovereign debt accumulation. \\
$\mathcal{A}$ & Architecture variable & Extraction architecture: the specific legal, monetary, and institutional substrate on which the Predatory Min-Max Function executes at a given historical moment ($\mathcal{A}_{\text{old}}$, $\mathcal{A}_{\text{new}}$). See eq.~\eqref{eq:13.16-polymorphic-reboot-operator}. & Argument of the Polymorphic Reboot Operator. \\
$\Delta_E$ & Continuity invariant & Change in $E$'s asset base across an Interface Swap: $\Delta_E = \mathrm{Assets}(E, \mathcal{A}_{\text{new}}) - \mathrm{Assets}(E, \mathcal{A}_{\text{old}}) \geq 0$. See eq.~\eqref{eq:13.17-elite-asset-continuity-invariant}. & Hard constraint on every Polymorphic Interface Swap; $\Delta_E < 0$ falsifies the reboot operator. \\
$\mathcal{L}_{\text{Dalio}}$ & Cycle index & Dalio Empire Index: 8-factor composite measuring national power (education, innovation, competitiveness, output, trade share, military, financial center, reserve currency), normalized $[0,1]$. See eq.~\eqref{eq:10.20-dalio-extraction-bound}. & Macro-clock for the 250-year empire cycle; decline below $\mathcal{L}_{\text{Dalio}} \approx 0.75$ correlates with Demographic Paradox Limit activation. \\
\bottomrule
\end{longtable}

\textbf{Non-reuse constraint}: $E$ denotes the Elite class set only; $\mathcal{E}(t)$ denotes the extraction output function. $\psi$ denotes psychological wage only. All phase operations use $\phi/\Phi$ notation only. The symbol $\mathcal{R}$ is disambiguated by arity: $\mathcal{R}$ without arguments (as in the Concession Invariant, $\S$\ref{sec:equation_registry}) denotes the set of non-kinetic reforms in the 1450--2026 dataset; $\mathcal{R}(x_i)$ with an individual argument denotes the reclassification operator defined in equation~\eqref{eq:10.13-reclassification-operator}. The two occupy orthogonal mathematical categories (set vs.\ function) and cannot be confused in well-formed expressions.

\section{Era-Level Interference Calibration Matrix}

The following ranges are bounded internal calibration intervals on a normalized [0,1] scale. They calibrate comparability across eras at a normalized scale without asserting fine-grained precision.

\begin{longtable}{>{\raggedright}p{3.4cm} >{\raggedright}p{2.3cm} >{\raggedright}p{2.1cm} >{\raggedright}p{2.4cm} >{\raggedright\arraybackslash}p{4.9cm}}
\toprule
\textbf{Era / Runtime} & \textbf{$\Phi_{\text{load}}$} & \textbf{$\rho_{\tau}$} & \textbf{Confidence Tier} & \textbf{Evidence Pathway} \\
\midrule
\endfirsthead
\toprule
\textbf{Era / Runtime} & \textbf{$\Phi_{\text{load}}$} & \textbf{$\rho_{\tau}$} & \textbf{Confidence Tier} & \textbf{Evidence Pathway} \\
\midrule
\endhead
1486 (Lisbon) & [0.10, 0.20] & [0.70, 0.85] (calibrated against Bacon/Haiti anchors; see p.~\pageref{ch:empirical_methodology}) & Tier 2 & Chapter~\ref{ch:portugal} runtime + initial partition architecture + psychological wage initialization \cite{dubois}. \\
1676--1787 & [0.15, 0.55] & (1.00, 0.75]; \textbf{Anchor} ($\rho_\tau = 1.0$; see p.~\pageref{ch:empirical_methodology}) & Tier 1 & \textbf{Anchor} ($\rho_\tau = 1.0$): Bacon crash evidence + legal codification sequence + constitutional patch dynamics. \\
1704--1865 & [0.45, 0.60] & [0.45, 0.65] (calibrated against Bacon/Haiti anchors; see p.~\pageref{ch:empirical_methodology}) & Tier 2 & Enforcement-class consolidation + slave patrol and 13th Amendment continuity evidence. \\
Gendered axis & [0.35, 0.55] & [0.40, 0.60] (calibrated against Bacon/Haiti anchors; see p.~\pageref{ch:empirical_methodology}) & Tier 2 & Chapter~\ref{ch:gendered} runtime + coverture + reproductive criminalization evidence. \\
1870s--1960s & [0.55, 0.72] & [0.50, 0.68] (calibrated against Bacon/Haiti anchors; see p.~\pageref{ch:empirical_methodology}) & Tier 1 & Redlining containment + two-party scaling + capture dynamics \cite{rothstein, coatesreparations}. \\
1968--Present & [0.75, 0.92] & [0.80, 0.98] (calibrated against Bacon/Haiti anchors; see p.~\pageref{ch:empirical_methodology}) & Tier 1 & Variable Swap + carceral expansion + agenda-setting and elite-policy evidence \cite{gilens}. \\
2007--2012 (Great Recession Recompile) & [0.78, 0.90] & [0.82, 0.96] (calibrated against Bacon/Haiti anchors; see p.~\pageref{ch:empirical_methodology}) & Tier 1 & Housing crash + racialized wealth loss + bottom-half balance-sheet floor + asset recovery preserving $\Delta\max=0$ \cite{fed_dfa_2026,cbo_family_wealth_2024,pew_wealth_after_great_recession_2017,pfeffer_danziger_schoeni_2013}. \\
1992--1995 (Ruby Ridge / Waco / OKC) & [0.70, 0.85] & [0.72, 0.88] (calibrated against Bacon/Haiti anchors; see p.~\pageref{ch:empirical_methodology}) & Tier 1 & First live-action reclassification of Buffer-Class subjects under rewritten ROE; Patriot/militia misroute of kinetic response \cite{dojruby1994, oprruby2001, danforth2000, treasury1993waco, splc_patriot_timeline, pitcavage2001}. \\
Disarmament timeline & [0.60, 0.85] & [0.70, 0.95] (calibrated against Bacon/Haiti anchors; see p.~\pageref{ch:empirical_methodology}) & Tier 2 & Kinetic-control sequence as adjunct to interference suppression across periods. \\
Prescriptive stress test & [0.82, 0.95] & [0.92, 1.05] (calibrated against Bacon/Haiti anchors; see p.~\pageref{ch:empirical_methodology}) & Tier 2 & Reform paradox and concession dynamics under near-threshold conditions. \\
Global scaling & [0.50, 0.90] & [0.55, 0.95] (calibrated against Bacon/Haiti anchors; see p.~\pageref{ch:empirical_methodology}) & Tier 3 & Imperial containment variance by region and bloc coordination \cite{ohchr2026}. \\
Final output & [0.85, 0.97] & [0.95, 1.08] (calibrated against Bacon/Haiti anchors; see p.~\pageref{ch:empirical_methodology}) & Tier 2 & Book-level synthesis of terminal saturation and diminishing suppression returns. \\
\bottomrule
\end{longtable}

\textbf{Confidence tiers}: Tier 1 = multi-source quantitative alignment in-text; Tier 2 = mixed quantitative + structural diagnostics; Tier 3 = structurally supported but data-fragmented comparative estimate.

\chapter{Compiled Runtime Log}
\label{sec:runtime_log}

The following appendix concatenates all RUNTIME LOG boxes from the manuscript in chronological order, providing a single consolidated execution trace of the Predatory Min-Max Function across its five-century runtime. Each entry preserves the diagnostic format used in situ: system stress, capital output, interference state, variables deployed, and policy results.

\bigskip

\begin{tcolorbox}[colback=black!5, colframe=black!70, fonttitle=\bfseries\ttfamily, title=1. RUNTIME LOG: 1486 (LISBON{,} PORTUGAL) \normalfont\small--- Chapter \ref{ch:portugal}]
\ttfamily\small
\textbf{System Stress} ($\min$): HIGH --- Domestic labor shortage threatening agricultural capital.\\
\textbf{Capital} ($\max$): STAGNANT --- Sugar trade demands labor the current moral economy cannot supply.\\
\textbf{$\Phi_{\text{load}}$}: $[0.10, 0.20]$; $\rho_{\tau} \in [0.70, 0.85]$.\\
\textbf{Variables Deployed}: $E$, $O_{\text{racialized}}$, $I$.\\
\textbf{Result}: Extraction Algorithm initialized. \textit{Dum Diversas} (1452), \textit{Romanus Pontifex} (1455), \textit{Casa dos Escravos} (1486).
\end{tcolorbox}

\begin{tcolorbox}[colback=black!5, colframe=black!70, fonttitle=\bfseries\ttfamily, title=2. RUNTIME LOG: 1676--1787 (VIRGINIA $\rightarrow$ PHILADELPHIA) \normalfont\small--- Chapter \ref{ch:bacon}]
\ttfamily\small
\textbf{System Stress} ($\min$): \textcolor{red}{\textbf{CRITICAL}} --- Cross-racial labor solidarity detected. Jamestown burning.\\
\textbf{Capital} ($\max$): AT RISK --- Plantation economy destabilized by unified revolt.\\
\textbf{$\Phi_{\text{load}}$}: $[0.15, 0.55]$; $\rho_{\tau}>1.00$ at crash, then $[0.60, 0.75]$ post-patch.\\
\textbf{Variables Deployed}: $I_{\text{buffer}}$, $F_{\text{enforce}}^{\text{proto}}$, $P_{\text{uppet}}^{\text{v1.0}}$, $W = j\psi_s$ formalized.\\
\textbf{Result}: Virginia Slave Codes (1705), Three-Fifths Compromise (1787), constitutional front-end/back-end separation.
\end{tcolorbox}

\begin{tcolorbox}[colback=black!5, colframe=black!70, fonttitle=\bfseries\ttfamily, title=3. RUNTIME LOG: GENDERED AXIS (FOUNDATIONAL--PRESENT) \normalfont\small--- Chapter \ref{ch:gendered}]
\ttfamily\small
\textbf{System Stress} ($\min$): MODERATE --- $O_{\text{gendered}}$ isolated from economic and lethal autonomy.\\
\textbf{Capital} ($\max$): HIGH --- Reproductive labor optimized through coverture and coercive control.\\
\textbf{$\Phi_{\text{load}}$}: FRAGMENTED --- Intersection $O_{\text{racialized}} \cap O_{\text{gendered}}$ generates maximum $\alpha_{r,g}$.\\
\textbf{Variables Deployed}: $O_{\text{gendered}}$, $P_{\text{coverture}}$, $P_{\text{fetal\_personhood}}$.\\
\textbf{Result}: Coverture (1066--1974), Buck v.\ Bell (1927), Dobbs (2022) \cite{dobbs}.
\end{tcolorbox}

\begin{tcolorbox}[colback=black!5, colframe=black!70, fonttitle=\bfseries\ttfamily, title=4. RUNTIME LOG: 1704--1865 (AMERICAN SOUTH) \normalfont\small--- Chapter \ref{ch:enforcement}]
\ttfamily\small
\textbf{System Stress} ($\min$): MODERATE --- Buffer Class pacified by $j\psi_s$. Racial partition holding.\\
\textbf{Capital} ($\max$): EXPANDING --- Slave capitalism scaling. Human bodies as mortgageable assets.\\
\textbf{$\Phi_{\text{load}}$}: $[0.45, 0.60]$; $\rho_{\tau} \in [0.45, 0.65]$.\\
\textbf{Variables Deployed}: $F_{\text{enforce}}$, $QI$ (latent).\\
\textbf{Result}: Fugitive Slave Act (1850), 13th Amendment loophole (1865). $\max$ secured.
\end{tcolorbox}

\begin{tcolorbox}[colback=black!5, colframe=black!70, fonttitle=\bfseries\ttfamily, title=5. RUNTIME LOG: 1870s--1960s (RECONSTRUCTION $\rightarrow$ CIVIL RIGHTS) \normalfont\small--- Chapter \ref{ch:containment}]
\ttfamily\small
\textbf{System Stress} ($\min$): HIGH --- 13th Amendment reclassified $O_{\text{racialized}}$ as citizens. Civil Rights Movement breaching the interface.\\
\textbf{Capital} ($\max$): RESTRUCTURING --- Transitioning to indirect extraction via convict leasing, sharecropping, spatial containment.\\
\textbf{$\Phi_{\text{load}}$}: $[0.55, 0.72]$; $\rho_{\tau} \in [0.50, 0.68]$.\\
\textbf{Variables Deployed}: $P_{\text{uppet}}$ (scaled), $P_{\text{spatial}}$ (Redlining), Capture Variable, Tweedism Filter.\\
\textbf{Result}: HOLC Redlining (1934), Civil Rights Act (1964), Voting Rights Act (1965) --- interface dismantled, kernel preserved.
\end{tcolorbox}

\begin{tcolorbox}[colback=black!5, colframe=black!70, fonttitle=\bfseries\ttfamily, title=6. RUNTIME LOG: v5.0 (FRACTAL INTERFERENCE ENGINE) \normalfont\small--- Chapter \ref{ch:tweedism}]
\ttfamily\small
\textbf{Constraint}: Binary partition alone permits periodic cross-group coherence.\\
\textbf{Executed}: Compound phase shifting across intersecting axes (race/gender/orientation/ability).\\
\textbf{Diagnostic}: $\Phi_j = \sum_{k=1}^{K}\phi_{k,j}$; $S_{\text{total}}(t) = \sum A_j\sin(2\pi f_{\text{class}}t + \Phi_j)$.\\
\textbf{Result}: Destructive interference maximized. Kernel stable.
\end{tcolorbox}

\begin{tcolorbox}[colback=black!5, colframe=black!70, fonttitle=\bfseries\ttfamily, title=7. RUNTIME LOG: 1966 (SECURITY PATCH) \normalfont\small--- Chapter \ref{ch:recompile}]
\ttfamily\small
\textbf{System Stress} ($\min$): CRITICAL --- Civil Rights breaches legal kernel. IFAS Protocol detects 91\% breakthrough rate.\\
\textbf{Capital} ($\max$): THREATENED --- Legal desegregation disrupts extraction models.\\
\textbf{$\Phi_{\text{load}}$}: HIGH; $\rho_{\tau} \rightarrow 0.95$.\\
\textbf{Active Patch}: 1966 Moratorium, criminalizing LSD/Mescaline, recompiling Version 2.0 (Variable Swap).\\
\textbf{Result}: Epistemic Enclosure restored. Extraction migrated to race-neutral interface (War on Drugs).
\end{tcolorbox}

\begin{tcolorbox}[colback=black!5, colframe=black!70, fonttitle=\bfseries\ttfamily, title=8. RUNTIME LOG: 1968--1994 (THE RECOMPILE) \normalfont\small--- Chapter \ref{ch:recompile}]
\ttfamily\small
\textbf{System Stress} ($\min$): \textcolor{red}{\textbf{CRITICAL}} --- Demographic Paradox emerging. $I_{\text{buffer}}$ detecting contract breach.\\
\textbf{Capital} ($\max$): PEAK --- Carceral state fully industrialized. War on Drugs providing unlimited proxy criminalization.\\
\textbf{$\Phi_{\text{load}}$}: $[0.75, 0.92]$; $\rho_{\tau} \in [0.80, 0.98]$.\\
\textbf{Variables Deployed}: $P_{\text{criminal}}$, $P_{\text{lead}}$, $P_{\text{deindustrial}}$, $P_{\text{BrokenWindows}}$, Universal Latent Criminality.\\
\textbf{WARNING}: $\min$ VARIABLE FAILING. System entering terminal phase.
\end{tcolorbox}

\begin{tcolorbox}[colback=black!5, colframe=black!70, fonttitle=\bfseries\ttfamily, title=9. RUNTIME LOG: 2007--2012 (FINANCIAL RECOMPILE) \normalfont\small--- Chapter \ref{ch:full_algo}]
\ttfamily\small
\textbf{System Stress} ($\min$): HIGH --- household balance sheets collapsing; foreclosure shock destabilizing $O_{\text{racialized}}$ and lower $I_{\text{buffer}}$.\\
\textbf{Capital} ($\max$): RESTORED --- state liquidity and monetary intervention stabilize financial assets and institutional balance sheets.\\
\textbf{$\Phi_{\text{load}}$}: $[0.78, 0.90]$; $\rho_{\tau} \in [0.82, 0.96]$.\\
\textbf{Variables Deployed}: $P_{\text{debt}}$, $X_{\text{temporal}}$, foreclosure transfer, QE asset inflation.\\
\textbf{Result}: Bottom-tier wealth destroyed; racial wealth gap amplified; Elite asset base restored; $\Delta\max = 0$ preserved.
\end{tcolorbox}

\begin{tcolorbox}[colback=black!5, colframe=black!70, fonttitle=\bfseries\ttfamily, title=10. RUNTIME LOG: DISARMAMENT VARIABLE (FULL TIMELINE) \normalfont\small--- Chapter \ref{ch:kinetic}]
\ttfamily\small
\textbf{Origin}: Portuguese firearms superiority (1440s).\\
\textbf{Feedback Loop}: Arms-for-slaves trade creates exponential demand.\\
\textbf{$\Phi_{\text{load}}$}: $[0.60, 0.85]$; $\rho_{\tau} \in [0.70, 0.95]$.\\
\textbf{Disarmament Protocol}: Sullivan (1911) $\rightarrow$ NFA (1934) $\rightarrow$ Mulford (1967) $\rightarrow$ GCA (1968) $\rightarrow$ Hughes (1986) $\rightarrow$ Brady/AWB (1993--94) $\rightarrow$ Post-Bruen (2022--).\\
\textbf{Pattern}: Every federal gun control law correlates with $O_{\text{racialized}}$ approaching kinetic parity or $P_{\text{uppet}}$ being kinetically corrected.
\end{tcolorbox}

\begin{tcolorbox}[colback=black!5, colframe=black!70, fonttitle=\bfseries\ttfamily, title=11. RUNTIME LOG: PRESCRIPTIVE STRESS TEST \normalfont\small--- Chapter \ref{ch:contradiction}]
\ttfamily\small
\textbf{System Stress} ($\min$): \textcolor{red}{\textbf{FAILING}} --- $I_{\text{buffer}}$ awakening. Kinetic capacity of civilian population exceeds $F_{\text{enforce}}$.\\
\textbf{Capital} ($\max$): TERMINAL PHASE --- Extraction zone expanded into $I_{\text{buffer}}$.\\
\textbf{$\Phi_{\text{load}}$}: $[0.82, 0.95]$; $\rho_{\tau} \in [0.92, 1.05]$.\\
\textbf{WARNING}: A reform that reduces $\min$ while preserving $\max$ leaves the kernel in \textbf{maintenance} mode.
\end{tcolorbox}

\begin{tcolorbox}[colback=black!5, colframe=black!70, fonttitle=\bfseries\ttfamily, title=11. RUNTIME LOG: SCALING TO IMPERIAL ARCHITECTURE \normalfont\small--- Chapter \ref{ch:global}]
\ttfamily\small
\textbf{System Stress} ($\min$): Historically variable. Each crash temporarily spiked $\min$, forcing emergency patches.\\
\textbf{Capital} ($\max$): Monotonically increasing across five centuries. Scope expanding to global field.\\
\textbf{$\Phi_{\text{load}}$}: $[0.50, 0.90]$; $\rho_{\tau} \in [0.55, 0.95]$ (region-dependent).\\
\textbf{Executing}: Extending 5-tier hierarchy to international system. Testing Haitian Theorem against global containment.
\end{tcolorbox}

\begin{tcolorbox}[colback=black!5, colframe=black!70, fonttitle=\bfseries\ttfamily, title=12. RUNTIME LOG: FINAL OUTPUT \normalfont\small--- Chapters 11--12]
\ttfamily\small
\textbf{System Stress} ($\min$): Five centuries of runtime. $\Delta\max = 0$ for every non-kinetic reform $R_i$.\\
\textbf{Capital} ($\max$): $O$ expands ($O_{\text{final}} = \text{Everyone} \setminus E$). Global containment neutralizes peripheral liberation.\\
\textbf{$\Phi_{\text{load}}$}: $[0.85, 0.97]$; $\rho_{\tau} \in [0.95, 1.08]$. TERMINAL SATURATION.\\
\textbf{Variables (Complete)}: $E$/$E_{\text{global}}$, $P_{\text{uppet}}$/$P_{\text{uppet}}^{\text{global}}$, $F_{\text{enforce}}$/$F_{\text{enforce}}^{\text{global}}$, $I_{\text{buffer}}$/$I_{\text{buffer}}^{\text{global}}$, $O_{\text{racialized}}$/$O_{\text{global}}$.\\
\textbf{Status}: \texttt{return} $\vec{R}_{\text{acism}}$
\end{tcolorbox}

\chapter{Falsifiability Conditions for the Two Terminal Theorems}
\label{app:falsifiability}

This appendix specifies, in precise terms, the empirical conditions under which the Concession Theorem and the Haitian Theorem would be \textit{falsified}. The purpose is twofold. First, it demonstrates that both theorems are genuine empirical claims with specifiable defeat conditions. Second, it provides a reference standard against which critics can test proposed counter-examples, and against which future evidence can be assessed as the historical record develops past 2026.

\section*{I. Falsification Conditions for the Concession Theorem}

\textbf{The theorem's claim.} For every non-kinetic reform $R_i$ in the 1450--2026 dataset, $\Delta\max(R_i) = 0$: no non-kinetic reform has permanently reduced the Elite's extraction share as measured at the system apex.

\textbf{What falsification requires.} A genuine falsification of the Concession Theorem must satisfy \textit{all three} of the following conditions simultaneously:

\begin{enumerate}
    \item \textbf{Structural reduction in extraction share}: The reform must produce a measurable, permanent decline in the top-decile's share of total national wealth, with an absolute decrease in $E$'s extraction share sustained over a minimum of three decades. A reduction in the growth rate of inequality falls short of the required structural reduction.

    \item \textbf{No compensatory extraction elsewhere}: The reduction in the measured extraction share must not be offset by increased extraction through other mechanisms---capital flight to offshore jurisdictions, expansion of the extraction base to new populations, or substitution of financial extraction for labor extraction. The test is whether aggregate extraction by the apex ($E$) declined in real terms across all channels.

    \item \textbf{Sustained across threat cycles}: The reduction must persist through subsequent periods of low kinetic threat, demonstrating structural alteration of the extraction kernel. A gain that erodes once the threat subsides counts as evidence for the theorem.
\end{enumerate}

\textbf{The kinetic-precondition prediction.} The theorem additionally predicts the \textit{mechanism} by which concessions arise: reforms that reduce extraction load are offered when class resistance rises ($\min$ approaching $\tau$), their magnitude scales with the threat level, and their durability decays as the threat retreats. This prediction is testable on its own terms; the Great Compression arc---rising with the kinetic convergence of 1932--1945, decaying after 1973 as the threat environment dissipated---is its cleanest measured instance (Figure~\ref{fig:deltamax_invariant}). The presence of a kinetic precondition at adoption is evidence about the mechanism of concession; the three durability conditions above carry the falsification weight. A reform satisfying conditions 1--3 falsifies the theorem regardless of the threat conditions under which it was adopted.

\textbf{The current state of the record.} No reform in the 1450--2026 dataset satisfies all three conditions. The New Deal, the Civil Rights Act, OSHA, the Voting Rights Act, and comparable reforms in the dataset satisfy condition 1 partially and temporarily, but fail conditions 2 or 3 or both: the long-run wealth concentration data shows that $\Delta\max$ returned to zero within two to three decades in every case (Figure~\ref{fig:deltamax_invariant}), and each reform's adoption under elevated $\min$ conditions matches the kinetic-precondition prediction. The theorem remains unfalsified.

\textbf{What would falsify it.} A concrete example: if post-2026 evidence showed that the top 10\% wealth share in the United States had declined from approximately 70\% (2020 baseline) to below 60\% and remained there for thirty or more years following a non-kinetic reform (legislation, litigation, or electoral change), with no measurable compensatory extraction in other jurisdictions or asset classes, the Concession Theorem would require revision or scope restriction---whatever the threat environment at the moment of adoption.

\section*{II. Falsification Conditions for the Haitian Theorem}

\textbf{The theorem's claim.} For all non-kinetic interventions $R_i$ in the dataset: $\Delta\max(R_i) = 0$ (eq.~\ref{eq:12.5-haitian-theorem-nonkinetic}). There exists at least one kinetic action $K_j$ such that $\max(t_{\text{post}}) = 0$ locally (eq.~\ref{eq:12.6-haitian-theorem-kinetic}).

\textbf{What falsification requires.} A genuine falsification of the strong-form Haitian Theorem must produce a documented historical case in which \textit{all five} of the following conditions are satisfied:

\begin{enumerate}
    \item \textbf{Intra-national racial partition}: The society in question must have an installed, operative racial partition with a functioning $\psi$ mechanism---a buffer class holding status wages against a racially defined out-group. Systems without this architecture fall under the theorem's scope condition (Appendix~D, Section~I) and cannot falsify the strong form.

    \item \textbf{Absence of kinetic precondition}: The liberation must have been achieved without any credible kinetic threat operating in the background---no armed wing, no imminent insurrection, no defection within $F_{\text{enforce}}$, no external military force capable of enforcing a transition. If kinetic potential was present in any of these forms, the case resolves as the theorem's weak form (threat without discharge), which falls outside the falsification condition.

    \item \textbf{Kernel-level liberation}: The outcome must terminate the prior extraction kernel; desegregation of public facilities while the economic commanding heights remain in Elite ownership is an interface update. Kernel termination requires that the prior Elite's extraction capacity was structurally reduced to zero in the material apparatus---the plantation, the corporation, the financial system, or whatever constitutes the extraction mechanism---independently of any legal or political face change.

    \item \textbf{Durable outcome}: The liberation must persist through at least one full threat cycle without requiring new kinetic pressure to maintain it, which would demonstrate structural dismantlement of the kernel.

    \item \textbf{No Elite class substitution}: The liberation must not have resulted in a new extraction kernel installed by a replacement Elite---the Zimbabwe pattern (Appendix to the Boundary Conditions, Chapter~\ref{ch:contradiction}), in which kinetic liberation produces kernel transplant while leaving the extraction topology intact, does not falsify the theorem. Falsification requires permanent architectural dismantlement of the extraction kernel.
\end{enumerate}

\textbf{The current state of the record.} No case in the 1450--2026 dataset satisfies all five conditions. The boundary cases examined in Chapter~\ref{ch:contradiction}---South Africa, Scandinavia, India, Zimbabwe---each fail at least one condition; each failure illuminates the structural architecture of the theorem. The theorem remains unfalsified in strong form.

\textbf{What would falsify it.} A concrete example: if a documented historical case could be established in which a society with an operative intra-national racial partition (comparable to the American 1705--2026 architecture) achieved $\max(t_{\text{post}}) = 0$---the permanent dismantlement of the extraction kernel---through moral argument, electoral politics, litigation, or nonviolent civil disobedience, with no kinetic threat operating in any form, the Haitian Theorem's strong form would require revision.

\section*{III. Why These Conditions Have Not Been Met}

The two falsification specifications above share a common structure: they require permanent, kernel-level structural change in a society with an installed racial partition. The framework explains the empty counter-example set as a consequence of the architecture itself: the extraction kernel is a self-reinforcing closed-loop system whose Lyapunov stability (Section~\ref{sec:formal-containment-model}) absorbs every perturbation into a descent direction of the system's energy function. On this account, a non-kinetic intervention is a perturbation within the system's invariant set, and only a force that operates outside the system's control architecture---one that severs the spanning-tree topology itself---can produce the structural change that falsification requires. The explanation is itself a hypothesis, held accountable by the conditions above: it fails on the day a counter-example meets them.

The theorems are falsifiable: specifiable evidence can count against them. They remain unfalsified because the architectural constraint that generates their predictions is still operative. The Concession Theorem and the Haitian Theorem would be falsified by the existence of a historical counter-example meeting the specifications above. No such counter-example has been identified in the 1450--2026 dataset. The burden of proof falls on the critic to supply one.

\chapter{Geometric Algebra and the N-Dimensional Wage}
\label{apx:geometric_algebra}

The 2D complex number $W = \psi_m + j\psi_s$ models a single imaginary axis of oppression (e.g., race). Oppression is fractal in nature and operates across multiple simultaneous, orthogonal axes; the framework must therefore be extended to higher-dimensional algebraic structures.

This appendix is a speculative extension, in the same register as the photon model (Appendix~\ref{app:photon_model}). The algebra supplies a precise formal encoding of the directional claim that intersectional oppression compounds non-additively; it carries no calibrated measurement content, and Tier~3 status applies throughout. The empirical authority of the framework rests on the historical chapters and the tier-tagged case studies, with this material offered as a formal language for future development.

\section{The Quaternion Stage: Three Intersectional Axes}

For three primary imaginary axes (race, gender, sexuality), the standard complex number $a + bj$ is extended to a \emph{quaternion}:
\begin{equation}\label{eq:ga.quaternion}
W_Q = a + bi + cj + dk
\end{equation}
where $i, j, k$ are three distinct imaginary units satisfying:
\begin{equation}\label{eq:ga.quaternion-rules}
i^2 = j^2 = k^2 = ijk = -1
\end{equation}

\begin{center}
\begin{tabular}{lll}
\toprule
Component & EE Axis & Sociological Meaning \\
\midrule
$a$ & Real axis & Material Wage $\psi_m$ --- wealth, land, healthcare \\
$bi$ & Imaginary axis 1 & Racial Status Wage (e.g., whiteness) \\
$cj$ & Imaginary axis 2 & Gendered Status Wage (e.g., patriarchy/maleness) \\
$dk$ & Imaginary axis 3 & Cisnormative/Sexuality Status Wage \\
\bottomrule
\end{tabular}
\end{center}

\begin{theorem}[Multiplicative Intersection Compounding (Misogynoir)]
Because $i, j, k$ are mathematically orthogonal and multiply according to quaternion rules:
\begin{equation}\label{eq:ga.misogynoir}
i \times j = k, \quad j \times i = -k
\end{equation}
Under this algebra, intersectional oppression compounds multiplicatively. For a Black woman, the intersection of the racial wage axis ($i$) and the gendered wage axis ($j$) multiplies to create a distinct third vector $k$---a unique plane of systemic extraction requiring both axes for its full specification. This formalizes Kimberl\'e Crenshaw's theory of intersectionality as a quaternion multiplication law.\footnote{This theorem is classified as Tier~3 (ordinal/structural); the quaternion formalism provides a mathematical encoding of the directional claim that intersectional oppression is qualitatively distinct from the sum of its constituent axes. See the Empirical Methodology chapter (p.~\pageref{ch:empirical_methodology}).}
\end{theorem}

\section{The N-Dimensional Weighted Summation Equation}

Beyond three axes, the framework uses \emph{Geometric Algebra} (Clifford Algebra) over $N$ orthogonal imaginary bases $e_1, e_2, \ldots, e_N$. The full Complex Wage is:
\begin{equation}\label{eq:ga.ndim-allocation}
W = \psi_m + \sum_{k=1}^{N} j_k \left(\rho_k \cdot \psi_{s,k}\right)
\end{equation}
where:
\begin{itemize}
\item $\psi_m$: The Real Material Wage.
\item $j_k$: The $k$-th imaginary basis axis (e.g., $j_1$ = race, $j_2$ = gender, $j_3$ = ability, $j_4$ = sexuality, $j_5$ = neurology, $j_6$ = height/physicality, \ldots).
\item $\psi_{s,k}$: The baseline psychological wage on axis $k$.
\item $\rho_k$: The \textbf{Demographic Weight Coefficient}---the proportion of the population for whom axis $k$ is a primary axis of marginalization, scaled by institutional bandwidth devoted to that axis.
\end{itemize}

\begin{keyinsight}[The Transphobia Paradox (High Voltage, Low Current)]
The electrical analog of reactive power is $P = VI$ where $V$ is voltage and $I$ is current (population size). If $\rho_k$ is small (small population directly experiencing a given axis), but the Elite needs large reactive power $P_k$ from that axis to sustain the overall pacification of the Buffer Class, then:
\begin{equation}\label{eq:ga.transphobia-paradox}
P_k = V_k \cdot I_k \implies V_k = \frac{P_k}{\rho_k}
\end{equation}
As $\rho_k \to 0$, the required voltage $V_k \to \infty$. This mathematically predicts and explains the disproportionate legislative, media, and political bandwidth---hyper-inflated voltage---devoted to a statistically small demographic node (e.g., trans individuals). The Elite pumps astronomical voltage into a small-current node to generate sufficient reactive power to distract the Buffer Class from macro-level material extraction.\footnote{This derivation is classified as Tier~3 (ordinal/structural); it formalizes the directional claim that the Elite can achieve high reactive power from small demographic targets by intensifying the ideological voltage directed at them. See the Empirical Methodology chapter (p.~\pageref{ch:empirical_methodology}).}
\end{keyinsight}

\section{Bivectors and Intersectional Space}

In Geometric Algebra, multiplying two basis vectors produces a \emph{bivector}---a distinct two-dimensional plane:
\begin{equation}\label{eq:ga.bivector}
e_1 \wedge e_4 = e_{14}
\end{equation}
If $e_1$ is Race and $e_4$ is Disability, the bivector $e_{14}$ describes the unique extraction plane of a Black disabled person. This individual is located in $e_{14}$---a qualitatively distinct geometric zone where extraction is accelerated, enforcement violence is compounded, and systemic neglect is multiplicatively intensified. The bivector $e_{14}$ encodes an extraction plane qualitatively different from either axis alone.

\begin{theorem}[Fractal Consistency of Geometric Algebra]
The Geometric Algebra extension reproduces the same extraction geometry at the multi-axis level that the complex number $W = \psi_m + j\psi_s$ produces at the single-axis level. Specifically:
\begin{itemize}
\item At the single-axis level, Orthogonal Deflection is produced by the magnetic cross-product $q(\vec{v} \times \vec{B})$: the Buffer Class is deflected horizontally by the psychological wage field.
\item At the multi-axis level, the same deflection is produced by the superposition of $N$ cultural magnetic fields $\sum_{k=1}^{N} \rho_k \vec{B}_k$, each weighted by its demographic coefficient $\rho_k$.
\end{itemize}
The fractal is exact: the mathematics is the same mathematics operating at different dimensional resolutions.\footnote{This theorem is classified as Tier~3 (ordinal/structural). See the Empirical Methodology chapter (p.~\pageref{ch:empirical_methodology}).}
\end{theorem}

\chapter*{Empirical Validation Index}
\label{app:empirical_index}
\addcontentsline{toc}{chapter}{Empirical Validation Index}

% ── Empirical Validation Index ────────────────────────────────────────────
% Equations: 146 total | 146 fully specified | 0 with missing fields
% status=complete: 146 (abstract stats: N=146, M=146)
% Columns: label | case study location | tier (T1--T3 badge) | primary source | falsification
% Tier legend: 1=peer-reviewed quantitative  2=public dataset  3=ordinal/structural
% Bold label = Phase 3 headline equation
% ──────────────────────────────────────────────────────────────────────────

% empirical_index.tex — auto-generated by Paper/scripts/generate_index.py
% Do NOT edit manually. Regenerate with: make index
%
% Columns: equation label | case study location | confidence tier (badge) |
%          primary data source | falsification (truncated)
%
% Include in main document with:
%   % ── Empirical Validation Index ────────────────────────────────────────────
% Equations: 146 total | 146 fully specified | 0 with missing fields
% status=complete: 146 (abstract stats: N=146, M=146)
% Columns: label | case study location | tier (T1--T3 badge) | primary source | falsification
% Tier legend: 1=peer-reviewed quantitative  2=public dataset  3=ordinal/structural
% Bold label = Phase 3 headline equation
% ──────────────────────────────────────────────────────────────────────────

% empirical_index.tex — auto-generated by Paper/scripts/generate_index.py
% Do NOT edit manually. Regenerate with: make index
%
% Columns: equation label | case study location | confidence tier (badge) |
%          primary data source | falsification (truncated)
%
% Include in main document with:
%   % ── Empirical Validation Index ────────────────────────────────────────────
% Equations: 146 total | 146 fully specified | 0 with missing fields
% status=complete: 146 (abstract stats: N=146, M=146)
% Columns: label | case study location | tier (T1--T3 badge) | primary source | falsification
% Tier legend: 1=peer-reviewed quantitative  2=public dataset  3=ordinal/structural
% Bold label = Phase 3 headline equation
% ──────────────────────────────────────────────────────────────────────────

% empirical_index.tex — auto-generated by Paper/scripts/generate_index.py
% Do NOT edit manually. Regenerate with: make index
%
% Columns: equation label | case study location | confidence tier (badge) |
%          primary data source | falsification (truncated)
%
% Include in main document with:
%   \input{empirical_index}
%
% Requires: longtable, booktabs, xcolor, array (loaded in main preamble)

\begingroup
\small
\setlength{\LTpre}{6pt}
\setlength{\LTpost}{6pt}
% Equation labels are long unbreakable \texttt strings (up to 47 chars). Without
% explicit break opportunities they overrun the label column and print on top of
% the case-study column. \eqlab emits the label with zero-width discretionary
% breaks after ':', '-' and '.', inserted by the generator.
\providecommand{\eqlab}[1]{\texttt{#1}}
\begin{longtable}{@{}%
  >{\raggedright\arraybackslash}p{3.2cm}
  >{\raggedright\arraybackslash}p{2.5cm}
  >{\centering\arraybackslash}p{1.0cm}
  >{\raggedright\arraybackslash}p{3.3cm}
  >{\raggedright\arraybackslash}p{5.2cm}@{}}
\textbf{Equation label} &
\textbf{Case study location} &
\textbf{Tier} &
\textbf{Primary data source} &
\textbf{Falsification (truncated)} \\
\endfirsthead
\textbf{Equation label} &
\textbf{Case study location} &
\textbf{Tier} &
\textbf{Primary data source} &
\textbf{Falsification (truncated)} \\
\endhead
\multicolumn{5}{r}{\small\itshape continued on next page} \\
\endfoot
\endlastfoot
\multicolumn{5}{l}{\small\bfseries Chapter 1: Redefining Racism} \\
\eqlab{eq:\discretionary{}{}{}1.\discretionary{}{}{}1-\discretionary{}{}{}enclosure-\discretionary{}{}{}score} & Ch.~1, l.~241 & \textcolor{gray}{\textbf{T3}} & \textit{(ordinal)} & Falsified if a population with all three outlets blocked (e\_i=1) shows \ldots \\
\eqlab{eq:\discretionary{}{}{}1.\discretionary{}{}{}2-\discretionary{}{}{}total-\discretionary{}{}{}enclosure} & Ch.~1, l.~257 & \textcolor{gray}{\textbf{T3}} & \textit{(ordinal)} & Falsified if a fully blocked population (e\_1=e\_2=e\_3=1) achieves coordi\ldots \\
\eqlab{eq:\discretionary{}{}{}1.\discretionary{}{}{}3-\discretionary{}{}{}conventional-\discretionary{}{}{}causal-\discretionary{}{}{}arrow} & Ch.~1, l.~291 & \textcolor{gray}{\textbf{T3}} & \textit{(ordinal)} & Falsified if documented evidence shows individual prejudice predating E\ldots \\
\eqlab{eq:\discretionary{}{}{}1.\discretionary{}{}{}4-\discretionary{}{}{}elite-\discretionary{}{}{}causal-\discretionary{}{}{}arrow} & Ch.~1, l.~297 & \textcolor{gray}{\textbf{T3}} & Zurara & Falsified if racialization systems are shown to emerge without Elite ec\ldots \\
\textbf{\eqlab{eq:\discretionary{}{}{}1.\discretionary{}{}{}5-\discretionary{}{}{}kernel-\discretionary{}{}{}optimization}} & Ch.~1, l.~495 & \textcolor{green!60!black}{\textbf{T1}} & Piketty & Falsified if a documented period shows sustained decline in Elite extra\ldots \\
\eqlab{eq:\discretionary{}{}{}1.\discretionary{}{}{}6-\discretionary{}{}{}interface-\discretionary{}{}{}strategy-\discretionary{}{}{}set} & Ch.~1, l.~576 & \textcolor{gray}{\textbf{T3}} & \textit{(ordinal)} & Falsified if a documented Elite extraction system used a strategy outsi\ldots \\
\eqlab{eq:\discretionary{}{}{}1.\discretionary{}{}{}7-\discretionary{}{}{}interface-\discretionary{}{}{}optimizer} & Ch.~1, l.~608 & \textcolor{blue!70!black}{\textbf{T2}} & Dudziak (2000) --- Cold War Civil Rights & Falsified if a historical interface transition increased net cost C\_tot\ldots \\
\textbf{\eqlab{eq:\discretionary{}{}{}1.\discretionary{}{}{}8-\discretionary{}{}{}suppression-\discretionary{}{}{}envelope}} & Ch.~1, l.~714 & \textcolor{blue!70!black}{\textbf{T2}} & BLS Union Membership Historical Series & Falsified if M(t) consistently exceeds $\Sigma$\_sup(t) for extended periods wi\ldots \\
\eqlab{eq:\discretionary{}{}{}1.\discretionary{}{}{}9-\discretionary{}{}{}crash-\discretionary{}{}{}condition-\discretionary{}{}{}derivative} & Ch.~1, l.~710 & \textcolor{gray}{\textbf{T3}} & \textit{(ordinal)} & Falsified if a crash event (Bacon's Rebellion, Haitian Revolution) occu\ldots \\
\eqlab{eq:\discretionary{}{}{}1.\discretionary{}{}{}10-\discretionary{}{}{}effective-\discretionary{}{}{}coherence} & Ch.~1, l.~716 & \textcolor{blue!70!black}{\textbf{T2}} & ANES cross-racial coalition data 1948--2020 & Falsified if a crash event occurs when M\_eff(t) < $\tau$, or if no crash occ\ldots \\
\eqlab{eq:\discretionary{}{}{}1.\discretionary{}{}{}11-\discretionary{}{}{}phase-\discretionary{}{}{}loading} & Ch.~1, l.~838 & \textcolor{gray}{\textbf{T3}} & ANES Time Series Study 1948--2020 & Falsified if survey-measured cross-group solidarity remains high when $\Phi$\ldots \\
\multicolumn{5}{l}{\small\bfseries Chapter 2: Version 1.0: Initializing the Vector (15th-Century Portugal)} \\
\eqlab{eq:\discretionary{}{}{}2.\discretionary{}{}{}1-\discretionary{}{}{}racism-\discretionary{}{}{}vector} & Ch.~2, l.~1091 & \textcolor{gray}{\textbf{T3}} & \textit{(ordinal)} & Falsified if documented racism operates without a state-power direction\ldots \\
\eqlab{eq:\discretionary{}{}{}2.\discretionary{}{}{}2-\discretionary{}{}{}status-\discretionary{}{}{}suppression-\discretionary{}{}{}allocation} & Ch.~2, l.~1163 & \textcolor{gray}{\textbf{T3}} & Du Bois (1935) --- Black Reconstruction in America & Falsified if Buffer Class loyalty is maintained without a non-material \ldots \\
\eqlab{eq:\discretionary{}{}{}2.\discretionary{}{}{}3-\discretionary{}{}{}psi-\discretionary{}{}{}null-\discretionary{}{}{}in-\discretionary{}{}{}survival} & Ch.~2, l.~1167 & \textcolor{gray}{\textbf{T3}} & \textit{(ordinal)} & Falsified if material wages to Buffer Class are shown to be independent\ldots \\
\eqlab{eq:\discretionary{}{}{}2.\discretionary{}{}{}4-\discretionary{}{}{}roman-\discretionary{}{}{}membership-\discretionary{}{}{}function} & Ch.~2, l.~1208 & \textcolor{gray}{\textbf{T3}} & Bradley (1994) --- Slavery and Society at Rome & Falsified if Roman enslaved population is shown to be phenotypically ho\ldots \\
\eqlab{eq:\discretionary{}{}{}2.\discretionary{}{}{}5-\discretionary{}{}{}american-\discretionary{}{}{}racial-\discretionary{}{}{}function} & Ch.~2, l.~1217 & \textcolor{gray}{\textbf{T3}} & Hening's Statutes at Large (Virginia) Vol. 2 & Falsified if documented American slavery assigned status based on non-p\ldots \\
\eqlab{eq:\discretionary{}{}{}2.\discretionary{}{}{}6-\discretionary{}{}{}roman-\discretionary{}{}{}american-\discretionary{}{}{}structural-\discretionary{}{}{}parallel} & Ch.~2, l.~1235 & \textcolor{gray}{\textbf{T3}} & \textit{(ordinal)} & Falsified if the structural mapping between Roman and American slavery \ldots \\
\eqlab{eq:\discretionary{}{}{}2.\discretionary{}{}{}7-\discretionary{}{}{}three-\discretionary{}{}{}tier-\discretionary{}{}{}portugal-\discretionary{}{}{}inequality} & Ch.~2, l.~1280 & \textcolor{gray}{\textbf{T3}} & \textit{(ordinal)} & Falsified if measured material outcomes show Benefit(I\_buffer) $\geq$ Benefi\ldots \\
\eqlab{eq:\discretionary{}{}{}2.\discretionary{}{}{}8-\discretionary{}{}{}church-\discretionary{}{}{}science-\discretionary{}{}{}feedback-\discretionary{}{}{}loop} & Ch.~2, l.~1335 & \textcolor{gray}{\textbf{T3}} & Cartwright (1851) --- Diseases and Peculiarities of t\ldots & Falsified if the religious-scientific legitimation loop is disrupted fr\ldots \\
\multicolumn{5}{l}{\small\bfseries Chapter 3: The Application: Bacon's Rebellion, the Buffer Class, and the Constitutional Patch (1676--1787)} \\
\textbf{\eqlab{eq:\discretionary{}{}{}3.\discretionary{}{}{}1-\discretionary{}{}{}bacon-\discretionary{}{}{}solidarity-\discretionary{}{}{}condition}} & Ch.~3, l.~1580 & \textcolor{green!60!black}{\textbf{T1}} & Morgan (1975) --- American Slavery & Falsified if an antebellum cross-racial coalition reached M(t) > $\tau$ with\ldots \\
\eqlab{eq:\discretionary{}{}{}3.\discretionary{}{}{}2-\discretionary{}{}{}boundary-\discretionary{}{}{}enforcement-\discretionary{}{}{}1691} & Ch.~3, l.~1507 & \textcolor{gray}{\textbf{T3}} & Hening's Statutes at Large (Virginia) Vol. 3 & Falsified if racial boundary enforcement persisted without legal penalt\ldots \\
\eqlab{eq:\discretionary{}{}{}3.\discretionary{}{}{}3-\discretionary{}{}{}race-\discretionary{}{}{}destroys-\discretionary{}{}{}class} & Ch.~3, l.~1539 & \textcolor{gray}{\textbf{T3}} & Morgan (1975) --- American Slavery & Falsified if colonial-era racial coding produced class solidarity rathe\ldots \\
\eqlab{eq:\discretionary{}{}{}3.\discretionary{}{}{}4-\discretionary{}{}{}three-\discretionary{}{}{}tier-\discretionary{}{}{}post-\discretionary{}{}{}bacon} & Ch.~3, l.~1560 & \textcolor{gray}{\textbf{T3}} & \textit{(ordinal)} & Falsified if the post-Bacon ordering shows Buffer Class benefits exceed\ldots \\
\eqlab{eq:\discretionary{}{}{}3.\discretionary{}{}{}5-\discretionary{}{}{}kinetic-\discretionary{}{}{}necessary-\discretionary{}{}{}condition} & Ch.~3, l.~1595 & \textcolor{green!60!black}{\textbf{T1}} & Morgan (1975) --- American Slavery & Falsified if documented cases show Elite extraction maintained when K\_E\ldots \\
\multicolumn{5}{l}{\small\bfseries Chapter 5: The Gendered Axis: Coverture, Eugenics, and the Reproductive Extraction Kernel} \\
\eqlab{eq:\discretionary{}{}{}5.\discretionary{}{}{}1-\discretionary{}{}{}compliance-\discretionary{}{}{}differential} & Ch.~5, l.~2272 & \textcolor{blue!70!black}{\textbf{T2}} & Baptist (2014) --- The Half Has Never Been Told & Falsified if V\_c(x) $\leq$ V\_r(x) for subjects in documented extraction regi\ldots \\
\multicolumn{5}{l}{\small\bfseries Chapter 6: The Enforcement Engine: Slave Patrols, the 13th Amendment, and the Compounding Model (1704--1865)} \\
\eqlab{eq:\discretionary{}{}{}6.\discretionary{}{}{}1-\discretionary{}{}{}lethal-\discretionary{}{}{}autonomy-\discretionary{}{}{}gradient} & Ch.~6, l.~2889 & \textcolor{blue!70!black}{\textbf{T2}} & LEOSA --- 18 U.S.C. §926B statutory text & Falsified if Mapping Police Violence data shows parity in police killin\ldots \\
\textbf{\eqlab{eq:\discretionary{}{}{}6.\discretionary{}{}{}2-\discretionary{}{}{}four-\discretionary{}{}{}tier-\discretionary{}{}{}benefit-\discretionary{}{}{}ordering}} & Ch.~6, l.~2994 & \textcolor{gray}{\textbf{T3}} & Piketty-Saez-Zucman Distributional National Accounts & Falsified if comprehensive wealth, health, or carceral data shows Benef\ldots \\
\eqlab{eq:\discretionary{}{}{}6.\discretionary{}{}{}3-\discretionary{}{}{}thirteenth-\discretionary{}{}{}amendment-\discretionary{}{}{}exception} & Ch.~6, l.~3080 & \textcolor{gray}{\textbf{T3}} & Blackmon (2008) --- Slavery by Another Name & Falsified if 13th Amendment exception clause is shown to have been appl\ldots \\
\eqlab{eq:\discretionary{}{}{}6.\discretionary{}{}{}4-\discretionary{}{}{}additive-\discretionary{}{}{}policy-\discretionary{}{}{}model} & Ch.~6, l.~3144 & \textcolor{gray}{\textbf{T3}} & \textit{(ordinal)} & Falsified if additive accumulation predicts compound policy outcomes wi\ldots \\
\eqlab{eq:\discretionary{}{}{}6.\discretionary{}{}{}5-\discretionary{}{}{}compounding-\discretionary{}{}{}temporal-\discretionary{}{}{}model} & Ch.~6, l.~3154 & \textcolor{blue!70!black}{\textbf{T2}} & Hamilton \& Darity (2017) --- The Political Economy of\ldots & Falsified if policy shock effects on O\textasciicircum{}capacity are shown to be indepen\ldots \\
\eqlab{eq:\discretionary{}{}{}6.\discretionary{}{}{}6-\discretionary{}{}{}asymmetric-\discretionary{}{}{}enforcement-\discretionary{}{}{}multiplier} & Ch.~6, l.~3172 & \textcolor{green!60!black}{\textbf{T1}} & ACLU (2020) --- A Tale of Two Countries: Racially Tar\ldots & Falsified if ACLU cannabis arrest data shows no racial disparity in pol\ldots \\
\eqlab{eq:\discretionary{}{}{}6.\discretionary{}{}{}7-\discretionary{}{}{}policy-\discretionary{}{}{}sequence-\discretionary{}{}{}noncommutative} & Ch.~6, l.~3177 & \textcolor{gray}{\textbf{T3}} & \textit{(ordinal)} & Falsified if order of policy application (P\_1 then P\_2 vs. P\_2 then P\_1\ldots \\
\eqlab{eq:\discretionary{}{}{}6.\discretionary{}{}{}8-\discretionary{}{}{}capacity-\discretionary{}{}{}chain-\discretionary{}{}{}1619} & Ch.~6, l.~3193 & \textcolor{blue!70!black}{\textbf{T2}} & Darity \& Mullen (2020) --- From Here to Equality (rep\ldots & Falsified if O\_1619 capacity is non-significantly different from O\_1450\ldots \\
\eqlab{eq:\discretionary{}{}{}6.\discretionary{}{}{}9-\discretionary{}{}{}capacity-\discretionary{}{}{}chain-\discretionary{}{}{}1865} & Ch.~6, l.~3196 & \textcolor{blue!70!black}{\textbf{T2}} & Blackmon (2008) --- Slavery by Another Name & Falsified if O\_1865 capacity is not reduced from O\_1619 as measured by \ldots \\
\eqlab{eq:\discretionary{}{}{}6.\discretionary{}{}{}10-\discretionary{}{}{}capacity-\discretionary{}{}{}chain-\discretionary{}{}{}1934} & Ch.~6, l.~3199 & \textcolor{green!60!black}{\textbf{T1}} & Mapping Inequality --- HOLC Redlining Maps & Falsified if redlining-exposed tracts show no differential wealth accum\ldots \\
\eqlab{eq:\discretionary{}{}{}6.\discretionary{}{}{}11-\discretionary{}{}{}capacity-\discretionary{}{}{}chain-\discretionary{}{}{}1971} & Ch.~6, l.~3230 & \textcolor{green!60!black}{\textbf{T1}} & ACLU (2020) --- Cannabis Arrests Report & Falsified if War on Drugs enforcement shows no differential impact on O\ldots \\
\textbf{\eqlab{eq:\discretionary{}{}{}6.\discretionary{}{}{}12-\discretionary{}{}{}capacity-\discretionary{}{}{}compounding-\discretionary{}{}{}full}} & Ch.~6, l.~3169 & \textcolor{green!60!black}{\textbf{T1}} & ACLU (2020) --- A Tale of Two Countries & Falsified if any one policy shock ($\alpha$, $\beta$, $\gamma$, $\delta$) is shown to have zero ma\ldots \\
\multicolumn{5}{l}{\small\bfseries Chapter 7: The Containment: Pullman, Redlining, and the Wages of Whiteness (1894--1965)} \\
\eqlab{eq:\discretionary{}{}{}7.\discretionary{}{}{}1-\discretionary{}{}{}pullman-\discretionary{}{}{}corollary} & Ch.~7, l.~3624 & \textcolor{gray}{\textbf{T3}} & Roediger (1991) --- The Wages of Whiteness & Falsified if a documented case shows I\_buffer achieving collective gain\ldots \\
\multicolumn{5}{l}{\small\bfseries Chapter 8: Tweedism and the Puppet Class: The Algorithmic Filter on Democracy} \\
\eqlab{eq:\discretionary{}{}{}8.\discretionary{}{}{}1-\discretionary{}{}{}vertex-\discretionary{}{}{}partition-\discretionary{}{}{}graph} & Ch.~8, l.~4020 & \textcolor{gray}{\textbf{T3}} & \textit{(ordinal)} & Falsified if leader/follower partition fails to predict convergence in \ldots \\
\eqlab{eq:\discretionary{}{}{}8.\discretionary{}{}{}2-\discretionary{}{}{}tier-\discretionary{}{}{}to-\discretionary{}{}{}graph-\discretionary{}{}{}mapping} & Ch.~8, l.~4024 & \textcolor{gray}{\textbf{T3}} & \textit{(ordinal)} & Falsified if the tier-to-graph mapping produces inconsistent prediction\ldots \\
\eqlab{eq:\discretionary{}{}{}8.\discretionary{}{}{}3-\discretionary{}{}{}fractional-\discretionary{}{}{}agent-\discretionary{}{}{}dynamics} & Ch.~8, l.~4032 & \textcolor{gray}{\textbf{T3}} & Li et al. --- Mittag-Leffler stability results (mathe\ldots & Falsified if integer-order dynamics ($\alpha$=1) predict social-mobility conve\ldots \\
\eqlab{eq:\discretionary{}{}{}8.\discretionary{}{}{}4-\discretionary{}{}{}compounding-\discretionary{}{}{}chain-\discretionary{}{}{}formal} & Ch.~8, l.~3187 & \textcolor{blue!70!black}{\textbf{T2}} & Federal Reserve SCF racial wealth data & Falsified if multiplicative chain prediction diverges significantly fro\ldots \\
\eqlab{eq:\discretionary{}{}{}8.\discretionary{}{}{}5-\discretionary{}{}{}stationary-\discretionary{}{}{}leader-\discretionary{}{}{}condition} & Ch.~8, l.~4046 & \textcolor{gray}{\textbf{T3}} & \textit{(ordinal)} & Falsified if documented Elite actors show systematic response to u\_i (u\ldots \\
\eqlab{eq:\discretionary{}{}{}8.\discretionary{}{}{}6-\discretionary{}{}{}proximity-\discretionary{}{}{}connectivity-\discretionary{}{}{}condition} & Ch.~8, l.~4128 & \textcolor{gray}{\textbf{T3}} & \textit{(ordinal)} & Falsified if social-affinity graph connectivity is shown to be independ\ldots \\
\eqlab{eq:\discretionary{}{}{}8.\discretionary{}{}{}7-\discretionary{}{}{}navigation-\discretionary{}{}{}components} & Ch.~8, l.~4220 & \textcolor{gray}{\textbf{T3}} & \textit{(ordinal)} & Falsified if goal and constraint functions do not bound follower trajec\ldots \\
\eqlab{eq:\discretionary{}{}{}8.\discretionary{}{}{}8-\discretionary{}{}{}navigation-\discretionary{}{}{}function} & Ch.~8, l.~4223 & \textcolor{gray}{\textbf{T3}} & \textit{(ordinal)} & Falsified if $\phi$\_i fails to produce convergence in any graph satisfying t\ldots \\
\eqlab{eq:\discretionary{}{}{}8.\discretionary{}{}{}9-\discretionary{}{}{}gradient-\discretionary{}{}{}control-\discretionary{}{}{}law} & Ch.~8, l.~4228 & \textcolor{gray}{\textbf{T3}} & \textit{(ordinal)} & Falsified if negative gradient of $\phi$\_i fails to drive convergence in any\ldots \\
\eqlab{eq:\discretionary{}{}{}8.\discretionary{}{}{}10-\discretionary{}{}{}convex-\discretionary{}{}{}hull-\discretionary{}{}{}convergence} & Ch.~8, l.~4316 & \textcolor{gray}{\textbf{T3}} & \textit{(ordinal)} & Falsified if followers fail to converge to Elite convex hull in any gra\ldots \\
\eqlab{eq:\discretionary{}{}{}8.\discretionary{}{}{}11-\discretionary{}{}{}mittag-\discretionary{}{}{}leffler-\discretionary{}{}{}stability} & Ch.~8, l.~4402 & \textcolor{gray}{\textbf{T3}} & Chetty et al. (2018) --- Race and Economic Opportunit\ldots & Falsified if social-mobility data for $\alpha$ $\in$ (0,1) shows exponential rathe\ldots \\
\textbf{\eqlab{eq:\discretionary{}{}{}8.\discretionary{}{}{}12-\discretionary{}{}{}class-\discretionary{}{}{}alignment-\discretionary{}{}{}base-\discretionary{}{}{}waves}} & Ch.~8, l.~4968 & \textcolor{green!60!black}{\textbf{T1}} & Google Trends & Falsified if class alignment signal is not measurable as a distinct fre\ldots \\
\eqlab{eq:\discretionary{}{}{}8.\discretionary{}{}{}13-\discretionary{}{}{}subgroup-\discretionary{}{}{}compound-\discretionary{}{}{}phase} & Ch.~8, l.~4968 & \textcolor{blue!70!black}{\textbf{T2}} & Google Trends & Falsified if compound phase $\Phi$\_j fails to predict subgroup political ali\ldots \\
\eqlab{eq:\discretionary{}{}{}8.\discretionary{}{}{}14-\discretionary{}{}{}net-\discretionary{}{}{}solidarity-\discretionary{}{}{}signal} & Ch.~8, l.~4968 & \textcolor{green!60!black}{\textbf{T1}} & Google Trends & Falsified if net solidarity signal cannot be measured as distinct from \ldots \\
\eqlab{eq:\discretionary{}{}{}8.\discretionary{}{}{}15-\discretionary{}{}{}solidarity-\discretionary{}{}{}collapse-\discretionary{}{}{}condition} & Ch.~8, l.~4968 & \textcolor{green!60!black}{\textbf{T1}} & Google Trends & Falsified if M(t) remains above $\tau$ even when S\_class approaches zero acr\ldots \\
\eqlab{eq:\discretionary{}{}{}8.\discretionary{}{}{}16-\discretionary{}{}{}interference-\discretionary{}{}{}control-\discretionary{}{}{}objective} & Ch.~8, l.~4968 & \textcolor{green!60!black}{\textbf{T1}} & Google Trends & Falsified if observed interference engine output increases S\_class rath\ldots \\
\eqlab{eq:\discretionary{}{}{}8.\discretionary{}{}{}17-\discretionary{}{}{}circular-\discretionary{}{}{}dispersion-\discretionary{}{}{}operator} & Ch.~8, l.~4968 & \textcolor{blue!70!black}{\textbf{T2}} & ANES Time Series --- cross-group solidarity items & Falsified if circular dispersion of phase values fails to predict cross\ldots \\
\textbf{\eqlab{eq:\discretionary{}{}{}8.\discretionary{}{}{}18-\discretionary{}{}{}tweedism-\discretionary{}{}{}agenda-\discretionary{}{}{}path}} & Ch.~8, l.~4844 & \textcolor{green!60!black}{\textbf{T1}} & Gilens \& Page (2014) --- Testing Theories of American\ldots & Falsified if Gilens-Page analysis shows median voter preferences predic\ldots \\
\multicolumn{5}{l}{\small\bfseries Chapter 9: The Recompile: COINTELPRO, the Variable Swap, and the War on Drugs (1968--1994)} \\
\textbf{\eqlab{eq:\discretionary{}{}{}9.\discretionary{}{}{}1-\discretionary{}{}{}lead-\discretionary{}{}{}crime-\discretionary{}{}{}compounding}} & Ch.~9, l.~6344 & \textcolor{green!60!black}{\textbf{T1}} & Reyes (2007) --- Environmental Policy as Social Policy & Falsified if Reyes (2007) blood-lead time-series fails to predict crime\ldots \\
\eqlab{eq:\discretionary{}{}{}9.\discretionary{}{}{}2-\discretionary{}{}{}epistemic-\discretionary{}{}{}suppression-\discretionary{}{}{}variable} & Ch.~9, l.~6344 & \textcolor{blue!70!black}{\textbf{T2}} & Kitman (2000) --- The Secret History of Lead (The Nat\ldots & Falsified if documented lead industry cover-up timeline shows no correl\ldots \\
\eqlab{eq:\discretionary{}{}{}9.\discretionary{}{}{}3-\discretionary{}{}{}multi-\discretionary{}{}{}vector-\discretionary{}{}{}lead-\discretionary{}{}{}exposure} & Ch.~9, l.~6344 & \textcolor{green!60!black}{\textbf{T1}} & EPA Air Quality Monitoring Data & Falsified if cumulative lead exposure from all three vectors is not sig\ldots \\
\eqlab{eq:\discretionary{}{}{}9.\discretionary{}{}{}4-\discretionary{}{}{}redlining-\discretionary{}{}{}containment-\discretionary{}{}{}field} & Ch.~9, l.~6344 & \textcolor{green!60!black}{\textbf{T1}} & Mapping Inequality HOLC maps & Falsified if HOLC map grades show no statistically significant correlat\ldots \\
\eqlab{eq:\discretionary{}{}{}9.\discretionary{}{}{}5-\discretionary{}{}{}school-\discretionary{}{}{}funding-\discretionary{}{}{}property-\discretionary{}{}{}value} & Ch.~9, l.~6337 & \textcolor{green!60!black}{\textbf{T1}} & NCES Public School Finance Survey & Falsified if school funding is shown to be independent of property valu\ldots \\
\eqlab{eq:\discretionary{}{}{}9.\discretionary{}{}{}6-\discretionary{}{}{}infrastructure-\discretionary{}{}{}quality-\discretionary{}{}{}funding} & Ch.~9, l.~6337 & \textcolor{blue!70!black}{\textbf{T2}} & NCES School Facility Condition Survey & Falsified if school infrastructure quality is shown to be unrelated to \ldots \\
\eqlab{eq:\discretionary{}{}{}9.\discretionary{}{}{}7-\discretionary{}{}{}school-\discretionary{}{}{}lead-\discretionary{}{}{}exposure-\discretionary{}{}{}inverse} & Ch.~9, l.~6337 & \textcolor{green!60!black}{\textbf{T1}} & EPA 3Ts for Reducing Lead in Drinking Water in Scho\ldots & Falsified if school infrastructure quality is shown to be unrelated to \ldots \\
\eqlab{eq:\discretionary{}{}{}9.\discretionary{}{}{}8-\discretionary{}{}{}community-\discretionary{}{}{}capacity-\discretionary{}{}{}lead-\discretionary{}{}{}reduction} & Ch.~9, l.~6337 & \textcolor{blue!70!black}{\textbf{T2}} & Aizer \& Currie (2019) --- Lead and Juvenile Delinquen\ldots & Falsified if community economic capacity is shown to be independent of \ldots \\
\eqlab{eq:\discretionary{}{}{}9.\discretionary{}{}{}9-\discretionary{}{}{}property-\discretionary{}{}{}value-\discretionary{}{}{}capacity-\discretionary{}{}{}feedback} & Ch.~9, l.~6337 & \textcolor{green!60!black}{\textbf{T1}} & Mapping Inequality HOLC maps + ACS property value d\ldots & Falsified if property values in subsequent periods are shown to be inde\ldots \\
\eqlab{eq:\discretionary{}{}{}9.\discretionary{}{}{}10-\discretionary{}{}{}highway-\discretionary{}{}{}lead-\discretionary{}{}{}spatial-\discretionary{}{}{}concentration} & Ch.~9, l.~7178 & \textcolor{green!60!black}{\textbf{T1}} & Mapping Inequality DSL (Nelson et al.) --- HOLC grade\ldots & Falsified if HOLC-D tracts show no excess firearm or lead burden relati\ldots \\
\multicolumn{5}{l}{\small\bfseries Chapter 10: The Full Algorithm: Demographic Paradox, Cannibalization, and the 5-Tier Reveal (1994--Present)} \\
\eqlab{eq:\discretionary{}{}{}10.\discretionary{}{}{}1-\discretionary{}{}{}outgroup-\discretionary{}{}{}racialized-\discretionary{}{}{}subset} & Ch.~10, l.~7051 & \textcolor{gray}{\textbf{T3}} & US Census Bureau demographic data & Falsified if O\_racialized is shown to be coextensive with O rather than\ldots \\
\eqlab{eq:\discretionary{}{}{}10.\discretionary{}{}{}2-\discretionary{}{}{}buffer-\discretionary{}{}{}class-\discretionary{}{}{}shrinkage} & Ch.~10, l.~7055 & \textcolor{blue!70!black}{\textbf{T2}} & Pew Research Center --- America's Shrinking Middle Cl\ldots & Falsified if Buffer Class size is shown to be expanding rather than con\ldots \\
\eqlab{eq:\discretionary{}{}{}10.\discretionary{}{}{}3-\discretionary{}{}{}cannibalization-\discretionary{}{}{}equation} & Ch.~10, l.~7153 & \textcolor{gray}{\textbf{T3}} & Saez-Zucman wealth share series & Falsified if documented Elite extraction rate increases without corresp\ldots \\
\eqlab{eq:\discretionary{}{}{}10.\discretionary{}{}{}4-\discretionary{}{}{}predatory-\discretionary{}{}{}minmax-\discretionary{}{}{}full-\discretionary{}{}{}definition} & Ch.~10, l.~7303 & \textcolor{gray}{\textbf{T3}} & \textit{(ordinal)} & Falsified if the five-tier hierarchy fails to map onto any contemporary\ldots \\
\eqlab{eq:\discretionary{}{}{}10.\discretionary{}{}{}5-\discretionary{}{}{}kernel-\discretionary{}{}{}objective-\discretionary{}{}{}canonical} & Ch.~10, l.~7477 & \textcolor{green!60!black}{\textbf{T1}} & Piketty-Saez-Zucman Distributional National Accounts & Falsified if a documented period shows sustained Elite wealth-share dec\ldots \\
\eqlab{eq:\discretionary{}{}{}10.\discretionary{}{}{}6-\discretionary{}{}{}effective-\discretionary{}{}{}resistance-\discretionary{}{}{}variable} & Ch.~10, l.~7309 & \textcolor{gray}{\textbf{T3}} & \textit{(ordinal)} & Falsified if M\_eff fails to capture the difference between raw class co\ldots \\
\eqlab{eq:\discretionary{}{}{}10.\discretionary{}{}{}7-\discretionary{}{}{}suppression-\discretionary{}{}{}envelope-\discretionary{}{}{}canonical} & Ch.~10, l.~7345 & \textcolor{blue!70!black}{\textbf{T2}} & \textit{(pending)} & Falsified if documented suppression tools outside this set constitute t\ldots \\
\eqlab{eq:\discretionary{}{}{}10.\discretionary{}{}{}8-\discretionary{}{}{}crash-\discretionary{}{}{}condition-\discretionary{}{}{}canonical} & Ch.~10, l.~7319 & \textcolor{gray}{\textbf{T3}} & \textit{(ordinal)} & Falsified if crash events occur in documented cases without M\_eff(t) > \ldots \\
\eqlab{eq:\discretionary{}{}{}10.\discretionary{}{}{}9-\discretionary{}{}{}tier-\discretionary{}{}{}benefit-\discretionary{}{}{}ordering-\discretionary{}{}{}canonical} & Ch.~10, l.~7324 & \textcolor{gray}{\textbf{T3}} & \textit{(ordinal)} & Falsified if any two-tier comparison fails to show stated ordering in c\ldots \\
\eqlab{eq:\discretionary{}{}{}10.\discretionary{}{}{}10-\discretionary{}{}{}system-\discretionary{}{}{}objective-\discretionary{}{}{}extraction-\discretionary{}{}{}ratio} & Ch.~10, l.~7405 & \textcolor{gray}{\textbf{T3}} & \textit{(ordinal)} & Falsified if documented extraction systems show this ratio declining wi\ldots \\
\eqlab{eq:\discretionary{}{}{}10.\discretionary{}{}{}11-\discretionary{}{}{}system-\discretionary{}{}{}stability-\discretionary{}{}{}algorithm} & Ch.~10, l.~7405 & \textcolor{gray}{\textbf{T3}} & \textit{(ordinal)} & Falsified if system stability is maintained when extraction/resistance \ldots \\
\textbf{\eqlab{eq:\discretionary{}{}{}10.\discretionary{}{}{}12-\discretionary{}{}{}o-\discretionary{}{}{}final-\discretionary{}{}{}construction}} & Ch.~10, l.~6828 & \textcolor{green!60!black}{\textbf{T1}} & Bureau of Justice Statistics incarceration by race \ldots & Falsified if modern mass-incarceration demographics fail to reproduce t\ldots \\
\eqlab{eq:\discretionary{}{}{}10.\discretionary{}{}{}13-\discretionary{}{}{}reclassification-\discretionary{}{}{}operator} & Ch.~10, l.~7734 & \textcolor{gray}{\textbf{T3}} & \textit{(ordinal)} & Falsified if the operator fails to predict observed ideological sorting\ldots \\
\eqlab{eq:\discretionary{}{}{}10.\discretionary{}{}{}14-\discretionary{}{}{}kinetic-\discretionary{}{}{}power-\discretionary{}{}{}distribution} & Ch.~10, l.~7768 & \textcolor{blue!70!black}{\textbf{T2}} & Pew Research Center --- Gun Ownership Survey 2017 & Falsified if kinetic power is not concentrated in F\_enforce as a propor\ldots \\
\eqlab{eq:\discretionary{}{}{}10.\discretionary{}{}{}15-\discretionary{}{}{}temporal-\discretionary{}{}{}enclosure-\discretionary{}{}{}definition} & Ch.~10, l.~7820 & \textcolor{gray}{\textbf{T3}} & Reinhart-Sbrancia 2015 (Economic Policy) & Falsified if sovereign debt accumulation is demonstrably unrelated to p\ldots \\
\textbf{\eqlab{eq:\discretionary{}{}{}10.\discretionary{}{}{}16-\discretionary{}{}{}financial-\discretionary{}{}{}repression-\discretionary{}{}{}condition}} & Ch.~10, l.~7825 & \textcolor{green!60!black}{\textbf{T1}} & FRED 10Y Treasury Constant Maturity Rate (GS10) & Falsified if: (1) real interest rates remain positive throughout the 19\ldots \\
\eqlab{eq:\discretionary{}{}{}10.\discretionary{}{}{}17-\discretionary{}{}{}temporal-\discretionary{}{}{}extraction-\discretionary{}{}{}rate} & Ch.~10, l.~7825 & \textcolor{green!60!black}{\textbf{T1}} & FRED: Federal Debt Total Public Debt as \% of GDP (G\ldots & Falsified if the indicator function 1[r < g] is inactive during documen\ldots \\
\textbf{\eqlab{eq:\discretionary{}{}{}10.\discretionary{}{}{}18-\discretionary{}{}{}psi-\discretionary{}{}{}degradation-\discretionary{}{}{}function}} & Ch.~10, l.~7905 & \textcolor{green!60!black}{\textbf{T1}} & FRED: Real Median Hourly Compensation (B4701C0A052N\ldots & Falsified if: (1) Chow test at 1971 shows no statistically significant \ldots \\
\eqlab{eq:\discretionary{}{}{}10.\discretionary{}{}{}19-\discretionary{}{}{}saeculum-\discretionary{}{}{}phase-\discretionary{}{}{}map} & Ch.~10, l.~7900 & \textcolor{blue!70!black}{\textbf{T2}} & Strauss \& Howe, The Fourth Turning (1997) & Falsified if: (1) the generational archetype sequence does not follow t\ldots \\
\textbf{\eqlab{eq:\discretionary{}{}{}10.\discretionary{}{}{}20-\discretionary{}{}{}dalio-\discretionary{}{}{}extraction-\discretionary{}{}{}bound}} & Ch.~10, l.~7950 & \textcolor{blue!70!black}{\textbf{T2}} & Maddison Project Database 2020 (GDP, trade output) & Falsified if: (1) the reconstructed US Empire Index does not show a pea\ldots \\
\multicolumn{5}{l}{\small\bfseries Chapter 11: The Kinetic Guarantee: Arms Asymmetry, the Second Amendment, and the Disarmament Timeline} \\
\textbf{\eqlab{eq:\discretionary{}{}{}11.\discretionary{}{}{}1-\discretionary{}{}{}extraction-\discretionary{}{}{}precondition}} & Ch.~11, l.~7992 & \textcolor{green!60!black}{\textbf{T1}} & Primary legislation texts (sullivanlaw, nfa1934, mu\ldots & Falsified if a documented extraction system maintained stability with L\ldots \\
\eqlab{eq:\discretionary{}{}{}11.\discretionary{}{}{}2-\discretionary{}{}{}disarmament-\discretionary{}{}{}agenda-\discretionary{}{}{}path} & Ch.~11, l.~7992 & \textcolor{green!60!black}{\textbf{T1}} & Primary legislation and ruling texts (11 events) & Falsified if documented gun-control legislation passed under high $\Phi$\_loa\ldots \\
\eqlab{eq:\discretionary{}{}{}11.\discretionary{}{}{}3-\discretionary{}{}{}restriction-\discretionary{}{}{}precedent-\discretionary{}{}{}accumulation} & Ch.~11, l.~7992 & \textcolor{green!60!black}{\textbf{T1}} & Primary legislation and ruling texts (11 events) & Falsified if any documented legislative or ruling event reduced cumulat\ldots \\
\eqlab{eq:\discretionary{}{}{}11.\discretionary{}{}{}4-\discretionary{}{}{}restriction-\discretionary{}{}{}scope-\discretionary{}{}{}expansion} & Ch.~11, l.~7992 & \textcolor{green!60!black}{\textbf{T1}} & Primary legislation texts and enforcement records & Falsified if the historical record showed restriction scope declining o\ldots \\
\multicolumn{5}{l}{\small\bfseries Chapter 12: The Contradiction: Why Reform Serves the Algorithm} \\
\eqlab{eq:\discretionary{}{}{}12.\discretionary{}{}{}1-\discretionary{}{}{}reform-\discretionary{}{}{}absorption-\discretionary{}{}{}mechanism} & Ch.~12, l.~8859 & \textcolor{gray}{\textbf{T3}} & Piketty-Saez top income share data & Falsified if a documented reform reduces Elite extraction share (max) r\ldots \\
\eqlab{eq:\discretionary{}{}{}12.\discretionary{}{}{}2-\discretionary{}{}{}concession-\discretionary{}{}{}theorem} & Ch.~12, l.~8920 & \textcolor{blue!70!black}{\textbf{T2}} & Piven \& Cloward (1977) --- Poor People's Movements & Falsified if Elite-initiated reforms consistently exceed the minimum re\ldots \\
\eqlab{eq:\discretionary{}{}{}12.\discretionary{}{}{}3-\discretionary{}{}{}judiciary-\discretionary{}{}{}discrimination-\discretionary{}{}{}detection} & Ch.~12, l.~9071 & \textcolor{gray}{\textbf{T3}} & Washington v. Davis & Falsified if judicial detection function produces false positives for p\ldots \\
\eqlab{eq:\discretionary{}{}{}12.\discretionary{}{}{}4-\discretionary{}{}{}proxy-\discretionary{}{}{}discrimination-\discretionary{}{}{}equivalence} & Ch.~12, l.~9111 & \textcolor{gray}{\textbf{T3}} & Alexander (2010) --- The New Jim Crow & Falsified if P\_proxy achieves racial separation at rates significantly \ldots \\
\eqlab{eq:\discretionary{}{}{}12.\discretionary{}{}{}4a-\discretionary{}{}{}judicial-\discretionary{}{}{}double-\discretionary{}{}{}agent} & Ch.~12, l.~10000 & \textcolor{blue!70!black}{\textbf{T2}} & Local SCOTUS markdown corpus & Falsified if the SCOTUS corpus does not separate civil-rights/voting ca\ldots \\
\textbf{\eqlab{eq:\discretionary{}{}{}12.\discretionary{}{}{}5-\discretionary{}{}{}haitian-\discretionary{}{}{}theorem-\discretionary{}{}{}nonkinetic}} & Ch.~12, l.~8507 & \textcolor{green!60!black}{\textbf{T1}} & James (1938) The Black Jacobins & Falsified if any documented non-kinetic reform achieves $\Delta$max < 0 (susta\ldots \\
\textbf{\eqlab{eq:\discretionary{}{}{}12.\discretionary{}{}{}6-\discretionary{}{}{}haitian-\discretionary{}{}{}theorem-\discretionary{}{}{}kinetic}} & Ch.~12, l.~8507 & \textcolor{green!60!black}{\textbf{T1}} & James (1938) The Black Jacobins & Falsified if any kinetic event fails to produce at least temporary loca\ldots \\
\eqlab{eq:\discretionary{}{}{}12.\discretionary{}{}{}7-\discretionary{}{}{}gendered-\discretionary{}{}{}outgroup-\discretionary{}{}{}set} & Ch.~12, l.~9597 & \textcolor{gray}{\textbf{T3}} & \textit{(ordinal)} & Falsified if O\_gendered is shown to be structured differently from O\_ra\ldots \\
\eqlab{eq:\discretionary{}{}{}12.\discretionary{}{}{}8-\discretionary{}{}{}intersection-\discretionary{}{}{}double-\discretionary{}{}{}partition} & Ch.~12, l.~9602 & \textcolor{gray}{\textbf{T3}} & Crenshaw (1989) --- Demarginalizing the Intersection \ldots & Falsified if intersection O\_racialized $\cap$ O\_gendered does not show compo\ldots \\
\eqlab{eq:\discretionary{}{}{}12.\discretionary{}{}{}9-\discretionary{}{}{}multiplicative-\discretionary{}{}{}intersection-\discretionary{}{}{}compounding} & Ch.~12, l.~9607 & \textcolor{blue!70!black}{\textbf{T2}} & AAUW pay-gap data by race and gender & Falsified if Black women's wealth gap from white men is not multiplicat\ldots \\
\eqlab{eq:\discretionary{}{}{}12.\discretionary{}{}{}10-\discretionary{}{}{}nonviolence-\discretionary{}{}{}mandate-\discretionary{}{}{}function} & Ch.~12, l.~9709 & \textcolor{gray}{\textbf{T3}} & Carson (1981) --- In Struggle: SNCC and the Black Awa\ldots & Falsified if documented liberation movements adopting nonviolence achie\ldots \\
\multicolumn{5}{l}{\small\bfseries Chapter 13: The Global Containment Field: Scaling the Algorithm} \\
\eqlab{eq:\discretionary{}{}{}13.\discretionary{}{}{}1-\discretionary{}{}{}global-\discretionary{}{}{}compounding-\discretionary{}{}{}capacity} & Ch.~13, l.~9925 & \textcolor{blue!70!black}{\textbf{T2}} & World Bank development indicators --- GDP per capita \ldots & Falsified if post-colonial nation's capacity trajectory shows exponenti\ldots \\
\eqlab{eq:\discretionary{}{}{}13.\discretionary{}{}{}2-\discretionary{}{}{}global-\discretionary{}{}{}predatory-\discretionary{}{}{}minmax} & Ch.~13, l.~9973 & \textcolor{gray}{\textbf{T3}} & Klein (2007) --- The Shock Doctrine & Falsified if documented international economic relationships fail to ma\ldots \\
\eqlab{eq:\discretionary{}{}{}13.\discretionary{}{}{}3-\discretionary{}{}{}externalized-\discretionary{}{}{}enforcement-\discretionary{}{}{}envelope} & Ch.~13, l.~9993 & \textcolor{gray}{\textbf{T3}} & UN MINUSTAH mission records & Falsified if Haiti's domestic coercive capacity is shown to be structur\ldots \\
\eqlab{eq:\discretionary{}{}{}13.\discretionary{}{}{}4-\discretionary{}{}{}killick-\discretionary{}{}{}extraction-\discretionary{}{}{}equation} & Ch.~13, l.~10021 & \textcolor{gray}{\textbf{T3}} & NYT 'The Ransom' (2022) & Falsified if the extraction equation fails to predict subsequent debt i\ldots \\
\eqlab{eq:\discretionary{}{}{}13.\discretionary{}{}{}5-\discretionary{}{}{}legitimation-\discretionary{}{}{}constraint-\discretionary{}{}{}set} & Ch.~13, l.~10039 & \textcolor{gray}{\textbf{T3}} & \textit{(ordinal)} & Falsified if E\_global actors consistently operate outside their stated \ldots \\
\eqlab{eq:\discretionary{}{}{}13.\discretionary{}{}{}6-\discretionary{}{}{}authorization-\discretionary{}{}{}precondition} & Ch.~13, l.~10050 & \textcolor{gray}{\textbf{T3}} & \textit{(ordinal)} & Falsified if E\_global achieves extraction without satisfying authorizat\ldots \\
\eqlab{eq:\discretionary{}{}{}13.\discretionary{}{}{}7-\discretionary{}{}{}coercion-\discretionary{}{}{}legitimacy-\discretionary{}{}{}frontier} & Ch.~13, l.~10056 & \textcolor{gray}{\textbf{T3}} & \textit{(ordinal)} & Falsified if observed coercion levels in 'voluntary' agreements consist\ldots \\
\eqlab{eq:\discretionary{}{}{}13.\discretionary{}{}{}8-\discretionary{}{}{}firmin-\discretionary{}{}{}extraction-\discretionary{}{}{}threshold} & Ch.~13, l.~10062 & \textcolor{gray}{\textbf{T3}} & \textit{(ordinal)} & Falsified if E\_global achieves extraction with deployed force below F* \ldots \\
\eqlab{eq:\discretionary{}{}{}13.\discretionary{}{}{}9-\discretionary{}{}{}infeasible-\discretionary{}{}{}constraint-\discretionary{}{}{}set} & Ch.~13, l.~10075 & \textcolor{gray}{\textbf{T3}} & \textit{(ordinal)} & Falsified if E\_global achieves both extraction success and legitimation\ldots \\
\eqlab{eq:\discretionary{}{}{}13.\discretionary{}{}{}10-\discretionary{}{}{}asymmetric-\discretionary{}{}{}leverage-\discretionary{}{}{}ratio} & Ch.~13, l.~10100 & \textcolor{blue!70!black}{\textbf{T2}} & Firmin (1885) --- De l'égalité des races humaines & Falsified if a weaker actor's invocation of E\_global's legitimation arc\ldots \\
\eqlab{eq:\discretionary{}{}{}13.\discretionary{}{}{}11-\discretionary{}{}{}legitimation-\discretionary{}{}{}abandonment-\discretionary{}{}{}condition} & Ch.~13, l.~10107 & \textcolor{gray}{\textbf{T3}} & \textit{(ordinal)} & Falsified if E\_global maintains legitimation constraints even when extr\ldots \\
\eqlab{eq:\discretionary{}{}{}13.\discretionary{}{}{}12-\discretionary{}{}{}global-\discretionary{}{}{}containment-\discretionary{}{}{}strategy} & Ch.~13, l.~10123 & \textcolor{gray}{\textbf{T3}} & \textit{(ordinal)} & Falsified if global containment strategies differ structurally from the\ldots \\
\eqlab{eq:\discretionary{}{}{}13.\discretionary{}{}{}13-\discretionary{}{}{}un-\discretionary{}{}{}vote-\discretionary{}{}{}distribution} & Ch.~13, l.~10162 & \textcolor{gray}{\textbf{T3}} & UN Digital Library --- UNSC vote records & Falsified if UN Security Council voting distribution fails to map onto \ldots \\
\textbf{\eqlab{eq:\discretionary{}{}{}13.\discretionary{}{}{}14-\discretionary{}{}{}imperial-\discretionary{}{}{}core-\discretionary{}{}{}collapse}} & Ch.~13, l.~9208 & \textcolor{green!60!black}{\textbf{T1}} & World Bank Development Indicators --- GDP trajectory \ldots & Falsified if documented rising powers (China, OPEC, Asian Tigers) achie\ldots \\
\textbf{\eqlab{eq:\discretionary{}{}{}13.\discretionary{}{}{}15-\discretionary{}{}{}interface-\discretionary{}{}{}swap-\discretionary{}{}{}trigger}} & Ch.~13, l.~10280 & \textcolor{green!60!black}{\textbf{T1}} & Federal Reserve H.4.1: Historical Gold Stock Data (\ldots & Falsified if the gold-cover ratio was still above 80\% at the time of th\ldots \\
\eqlab{eq:\discretionary{}{}{}13.\discretionary{}{}{}16-\discretionary{}{}{}polymorphic-\discretionary{}{}{}reboot-\discretionary{}{}{}operator} & Ch.~13, l.~10310 & \textcolor{blue!70!black}{\textbf{T2}} & Gowa 1983 (Closing the Gold Window) & Falsified if Assets(E, A\_new) < Assets(E, A\_old) across any of the thre\ldots \\
\textbf{\eqlab{eq:\discretionary{}{}{}13.\discretionary{}{}{}17-\discretionary{}{}{}elite-\discretionary{}{}{}asset-\discretionary{}{}{}continuity-\discretionary{}{}{}invariant}} & Ch.~13, l.~10320 & \textcolor{green!60!black}{\textbf{T1}} & WID.world: Top-0.1\% US Wealth Share 1913-present (P\ldots & Falsified if the top-0.1\% US wealth share declines STRUCTURALLY (> 3 pe\ldots \\
\multicolumn{5}{l}{\small\bfseries Chapter 14: The Algorithmic Epoch: Real-Time Subjugation and the Necessity of the Counter-Virus} \\
\eqlab{eq:\discretionary{}{}{}14.\discretionary{}{}{}1-\discretionary{}{}{}algo-\discretionary{}{}{}prior-\discretionary{}{}{}inheritance} & Ch.~14, l.~10757 & \textcolor{blue!70!black}{\textbf{T2}} & Angwin et al. (2016) --- Machine Bias (ProPublica) & Falsified if ML systems trained on pre-2024 US data show no racial disp\ldots \\
\eqlab{eq:\discretionary{}{}{}14.\discretionary{}{}{}2-\discretionary{}{}{}realtime-\discretionary{}{}{}proximity-\discretionary{}{}{}estimate} & Ch.~14, l.~10777 & \textcolor{blue!70!black}{\textbf{T2}} & Twitter/X political affinity graph data (academic d\ldots & Falsified if real-time social-affinity graph G(t) shows no correlation \ldots \\
\eqlab{eq:\discretionary{}{}{}14.\discretionary{}{}{}3-\discretionary{}{}{}orthogonal-\discretionary{}{}{}vector-\discretionary{}{}{}injection} & Ch.~14, l.~10793 & \textcolor{gray}{\textbf{T3}} & \textit{(ordinal)} & Falsified if documented interference engine operations increase class-a\ldots \\
\eqlab{eq:\discretionary{}{}{}14.\discretionary{}{}{}4-\discretionary{}{}{}correction-\discretionary{}{}{}signal-\discretionary{}{}{}damping} & Ch.~14, l.~10807 & \textcolor{gray}{\textbf{T3}} & \textit{(ordinal)} & Falsified if scale-accurate counter-signals are shown to reduce identit\ldots \\
\eqlab{eq:\discretionary{}{}{}14.\discretionary{}{}{}5-\discretionary{}{}{}optical-\discretionary{}{}{}enclosure} & Ch.~14, l.~10820 & \textcolor{gray}{\textbf{T3}} & Pew Research --- Police-community relations by race & Falsified if documented perception studies show O's population maintain\ldots \\
\eqlab{eq:\discretionary{}{}{}14.\discretionary{}{}{}6-\discretionary{}{}{}decoy-\discretionary{}{}{}transfer-\discretionary{}{}{}coefficient} & Ch.~14, l.~10834 & \textcolor{blue!70!black}{\textbf{T2}} & \textit{(pending)} & Falsified if documented cases show kinetic class energy passing to E at\ldots \\
\eqlab{eq:\discretionary{}{}{}14.\discretionary{}{}{}7-\discretionary{}{}{}ddos-\discretionary{}{}{}bandwidth-\discretionary{}{}{}ceiling} & Ch.~14, l.~10858 & \textcolor{blue!70!black}{\textbf{T2}} & Armed Conflict Location \& Event Data (ACLED) --- US p\ldots & Falsified if F\_enforce simultaneously suppresses n > $\Gamma$ nodes in any doc\ldots \\
\eqlab{eq:\discretionary{}{}{}14.\discretionary{}{}{}8-\discretionary{}{}{}armed-\discretionary{}{}{}swarm-\discretionary{}{}{}bandwidth} & Ch.~14, l.~10869 & \textcolor{blue!70!black}{\textbf{T2}} & \textit{(pending)} & Falsified if F\_enforce simultaneously neutralizes armed nodes with forc\ldots \\
\eqlab{eq:\discretionary{}{}{}14.\discretionary{}{}{}9-\discretionary{}{}{}kinetic-\discretionary{}{}{}capital-\discretionary{}{}{}asymmetry} & Ch.~14, l.~10887 & \textcolor{blue!70!black}{\textbf{T2}} & Pew Research Center --- Gun Ownership Survey 2017 & Falsified if gun-ownership data shows no racial asymmetry in kinetic ca\ldots \\
\eqlab{eq:\discretionary{}{}{}14.\discretionary{}{}{}10-\discretionary{}{}{}sustained-\discretionary{}{}{}siege-\discretionary{}{}{}cost} & Ch.~14, l.~10909 & \textcolor{blue!70!black}{\textbf{T2}} & \textit{(pending)} & Falsified if F\_enforce besieges multiple simultaneous armed nodes with \ldots \\
\eqlab{eq:\discretionary{}{}{}14.\discretionary{}{}{}11-\discretionary{}{}{}escalation-\discretionary{}{}{}ceiling-\discretionary{}{}{}tiers} & Ch.~14, l.~10919 & \textcolor{gray}{\textbf{T3}} & \textit{(ordinal)} & Falsified if F\_enforce escalation tiers do not follow the predicted ban\ldots \\
\eqlab{eq:\discretionary{}{}{}14.\discretionary{}{}{}12-\discretionary{}{}{}inclusion-\discretionary{}{}{}predicate} & Ch.~14, l.~10948 & \textcolor{blue!70!black}{\textbf{T2}} & Bureau of Labor Statistics contingent worker survey & Falsified if data shows V(x,t) $\leq$ R(x,t) for majority members of I\_buffe\ldots \\
\eqlab{eq:\discretionary{}{}{}14.\discretionary{}{}{}13-\discretionary{}{}{}terminal-\discretionary{}{}{}interface-\discretionary{}{}{}swap} & Ch.~14, l.~10959 & \textcolor{gray}{\textbf{T3}} & \textit{(ordinal)} & Falsified if documented late-stage extraction shows I\_buffer benefits m\ldots \\
\eqlab{eq:\discretionary{}{}{}14.\discretionary{}{}{}14-\discretionary{}{}{}liability-\discretionary{}{}{}phase-\discretionary{}{}{}optimization} & Ch.~14, l.~10969 & \textcolor{gray}{\textbf{T3}} & \textit{(ordinal)} & Falsified if documented terminal-phase extraction shows E maintaining I\ldots \\
\eqlab{eq:\discretionary{}{}{}14.\discretionary{}{}{}15-\discretionary{}{}{}noise-\discretionary{}{}{}spectrum-\discretionary{}{}{}index} & Ch.~14, l.~11017 & \textcolor{blue!70!black}{\textbf{T2}} & ACLED --- US social movements 2011--2020 & Falsified if noise-spectrum index of a documented swarm fails to predic\ldots \\
\eqlab{eq:\discretionary{}{}{}14.\discretionary{}{}{}16-\discretionary{}{}{}agnostic-\discretionary{}{}{}swarm-\discretionary{}{}{}payload} & Ch.~14, l.~11036 & \textcolor{gray}{\textbf{T3}} & \textit{(ordinal)} & Falsified if documented multi-node mobilizations with ideological diver\ldots \\
\eqlab{eq:\discretionary{}{}{}14.\discretionary{}{}{}17-\discretionary{}{}{}botnet-\discretionary{}{}{}load-\discretionary{}{}{}theorem} & Ch.~14, l.~11051 & \textcolor{blue!70!black}{\textbf{T2}} & ACLED --- 2020 US protest dataset & Falsified if enforcement grid contains a swarm with n > $\Gamma$\_armed nodes a\ldots \\
\eqlab{eq:\discretionary{}{}{}14.\discretionary{}{}{}18-\discretionary{}{}{}ideological-\discretionary{}{}{}routing-\discretionary{}{}{}operator} & Ch.~14, l.~11072 & \textcolor{gray}{\textbf{T3}} & \textit{(ordinal)} & Falsified if nodes in documented distributed mobilizations fail to rout\ldots \\
\eqlab{eq:\discretionary{}{}{}14.\discretionary{}{}{}19-\discretionary{}{}{}zugzwang-\discretionary{}{}{}payoff-\discretionary{}{}{}function} & Ch.~14, l.~11122 & \textcolor{blue!70!black}{\textbf{T2}} & \textit{(pending)} & Falsified if Elite actors systematically choose attack over stand-down \ldots \\
\eqlab{eq:\discretionary{}{}{}14.\discretionary{}{}{}20-\discretionary{}{}{}defection-\discretionary{}{}{}cascade} & Ch.~14, l.~11164 & \textcolor{blue!70!black}{\textbf{T2}} & James (1938) --- The Black Jacobins (Polish Legion de\ldots & Falsified if documented enforcement nodes fail to defect when community\ldots \\
\eqlab{eq:\discretionary{}{}{}14.\discretionary{}{}{}21-\discretionary{}{}{}empathy-\discretionary{}{}{}permeability} & Ch.~14, l.~11196 & \textcolor{blue!70!black}{\textbf{T2}} & James (1938) --- The Black Jacobins & Falsified if trauma overlap between enforcement and out-group nodes fai\ldots \\
\eqlab{eq:\discretionary{}{}{}14.\discretionary{}{}{}22-\discretionary{}{}{}semantic-\discretionary{}{}{}overwrite} & Ch.~14, l.~11222 & \textcolor{gray}{\textbf{T3}} & Haitian Constitution of 1805 (primary source) & Falsified if 1805 Haitian Constitution's ideological redefinition faile\ldots \\
\multicolumn{5}{l}{\small\bfseries Chapter 17: Conclusion} \\
\eqlab{eq:\discretionary{}{}{}17.\discretionary{}{}{}1-\discretionary{}{}{}racism-\discretionary{}{}{}as-\discretionary{}{}{}institutional-\discretionary{}{}{}vector} & Ch.~17, l.~11531 & \textcolor{gray}{\textbf{T3}} & \textit{(ordinal)} & Falsified if documented cases of marginalized-group prejudice achieve s\ldots \\
\eqlab{eq:\discretionary{}{}{}17.\discretionary{}{}{}2-\discretionary{}{}{}o-\discretionary{}{}{}terminal-\discretionary{}{}{}expansion} & Ch.~17, l.~11586 & \textcolor{blue!70!black}{\textbf{T2}} & US Census Bureau 2020 demographic projections & Falsified if documented demographic trends show O contracting rather th\ldots \\
\multicolumn{5}{l}{\small\bfseries Chapter 19: Equation Registry and Era-Level Calibration} \\
\eqlab{eq:\discretionary{}{}{}19.\discretionary{}{}{}1-\discretionary{}{}{}kernel-\discretionary{}{}{}objective-\discretionary{}{}{}registry-\discretionary{}{}{}entry} & Ch.~19, l.~11717 & \textcolor{gray}{\textbf{T3}} & \textit{(ordinal)} & Redundant with eq:5 and eq:56; falsified by the same conditions as thos\ldots \\
\eqlab{eq:\discretionary{}{}{}19.\discretionary{}{}{}2-\discretionary{}{}{}suppression-\discretionary{}{}{}envelope-\discretionary{}{}{}registry-\discretionary{}{}{}entry} & Ch.~19, l.~11721 & \textcolor{gray}{\textbf{T3}} & \textit{(ordinal)} & Redundant with eq:8 and eq:58; falsified by the same conditions as thos\ldots \\
\eqlab{eq:\discretionary{}{}{}19.\discretionary{}{}{}3-\discretionary{}{}{}crash-\discretionary{}{}{}condition-\discretionary{}{}{}registry-\discretionary{}{}{}entry} & Ch.~19, l.~11725 & \textcolor{gray}{\textbf{T3}} & \textit{(ordinal)} & Redundant with eq:9 and eq:59; falsified by the same conditions as thos\ldots \\
\eqlab{eq:\discretionary{}{}{}19.\discretionary{}{}{}4-\discretionary{}{}{}tier-\discretionary{}{}{}ordering-\discretionary{}{}{}registry-\discretionary{}{}{}entry} & Ch.~19, l.~11729 & \textcolor{gray}{\textbf{T3}} & \textit{(ordinal)} & Redundant with eq:60; falsified if any tier ordering is reversed in com\ldots \\
\end{longtable}
\endgroup

%
% Requires: longtable, booktabs, xcolor, array (loaded in main preamble)

\begingroup
\small
\setlength{\LTpre}{6pt}
\setlength{\LTpost}{6pt}
% Equation labels are long unbreakable \texttt strings (up to 47 chars). Without
% explicit break opportunities they overrun the label column and print on top of
% the case-study column. \eqlab emits the label with zero-width discretionary
% breaks after ':', '-' and '.', inserted by the generator.
\providecommand{\eqlab}[1]{\texttt{#1}}
\begin{longtable}{@{}%
  >{\raggedright\arraybackslash}p{3.2cm}
  >{\raggedright\arraybackslash}p{2.5cm}
  >{\centering\arraybackslash}p{1.0cm}
  >{\raggedright\arraybackslash}p{3.3cm}
  >{\raggedright\arraybackslash}p{5.2cm}@{}}
\textbf{Equation label} &
\textbf{Case study location} &
\textbf{Tier} &
\textbf{Primary data source} &
\textbf{Falsification (truncated)} \\
\endfirsthead
\textbf{Equation label} &
\textbf{Case study location} &
\textbf{Tier} &
\textbf{Primary data source} &
\textbf{Falsification (truncated)} \\
\endhead
\multicolumn{5}{r}{\small\itshape continued on next page} \\
\endfoot
\endlastfoot
\multicolumn{5}{l}{\small\bfseries Chapter 1: Redefining Racism} \\
\eqlab{eq:\discretionary{}{}{}1.\discretionary{}{}{}1-\discretionary{}{}{}enclosure-\discretionary{}{}{}score} & Ch.~1, l.~241 & \textcolor{gray}{\textbf{T3}} & \textit{(ordinal)} & Falsified if a population with all three outlets blocked (e\_i=1) shows \ldots \\
\eqlab{eq:\discretionary{}{}{}1.\discretionary{}{}{}2-\discretionary{}{}{}total-\discretionary{}{}{}enclosure} & Ch.~1, l.~257 & \textcolor{gray}{\textbf{T3}} & \textit{(ordinal)} & Falsified if a fully blocked population (e\_1=e\_2=e\_3=1) achieves coordi\ldots \\
\eqlab{eq:\discretionary{}{}{}1.\discretionary{}{}{}3-\discretionary{}{}{}conventional-\discretionary{}{}{}causal-\discretionary{}{}{}arrow} & Ch.~1, l.~291 & \textcolor{gray}{\textbf{T3}} & \textit{(ordinal)} & Falsified if documented evidence shows individual prejudice predating E\ldots \\
\eqlab{eq:\discretionary{}{}{}1.\discretionary{}{}{}4-\discretionary{}{}{}elite-\discretionary{}{}{}causal-\discretionary{}{}{}arrow} & Ch.~1, l.~297 & \textcolor{gray}{\textbf{T3}} & Zurara & Falsified if racialization systems are shown to emerge without Elite ec\ldots \\
\textbf{\eqlab{eq:\discretionary{}{}{}1.\discretionary{}{}{}5-\discretionary{}{}{}kernel-\discretionary{}{}{}optimization}} & Ch.~1, l.~495 & \textcolor{green!60!black}{\textbf{T1}} & Piketty & Falsified if a documented period shows sustained decline in Elite extra\ldots \\
\eqlab{eq:\discretionary{}{}{}1.\discretionary{}{}{}6-\discretionary{}{}{}interface-\discretionary{}{}{}strategy-\discretionary{}{}{}set} & Ch.~1, l.~576 & \textcolor{gray}{\textbf{T3}} & \textit{(ordinal)} & Falsified if a documented Elite extraction system used a strategy outsi\ldots \\
\eqlab{eq:\discretionary{}{}{}1.\discretionary{}{}{}7-\discretionary{}{}{}interface-\discretionary{}{}{}optimizer} & Ch.~1, l.~608 & \textcolor{blue!70!black}{\textbf{T2}} & Dudziak (2000) --- Cold War Civil Rights & Falsified if a historical interface transition increased net cost C\_tot\ldots \\
\textbf{\eqlab{eq:\discretionary{}{}{}1.\discretionary{}{}{}8-\discretionary{}{}{}suppression-\discretionary{}{}{}envelope}} & Ch.~1, l.~714 & \textcolor{blue!70!black}{\textbf{T2}} & BLS Union Membership Historical Series & Falsified if M(t) consistently exceeds $\Sigma$\_sup(t) for extended periods wi\ldots \\
\eqlab{eq:\discretionary{}{}{}1.\discretionary{}{}{}9-\discretionary{}{}{}crash-\discretionary{}{}{}condition-\discretionary{}{}{}derivative} & Ch.~1, l.~710 & \textcolor{gray}{\textbf{T3}} & \textit{(ordinal)} & Falsified if a crash event (Bacon's Rebellion, Haitian Revolution) occu\ldots \\
\eqlab{eq:\discretionary{}{}{}1.\discretionary{}{}{}10-\discretionary{}{}{}effective-\discretionary{}{}{}coherence} & Ch.~1, l.~716 & \textcolor{blue!70!black}{\textbf{T2}} & ANES cross-racial coalition data 1948--2020 & Falsified if a crash event occurs when M\_eff(t) < $\tau$, or if no crash occ\ldots \\
\eqlab{eq:\discretionary{}{}{}1.\discretionary{}{}{}11-\discretionary{}{}{}phase-\discretionary{}{}{}loading} & Ch.~1, l.~838 & \textcolor{gray}{\textbf{T3}} & ANES Time Series Study 1948--2020 & Falsified if survey-measured cross-group solidarity remains high when $\Phi$\ldots \\
\multicolumn{5}{l}{\small\bfseries Chapter 2: Version 1.0: Initializing the Vector (15th-Century Portugal)} \\
\eqlab{eq:\discretionary{}{}{}2.\discretionary{}{}{}1-\discretionary{}{}{}racism-\discretionary{}{}{}vector} & Ch.~2, l.~1091 & \textcolor{gray}{\textbf{T3}} & \textit{(ordinal)} & Falsified if documented racism operates without a state-power direction\ldots \\
\eqlab{eq:\discretionary{}{}{}2.\discretionary{}{}{}2-\discretionary{}{}{}status-\discretionary{}{}{}suppression-\discretionary{}{}{}allocation} & Ch.~2, l.~1163 & \textcolor{gray}{\textbf{T3}} & Du Bois (1935) --- Black Reconstruction in America & Falsified if Buffer Class loyalty is maintained without a non-material \ldots \\
\eqlab{eq:\discretionary{}{}{}2.\discretionary{}{}{}3-\discretionary{}{}{}psi-\discretionary{}{}{}null-\discretionary{}{}{}in-\discretionary{}{}{}survival} & Ch.~2, l.~1167 & \textcolor{gray}{\textbf{T3}} & \textit{(ordinal)} & Falsified if material wages to Buffer Class are shown to be independent\ldots \\
\eqlab{eq:\discretionary{}{}{}2.\discretionary{}{}{}4-\discretionary{}{}{}roman-\discretionary{}{}{}membership-\discretionary{}{}{}function} & Ch.~2, l.~1208 & \textcolor{gray}{\textbf{T3}} & Bradley (1994) --- Slavery and Society at Rome & Falsified if Roman enslaved population is shown to be phenotypically ho\ldots \\
\eqlab{eq:\discretionary{}{}{}2.\discretionary{}{}{}5-\discretionary{}{}{}american-\discretionary{}{}{}racial-\discretionary{}{}{}function} & Ch.~2, l.~1217 & \textcolor{gray}{\textbf{T3}} & Hening's Statutes at Large (Virginia) Vol. 2 & Falsified if documented American slavery assigned status based on non-p\ldots \\
\eqlab{eq:\discretionary{}{}{}2.\discretionary{}{}{}6-\discretionary{}{}{}roman-\discretionary{}{}{}american-\discretionary{}{}{}structural-\discretionary{}{}{}parallel} & Ch.~2, l.~1235 & \textcolor{gray}{\textbf{T3}} & \textit{(ordinal)} & Falsified if the structural mapping between Roman and American slavery \ldots \\
\eqlab{eq:\discretionary{}{}{}2.\discretionary{}{}{}7-\discretionary{}{}{}three-\discretionary{}{}{}tier-\discretionary{}{}{}portugal-\discretionary{}{}{}inequality} & Ch.~2, l.~1280 & \textcolor{gray}{\textbf{T3}} & \textit{(ordinal)} & Falsified if measured material outcomes show Benefit(I\_buffer) $\geq$ Benefi\ldots \\
\eqlab{eq:\discretionary{}{}{}2.\discretionary{}{}{}8-\discretionary{}{}{}church-\discretionary{}{}{}science-\discretionary{}{}{}feedback-\discretionary{}{}{}loop} & Ch.~2, l.~1335 & \textcolor{gray}{\textbf{T3}} & Cartwright (1851) --- Diseases and Peculiarities of t\ldots & Falsified if the religious-scientific legitimation loop is disrupted fr\ldots \\
\multicolumn{5}{l}{\small\bfseries Chapter 3: The Application: Bacon's Rebellion, the Buffer Class, and the Constitutional Patch (1676--1787)} \\
\textbf{\eqlab{eq:\discretionary{}{}{}3.\discretionary{}{}{}1-\discretionary{}{}{}bacon-\discretionary{}{}{}solidarity-\discretionary{}{}{}condition}} & Ch.~3, l.~1580 & \textcolor{green!60!black}{\textbf{T1}} & Morgan (1975) --- American Slavery & Falsified if an antebellum cross-racial coalition reached M(t) > $\tau$ with\ldots \\
\eqlab{eq:\discretionary{}{}{}3.\discretionary{}{}{}2-\discretionary{}{}{}boundary-\discretionary{}{}{}enforcement-\discretionary{}{}{}1691} & Ch.~3, l.~1507 & \textcolor{gray}{\textbf{T3}} & Hening's Statutes at Large (Virginia) Vol. 3 & Falsified if racial boundary enforcement persisted without legal penalt\ldots \\
\eqlab{eq:\discretionary{}{}{}3.\discretionary{}{}{}3-\discretionary{}{}{}race-\discretionary{}{}{}destroys-\discretionary{}{}{}class} & Ch.~3, l.~1539 & \textcolor{gray}{\textbf{T3}} & Morgan (1975) --- American Slavery & Falsified if colonial-era racial coding produced class solidarity rathe\ldots \\
\eqlab{eq:\discretionary{}{}{}3.\discretionary{}{}{}4-\discretionary{}{}{}three-\discretionary{}{}{}tier-\discretionary{}{}{}post-\discretionary{}{}{}bacon} & Ch.~3, l.~1560 & \textcolor{gray}{\textbf{T3}} & \textit{(ordinal)} & Falsified if the post-Bacon ordering shows Buffer Class benefits exceed\ldots \\
\eqlab{eq:\discretionary{}{}{}3.\discretionary{}{}{}5-\discretionary{}{}{}kinetic-\discretionary{}{}{}necessary-\discretionary{}{}{}condition} & Ch.~3, l.~1595 & \textcolor{green!60!black}{\textbf{T1}} & Morgan (1975) --- American Slavery & Falsified if documented cases show Elite extraction maintained when K\_E\ldots \\
\multicolumn{5}{l}{\small\bfseries Chapter 5: The Gendered Axis: Coverture, Eugenics, and the Reproductive Extraction Kernel} \\
\eqlab{eq:\discretionary{}{}{}5.\discretionary{}{}{}1-\discretionary{}{}{}compliance-\discretionary{}{}{}differential} & Ch.~5, l.~2272 & \textcolor{blue!70!black}{\textbf{T2}} & Baptist (2014) --- The Half Has Never Been Told & Falsified if V\_c(x) $\leq$ V\_r(x) for subjects in documented extraction regi\ldots \\
\multicolumn{5}{l}{\small\bfseries Chapter 6: The Enforcement Engine: Slave Patrols, the 13th Amendment, and the Compounding Model (1704--1865)} \\
\eqlab{eq:\discretionary{}{}{}6.\discretionary{}{}{}1-\discretionary{}{}{}lethal-\discretionary{}{}{}autonomy-\discretionary{}{}{}gradient} & Ch.~6, l.~2889 & \textcolor{blue!70!black}{\textbf{T2}} & LEOSA --- 18 U.S.C. §926B statutory text & Falsified if Mapping Police Violence data shows parity in police killin\ldots \\
\textbf{\eqlab{eq:\discretionary{}{}{}6.\discretionary{}{}{}2-\discretionary{}{}{}four-\discretionary{}{}{}tier-\discretionary{}{}{}benefit-\discretionary{}{}{}ordering}} & Ch.~6, l.~2994 & \textcolor{gray}{\textbf{T3}} & Piketty-Saez-Zucman Distributional National Accounts & Falsified if comprehensive wealth, health, or carceral data shows Benef\ldots \\
\eqlab{eq:\discretionary{}{}{}6.\discretionary{}{}{}3-\discretionary{}{}{}thirteenth-\discretionary{}{}{}amendment-\discretionary{}{}{}exception} & Ch.~6, l.~3080 & \textcolor{gray}{\textbf{T3}} & Blackmon (2008) --- Slavery by Another Name & Falsified if 13th Amendment exception clause is shown to have been appl\ldots \\
\eqlab{eq:\discretionary{}{}{}6.\discretionary{}{}{}4-\discretionary{}{}{}additive-\discretionary{}{}{}policy-\discretionary{}{}{}model} & Ch.~6, l.~3144 & \textcolor{gray}{\textbf{T3}} & \textit{(ordinal)} & Falsified if additive accumulation predicts compound policy outcomes wi\ldots \\
\eqlab{eq:\discretionary{}{}{}6.\discretionary{}{}{}5-\discretionary{}{}{}compounding-\discretionary{}{}{}temporal-\discretionary{}{}{}model} & Ch.~6, l.~3154 & \textcolor{blue!70!black}{\textbf{T2}} & Hamilton \& Darity (2017) --- The Political Economy of\ldots & Falsified if policy shock effects on O\textasciicircum{}capacity are shown to be indepen\ldots \\
\eqlab{eq:\discretionary{}{}{}6.\discretionary{}{}{}6-\discretionary{}{}{}asymmetric-\discretionary{}{}{}enforcement-\discretionary{}{}{}multiplier} & Ch.~6, l.~3172 & \textcolor{green!60!black}{\textbf{T1}} & ACLU (2020) --- A Tale of Two Countries: Racially Tar\ldots & Falsified if ACLU cannabis arrest data shows no racial disparity in pol\ldots \\
\eqlab{eq:\discretionary{}{}{}6.\discretionary{}{}{}7-\discretionary{}{}{}policy-\discretionary{}{}{}sequence-\discretionary{}{}{}noncommutative} & Ch.~6, l.~3177 & \textcolor{gray}{\textbf{T3}} & \textit{(ordinal)} & Falsified if order of policy application (P\_1 then P\_2 vs. P\_2 then P\_1\ldots \\
\eqlab{eq:\discretionary{}{}{}6.\discretionary{}{}{}8-\discretionary{}{}{}capacity-\discretionary{}{}{}chain-\discretionary{}{}{}1619} & Ch.~6, l.~3193 & \textcolor{blue!70!black}{\textbf{T2}} & Darity \& Mullen (2020) --- From Here to Equality (rep\ldots & Falsified if O\_1619 capacity is non-significantly different from O\_1450\ldots \\
\eqlab{eq:\discretionary{}{}{}6.\discretionary{}{}{}9-\discretionary{}{}{}capacity-\discretionary{}{}{}chain-\discretionary{}{}{}1865} & Ch.~6, l.~3196 & \textcolor{blue!70!black}{\textbf{T2}} & Blackmon (2008) --- Slavery by Another Name & Falsified if O\_1865 capacity is not reduced from O\_1619 as measured by \ldots \\
\eqlab{eq:\discretionary{}{}{}6.\discretionary{}{}{}10-\discretionary{}{}{}capacity-\discretionary{}{}{}chain-\discretionary{}{}{}1934} & Ch.~6, l.~3199 & \textcolor{green!60!black}{\textbf{T1}} & Mapping Inequality --- HOLC Redlining Maps & Falsified if redlining-exposed tracts show no differential wealth accum\ldots \\
\eqlab{eq:\discretionary{}{}{}6.\discretionary{}{}{}11-\discretionary{}{}{}capacity-\discretionary{}{}{}chain-\discretionary{}{}{}1971} & Ch.~6, l.~3230 & \textcolor{green!60!black}{\textbf{T1}} & ACLU (2020) --- Cannabis Arrests Report & Falsified if War on Drugs enforcement shows no differential impact on O\ldots \\
\textbf{\eqlab{eq:\discretionary{}{}{}6.\discretionary{}{}{}12-\discretionary{}{}{}capacity-\discretionary{}{}{}compounding-\discretionary{}{}{}full}} & Ch.~6, l.~3169 & \textcolor{green!60!black}{\textbf{T1}} & ACLU (2020) --- A Tale of Two Countries & Falsified if any one policy shock ($\alpha$, $\beta$, $\gamma$, $\delta$) is shown to have zero ma\ldots \\
\multicolumn{5}{l}{\small\bfseries Chapter 7: The Containment: Pullman, Redlining, and the Wages of Whiteness (1894--1965)} \\
\eqlab{eq:\discretionary{}{}{}7.\discretionary{}{}{}1-\discretionary{}{}{}pullman-\discretionary{}{}{}corollary} & Ch.~7, l.~3624 & \textcolor{gray}{\textbf{T3}} & Roediger (1991) --- The Wages of Whiteness & Falsified if a documented case shows I\_buffer achieving collective gain\ldots \\
\multicolumn{5}{l}{\small\bfseries Chapter 8: Tweedism and the Puppet Class: The Algorithmic Filter on Democracy} \\
\eqlab{eq:\discretionary{}{}{}8.\discretionary{}{}{}1-\discretionary{}{}{}vertex-\discretionary{}{}{}partition-\discretionary{}{}{}graph} & Ch.~8, l.~4020 & \textcolor{gray}{\textbf{T3}} & \textit{(ordinal)} & Falsified if leader/follower partition fails to predict convergence in \ldots \\
\eqlab{eq:\discretionary{}{}{}8.\discretionary{}{}{}2-\discretionary{}{}{}tier-\discretionary{}{}{}to-\discretionary{}{}{}graph-\discretionary{}{}{}mapping} & Ch.~8, l.~4024 & \textcolor{gray}{\textbf{T3}} & \textit{(ordinal)} & Falsified if the tier-to-graph mapping produces inconsistent prediction\ldots \\
\eqlab{eq:\discretionary{}{}{}8.\discretionary{}{}{}3-\discretionary{}{}{}fractional-\discretionary{}{}{}agent-\discretionary{}{}{}dynamics} & Ch.~8, l.~4032 & \textcolor{gray}{\textbf{T3}} & Li et al. --- Mittag-Leffler stability results (mathe\ldots & Falsified if integer-order dynamics ($\alpha$=1) predict social-mobility conve\ldots \\
\eqlab{eq:\discretionary{}{}{}8.\discretionary{}{}{}4-\discretionary{}{}{}compounding-\discretionary{}{}{}chain-\discretionary{}{}{}formal} & Ch.~8, l.~3187 & \textcolor{blue!70!black}{\textbf{T2}} & Federal Reserve SCF racial wealth data & Falsified if multiplicative chain prediction diverges significantly fro\ldots \\
\eqlab{eq:\discretionary{}{}{}8.\discretionary{}{}{}5-\discretionary{}{}{}stationary-\discretionary{}{}{}leader-\discretionary{}{}{}condition} & Ch.~8, l.~4046 & \textcolor{gray}{\textbf{T3}} & \textit{(ordinal)} & Falsified if documented Elite actors show systematic response to u\_i (u\ldots \\
\eqlab{eq:\discretionary{}{}{}8.\discretionary{}{}{}6-\discretionary{}{}{}proximity-\discretionary{}{}{}connectivity-\discretionary{}{}{}condition} & Ch.~8, l.~4128 & \textcolor{gray}{\textbf{T3}} & \textit{(ordinal)} & Falsified if social-affinity graph connectivity is shown to be independ\ldots \\
\eqlab{eq:\discretionary{}{}{}8.\discretionary{}{}{}7-\discretionary{}{}{}navigation-\discretionary{}{}{}components} & Ch.~8, l.~4220 & \textcolor{gray}{\textbf{T3}} & \textit{(ordinal)} & Falsified if goal and constraint functions do not bound follower trajec\ldots \\
\eqlab{eq:\discretionary{}{}{}8.\discretionary{}{}{}8-\discretionary{}{}{}navigation-\discretionary{}{}{}function} & Ch.~8, l.~4223 & \textcolor{gray}{\textbf{T3}} & \textit{(ordinal)} & Falsified if $\phi$\_i fails to produce convergence in any graph satisfying t\ldots \\
\eqlab{eq:\discretionary{}{}{}8.\discretionary{}{}{}9-\discretionary{}{}{}gradient-\discretionary{}{}{}control-\discretionary{}{}{}law} & Ch.~8, l.~4228 & \textcolor{gray}{\textbf{T3}} & \textit{(ordinal)} & Falsified if negative gradient of $\phi$\_i fails to drive convergence in any\ldots \\
\eqlab{eq:\discretionary{}{}{}8.\discretionary{}{}{}10-\discretionary{}{}{}convex-\discretionary{}{}{}hull-\discretionary{}{}{}convergence} & Ch.~8, l.~4316 & \textcolor{gray}{\textbf{T3}} & \textit{(ordinal)} & Falsified if followers fail to converge to Elite convex hull in any gra\ldots \\
\eqlab{eq:\discretionary{}{}{}8.\discretionary{}{}{}11-\discretionary{}{}{}mittag-\discretionary{}{}{}leffler-\discretionary{}{}{}stability} & Ch.~8, l.~4402 & \textcolor{gray}{\textbf{T3}} & Chetty et al. (2018) --- Race and Economic Opportunit\ldots & Falsified if social-mobility data for $\alpha$ $\in$ (0,1) shows exponential rathe\ldots \\
\textbf{\eqlab{eq:\discretionary{}{}{}8.\discretionary{}{}{}12-\discretionary{}{}{}class-\discretionary{}{}{}alignment-\discretionary{}{}{}base-\discretionary{}{}{}waves}} & Ch.~8, l.~4968 & \textcolor{green!60!black}{\textbf{T1}} & Google Trends & Falsified if class alignment signal is not measurable as a distinct fre\ldots \\
\eqlab{eq:\discretionary{}{}{}8.\discretionary{}{}{}13-\discretionary{}{}{}subgroup-\discretionary{}{}{}compound-\discretionary{}{}{}phase} & Ch.~8, l.~4968 & \textcolor{blue!70!black}{\textbf{T2}} & Google Trends & Falsified if compound phase $\Phi$\_j fails to predict subgroup political ali\ldots \\
\eqlab{eq:\discretionary{}{}{}8.\discretionary{}{}{}14-\discretionary{}{}{}net-\discretionary{}{}{}solidarity-\discretionary{}{}{}signal} & Ch.~8, l.~4968 & \textcolor{green!60!black}{\textbf{T1}} & Google Trends & Falsified if net solidarity signal cannot be measured as distinct from \ldots \\
\eqlab{eq:\discretionary{}{}{}8.\discretionary{}{}{}15-\discretionary{}{}{}solidarity-\discretionary{}{}{}collapse-\discretionary{}{}{}condition} & Ch.~8, l.~4968 & \textcolor{green!60!black}{\textbf{T1}} & Google Trends & Falsified if M(t) remains above $\tau$ even when S\_class approaches zero acr\ldots \\
\eqlab{eq:\discretionary{}{}{}8.\discretionary{}{}{}16-\discretionary{}{}{}interference-\discretionary{}{}{}control-\discretionary{}{}{}objective} & Ch.~8, l.~4968 & \textcolor{green!60!black}{\textbf{T1}} & Google Trends & Falsified if observed interference engine output increases S\_class rath\ldots \\
\eqlab{eq:\discretionary{}{}{}8.\discretionary{}{}{}17-\discretionary{}{}{}circular-\discretionary{}{}{}dispersion-\discretionary{}{}{}operator} & Ch.~8, l.~4968 & \textcolor{blue!70!black}{\textbf{T2}} & ANES Time Series --- cross-group solidarity items & Falsified if circular dispersion of phase values fails to predict cross\ldots \\
\textbf{\eqlab{eq:\discretionary{}{}{}8.\discretionary{}{}{}18-\discretionary{}{}{}tweedism-\discretionary{}{}{}agenda-\discretionary{}{}{}path}} & Ch.~8, l.~4844 & \textcolor{green!60!black}{\textbf{T1}} & Gilens \& Page (2014) --- Testing Theories of American\ldots & Falsified if Gilens-Page analysis shows median voter preferences predic\ldots \\
\multicolumn{5}{l}{\small\bfseries Chapter 9: The Recompile: COINTELPRO, the Variable Swap, and the War on Drugs (1968--1994)} \\
\textbf{\eqlab{eq:\discretionary{}{}{}9.\discretionary{}{}{}1-\discretionary{}{}{}lead-\discretionary{}{}{}crime-\discretionary{}{}{}compounding}} & Ch.~9, l.~6344 & \textcolor{green!60!black}{\textbf{T1}} & Reyes (2007) --- Environmental Policy as Social Policy & Falsified if Reyes (2007) blood-lead time-series fails to predict crime\ldots \\
\eqlab{eq:\discretionary{}{}{}9.\discretionary{}{}{}2-\discretionary{}{}{}epistemic-\discretionary{}{}{}suppression-\discretionary{}{}{}variable} & Ch.~9, l.~6344 & \textcolor{blue!70!black}{\textbf{T2}} & Kitman (2000) --- The Secret History of Lead (The Nat\ldots & Falsified if documented lead industry cover-up timeline shows no correl\ldots \\
\eqlab{eq:\discretionary{}{}{}9.\discretionary{}{}{}3-\discretionary{}{}{}multi-\discretionary{}{}{}vector-\discretionary{}{}{}lead-\discretionary{}{}{}exposure} & Ch.~9, l.~6344 & \textcolor{green!60!black}{\textbf{T1}} & EPA Air Quality Monitoring Data & Falsified if cumulative lead exposure from all three vectors is not sig\ldots \\
\eqlab{eq:\discretionary{}{}{}9.\discretionary{}{}{}4-\discretionary{}{}{}redlining-\discretionary{}{}{}containment-\discretionary{}{}{}field} & Ch.~9, l.~6344 & \textcolor{green!60!black}{\textbf{T1}} & Mapping Inequality HOLC maps & Falsified if HOLC map grades show no statistically significant correlat\ldots \\
\eqlab{eq:\discretionary{}{}{}9.\discretionary{}{}{}5-\discretionary{}{}{}school-\discretionary{}{}{}funding-\discretionary{}{}{}property-\discretionary{}{}{}value} & Ch.~9, l.~6337 & \textcolor{green!60!black}{\textbf{T1}} & NCES Public School Finance Survey & Falsified if school funding is shown to be independent of property valu\ldots \\
\eqlab{eq:\discretionary{}{}{}9.\discretionary{}{}{}6-\discretionary{}{}{}infrastructure-\discretionary{}{}{}quality-\discretionary{}{}{}funding} & Ch.~9, l.~6337 & \textcolor{blue!70!black}{\textbf{T2}} & NCES School Facility Condition Survey & Falsified if school infrastructure quality is shown to be unrelated to \ldots \\
\eqlab{eq:\discretionary{}{}{}9.\discretionary{}{}{}7-\discretionary{}{}{}school-\discretionary{}{}{}lead-\discretionary{}{}{}exposure-\discretionary{}{}{}inverse} & Ch.~9, l.~6337 & \textcolor{green!60!black}{\textbf{T1}} & EPA 3Ts for Reducing Lead in Drinking Water in Scho\ldots & Falsified if school infrastructure quality is shown to be unrelated to \ldots \\
\eqlab{eq:\discretionary{}{}{}9.\discretionary{}{}{}8-\discretionary{}{}{}community-\discretionary{}{}{}capacity-\discretionary{}{}{}lead-\discretionary{}{}{}reduction} & Ch.~9, l.~6337 & \textcolor{blue!70!black}{\textbf{T2}} & Aizer \& Currie (2019) --- Lead and Juvenile Delinquen\ldots & Falsified if community economic capacity is shown to be independent of \ldots \\
\eqlab{eq:\discretionary{}{}{}9.\discretionary{}{}{}9-\discretionary{}{}{}property-\discretionary{}{}{}value-\discretionary{}{}{}capacity-\discretionary{}{}{}feedback} & Ch.~9, l.~6337 & \textcolor{green!60!black}{\textbf{T1}} & Mapping Inequality HOLC maps + ACS property value d\ldots & Falsified if property values in subsequent periods are shown to be inde\ldots \\
\eqlab{eq:\discretionary{}{}{}9.\discretionary{}{}{}10-\discretionary{}{}{}highway-\discretionary{}{}{}lead-\discretionary{}{}{}spatial-\discretionary{}{}{}concentration} & Ch.~9, l.~7178 & \textcolor{green!60!black}{\textbf{T1}} & Mapping Inequality DSL (Nelson et al.) --- HOLC grade\ldots & Falsified if HOLC-D tracts show no excess firearm or lead burden relati\ldots \\
\multicolumn{5}{l}{\small\bfseries Chapter 10: The Full Algorithm: Demographic Paradox, Cannibalization, and the 5-Tier Reveal (1994--Present)} \\
\eqlab{eq:\discretionary{}{}{}10.\discretionary{}{}{}1-\discretionary{}{}{}outgroup-\discretionary{}{}{}racialized-\discretionary{}{}{}subset} & Ch.~10, l.~7051 & \textcolor{gray}{\textbf{T3}} & US Census Bureau demographic data & Falsified if O\_racialized is shown to be coextensive with O rather than\ldots \\
\eqlab{eq:\discretionary{}{}{}10.\discretionary{}{}{}2-\discretionary{}{}{}buffer-\discretionary{}{}{}class-\discretionary{}{}{}shrinkage} & Ch.~10, l.~7055 & \textcolor{blue!70!black}{\textbf{T2}} & Pew Research Center --- America's Shrinking Middle Cl\ldots & Falsified if Buffer Class size is shown to be expanding rather than con\ldots \\
\eqlab{eq:\discretionary{}{}{}10.\discretionary{}{}{}3-\discretionary{}{}{}cannibalization-\discretionary{}{}{}equation} & Ch.~10, l.~7153 & \textcolor{gray}{\textbf{T3}} & Saez-Zucman wealth share series & Falsified if documented Elite extraction rate increases without corresp\ldots \\
\eqlab{eq:\discretionary{}{}{}10.\discretionary{}{}{}4-\discretionary{}{}{}predatory-\discretionary{}{}{}minmax-\discretionary{}{}{}full-\discretionary{}{}{}definition} & Ch.~10, l.~7303 & \textcolor{gray}{\textbf{T3}} & \textit{(ordinal)} & Falsified if the five-tier hierarchy fails to map onto any contemporary\ldots \\
\eqlab{eq:\discretionary{}{}{}10.\discretionary{}{}{}5-\discretionary{}{}{}kernel-\discretionary{}{}{}objective-\discretionary{}{}{}canonical} & Ch.~10, l.~7477 & \textcolor{green!60!black}{\textbf{T1}} & Piketty-Saez-Zucman Distributional National Accounts & Falsified if a documented period shows sustained Elite wealth-share dec\ldots \\
\eqlab{eq:\discretionary{}{}{}10.\discretionary{}{}{}6-\discretionary{}{}{}effective-\discretionary{}{}{}resistance-\discretionary{}{}{}variable} & Ch.~10, l.~7309 & \textcolor{gray}{\textbf{T3}} & \textit{(ordinal)} & Falsified if M\_eff fails to capture the difference between raw class co\ldots \\
\eqlab{eq:\discretionary{}{}{}10.\discretionary{}{}{}7-\discretionary{}{}{}suppression-\discretionary{}{}{}envelope-\discretionary{}{}{}canonical} & Ch.~10, l.~7345 & \textcolor{blue!70!black}{\textbf{T2}} & \textit{(pending)} & Falsified if documented suppression tools outside this set constitute t\ldots \\
\eqlab{eq:\discretionary{}{}{}10.\discretionary{}{}{}8-\discretionary{}{}{}crash-\discretionary{}{}{}condition-\discretionary{}{}{}canonical} & Ch.~10, l.~7319 & \textcolor{gray}{\textbf{T3}} & \textit{(ordinal)} & Falsified if crash events occur in documented cases without M\_eff(t) > \ldots \\
\eqlab{eq:\discretionary{}{}{}10.\discretionary{}{}{}9-\discretionary{}{}{}tier-\discretionary{}{}{}benefit-\discretionary{}{}{}ordering-\discretionary{}{}{}canonical} & Ch.~10, l.~7324 & \textcolor{gray}{\textbf{T3}} & \textit{(ordinal)} & Falsified if any two-tier comparison fails to show stated ordering in c\ldots \\
\eqlab{eq:\discretionary{}{}{}10.\discretionary{}{}{}10-\discretionary{}{}{}system-\discretionary{}{}{}objective-\discretionary{}{}{}extraction-\discretionary{}{}{}ratio} & Ch.~10, l.~7405 & \textcolor{gray}{\textbf{T3}} & \textit{(ordinal)} & Falsified if documented extraction systems show this ratio declining wi\ldots \\
\eqlab{eq:\discretionary{}{}{}10.\discretionary{}{}{}11-\discretionary{}{}{}system-\discretionary{}{}{}stability-\discretionary{}{}{}algorithm} & Ch.~10, l.~7405 & \textcolor{gray}{\textbf{T3}} & \textit{(ordinal)} & Falsified if system stability is maintained when extraction/resistance \ldots \\
\textbf{\eqlab{eq:\discretionary{}{}{}10.\discretionary{}{}{}12-\discretionary{}{}{}o-\discretionary{}{}{}final-\discretionary{}{}{}construction}} & Ch.~10, l.~6828 & \textcolor{green!60!black}{\textbf{T1}} & Bureau of Justice Statistics incarceration by race \ldots & Falsified if modern mass-incarceration demographics fail to reproduce t\ldots \\
\eqlab{eq:\discretionary{}{}{}10.\discretionary{}{}{}13-\discretionary{}{}{}reclassification-\discretionary{}{}{}operator} & Ch.~10, l.~7734 & \textcolor{gray}{\textbf{T3}} & \textit{(ordinal)} & Falsified if the operator fails to predict observed ideological sorting\ldots \\
\eqlab{eq:\discretionary{}{}{}10.\discretionary{}{}{}14-\discretionary{}{}{}kinetic-\discretionary{}{}{}power-\discretionary{}{}{}distribution} & Ch.~10, l.~7768 & \textcolor{blue!70!black}{\textbf{T2}} & Pew Research Center --- Gun Ownership Survey 2017 & Falsified if kinetic power is not concentrated in F\_enforce as a propor\ldots \\
\eqlab{eq:\discretionary{}{}{}10.\discretionary{}{}{}15-\discretionary{}{}{}temporal-\discretionary{}{}{}enclosure-\discretionary{}{}{}definition} & Ch.~10, l.~7820 & \textcolor{gray}{\textbf{T3}} & Reinhart-Sbrancia 2015 (Economic Policy) & Falsified if sovereign debt accumulation is demonstrably unrelated to p\ldots \\
\textbf{\eqlab{eq:\discretionary{}{}{}10.\discretionary{}{}{}16-\discretionary{}{}{}financial-\discretionary{}{}{}repression-\discretionary{}{}{}condition}} & Ch.~10, l.~7825 & \textcolor{green!60!black}{\textbf{T1}} & FRED 10Y Treasury Constant Maturity Rate (GS10) & Falsified if: (1) real interest rates remain positive throughout the 19\ldots \\
\eqlab{eq:\discretionary{}{}{}10.\discretionary{}{}{}17-\discretionary{}{}{}temporal-\discretionary{}{}{}extraction-\discretionary{}{}{}rate} & Ch.~10, l.~7825 & \textcolor{green!60!black}{\textbf{T1}} & FRED: Federal Debt Total Public Debt as \% of GDP (G\ldots & Falsified if the indicator function 1[r < g] is inactive during documen\ldots \\
\textbf{\eqlab{eq:\discretionary{}{}{}10.\discretionary{}{}{}18-\discretionary{}{}{}psi-\discretionary{}{}{}degradation-\discretionary{}{}{}function}} & Ch.~10, l.~7905 & \textcolor{green!60!black}{\textbf{T1}} & FRED: Real Median Hourly Compensation (B4701C0A052N\ldots & Falsified if: (1) Chow test at 1971 shows no statistically significant \ldots \\
\eqlab{eq:\discretionary{}{}{}10.\discretionary{}{}{}19-\discretionary{}{}{}saeculum-\discretionary{}{}{}phase-\discretionary{}{}{}map} & Ch.~10, l.~7900 & \textcolor{blue!70!black}{\textbf{T2}} & Strauss \& Howe, The Fourth Turning (1997) & Falsified if: (1) the generational archetype sequence does not follow t\ldots \\
\textbf{\eqlab{eq:\discretionary{}{}{}10.\discretionary{}{}{}20-\discretionary{}{}{}dalio-\discretionary{}{}{}extraction-\discretionary{}{}{}bound}} & Ch.~10, l.~7950 & \textcolor{blue!70!black}{\textbf{T2}} & Maddison Project Database 2020 (GDP, trade output) & Falsified if: (1) the reconstructed US Empire Index does not show a pea\ldots \\
\multicolumn{5}{l}{\small\bfseries Chapter 11: The Kinetic Guarantee: Arms Asymmetry, the Second Amendment, and the Disarmament Timeline} \\
\textbf{\eqlab{eq:\discretionary{}{}{}11.\discretionary{}{}{}1-\discretionary{}{}{}extraction-\discretionary{}{}{}precondition}} & Ch.~11, l.~7992 & \textcolor{green!60!black}{\textbf{T1}} & Primary legislation texts (sullivanlaw, nfa1934, mu\ldots & Falsified if a documented extraction system maintained stability with L\ldots \\
\eqlab{eq:\discretionary{}{}{}11.\discretionary{}{}{}2-\discretionary{}{}{}disarmament-\discretionary{}{}{}agenda-\discretionary{}{}{}path} & Ch.~11, l.~7992 & \textcolor{green!60!black}{\textbf{T1}} & Primary legislation and ruling texts (11 events) & Falsified if documented gun-control legislation passed under high $\Phi$\_loa\ldots \\
\eqlab{eq:\discretionary{}{}{}11.\discretionary{}{}{}3-\discretionary{}{}{}restriction-\discretionary{}{}{}precedent-\discretionary{}{}{}accumulation} & Ch.~11, l.~7992 & \textcolor{green!60!black}{\textbf{T1}} & Primary legislation and ruling texts (11 events) & Falsified if any documented legislative or ruling event reduced cumulat\ldots \\
\eqlab{eq:\discretionary{}{}{}11.\discretionary{}{}{}4-\discretionary{}{}{}restriction-\discretionary{}{}{}scope-\discretionary{}{}{}expansion} & Ch.~11, l.~7992 & \textcolor{green!60!black}{\textbf{T1}} & Primary legislation texts and enforcement records & Falsified if the historical record showed restriction scope declining o\ldots \\
\multicolumn{5}{l}{\small\bfseries Chapter 12: The Contradiction: Why Reform Serves the Algorithm} \\
\eqlab{eq:\discretionary{}{}{}12.\discretionary{}{}{}1-\discretionary{}{}{}reform-\discretionary{}{}{}absorption-\discretionary{}{}{}mechanism} & Ch.~12, l.~8859 & \textcolor{gray}{\textbf{T3}} & Piketty-Saez top income share data & Falsified if a documented reform reduces Elite extraction share (max) r\ldots \\
\eqlab{eq:\discretionary{}{}{}12.\discretionary{}{}{}2-\discretionary{}{}{}concession-\discretionary{}{}{}theorem} & Ch.~12, l.~8920 & \textcolor{blue!70!black}{\textbf{T2}} & Piven \& Cloward (1977) --- Poor People's Movements & Falsified if Elite-initiated reforms consistently exceed the minimum re\ldots \\
\eqlab{eq:\discretionary{}{}{}12.\discretionary{}{}{}3-\discretionary{}{}{}judiciary-\discretionary{}{}{}discrimination-\discretionary{}{}{}detection} & Ch.~12, l.~9071 & \textcolor{gray}{\textbf{T3}} & Washington v. Davis & Falsified if judicial detection function produces false positives for p\ldots \\
\eqlab{eq:\discretionary{}{}{}12.\discretionary{}{}{}4-\discretionary{}{}{}proxy-\discretionary{}{}{}discrimination-\discretionary{}{}{}equivalence} & Ch.~12, l.~9111 & \textcolor{gray}{\textbf{T3}} & Alexander (2010) --- The New Jim Crow & Falsified if P\_proxy achieves racial separation at rates significantly \ldots \\
\eqlab{eq:\discretionary{}{}{}12.\discretionary{}{}{}4a-\discretionary{}{}{}judicial-\discretionary{}{}{}double-\discretionary{}{}{}agent} & Ch.~12, l.~10000 & \textcolor{blue!70!black}{\textbf{T2}} & Local SCOTUS markdown corpus & Falsified if the SCOTUS corpus does not separate civil-rights/voting ca\ldots \\
\textbf{\eqlab{eq:\discretionary{}{}{}12.\discretionary{}{}{}5-\discretionary{}{}{}haitian-\discretionary{}{}{}theorem-\discretionary{}{}{}nonkinetic}} & Ch.~12, l.~8507 & \textcolor{green!60!black}{\textbf{T1}} & James (1938) The Black Jacobins & Falsified if any documented non-kinetic reform achieves $\Delta$max < 0 (susta\ldots \\
\textbf{\eqlab{eq:\discretionary{}{}{}12.\discretionary{}{}{}6-\discretionary{}{}{}haitian-\discretionary{}{}{}theorem-\discretionary{}{}{}kinetic}} & Ch.~12, l.~8507 & \textcolor{green!60!black}{\textbf{T1}} & James (1938) The Black Jacobins & Falsified if any kinetic event fails to produce at least temporary loca\ldots \\
\eqlab{eq:\discretionary{}{}{}12.\discretionary{}{}{}7-\discretionary{}{}{}gendered-\discretionary{}{}{}outgroup-\discretionary{}{}{}set} & Ch.~12, l.~9597 & \textcolor{gray}{\textbf{T3}} & \textit{(ordinal)} & Falsified if O\_gendered is shown to be structured differently from O\_ra\ldots \\
\eqlab{eq:\discretionary{}{}{}12.\discretionary{}{}{}8-\discretionary{}{}{}intersection-\discretionary{}{}{}double-\discretionary{}{}{}partition} & Ch.~12, l.~9602 & \textcolor{gray}{\textbf{T3}} & Crenshaw (1989) --- Demarginalizing the Intersection \ldots & Falsified if intersection O\_racialized $\cap$ O\_gendered does not show compo\ldots \\
\eqlab{eq:\discretionary{}{}{}12.\discretionary{}{}{}9-\discretionary{}{}{}multiplicative-\discretionary{}{}{}intersection-\discretionary{}{}{}compounding} & Ch.~12, l.~9607 & \textcolor{blue!70!black}{\textbf{T2}} & AAUW pay-gap data by race and gender & Falsified if Black women's wealth gap from white men is not multiplicat\ldots \\
\eqlab{eq:\discretionary{}{}{}12.\discretionary{}{}{}10-\discretionary{}{}{}nonviolence-\discretionary{}{}{}mandate-\discretionary{}{}{}function} & Ch.~12, l.~9709 & \textcolor{gray}{\textbf{T3}} & Carson (1981) --- In Struggle: SNCC and the Black Awa\ldots & Falsified if documented liberation movements adopting nonviolence achie\ldots \\
\multicolumn{5}{l}{\small\bfseries Chapter 13: The Global Containment Field: Scaling the Algorithm} \\
\eqlab{eq:\discretionary{}{}{}13.\discretionary{}{}{}1-\discretionary{}{}{}global-\discretionary{}{}{}compounding-\discretionary{}{}{}capacity} & Ch.~13, l.~9925 & \textcolor{blue!70!black}{\textbf{T2}} & World Bank development indicators --- GDP per capita \ldots & Falsified if post-colonial nation's capacity trajectory shows exponenti\ldots \\
\eqlab{eq:\discretionary{}{}{}13.\discretionary{}{}{}2-\discretionary{}{}{}global-\discretionary{}{}{}predatory-\discretionary{}{}{}minmax} & Ch.~13, l.~9973 & \textcolor{gray}{\textbf{T3}} & Klein (2007) --- The Shock Doctrine & Falsified if documented international economic relationships fail to ma\ldots \\
\eqlab{eq:\discretionary{}{}{}13.\discretionary{}{}{}3-\discretionary{}{}{}externalized-\discretionary{}{}{}enforcement-\discretionary{}{}{}envelope} & Ch.~13, l.~9993 & \textcolor{gray}{\textbf{T3}} & UN MINUSTAH mission records & Falsified if Haiti's domestic coercive capacity is shown to be structur\ldots \\
\eqlab{eq:\discretionary{}{}{}13.\discretionary{}{}{}4-\discretionary{}{}{}killick-\discretionary{}{}{}extraction-\discretionary{}{}{}equation} & Ch.~13, l.~10021 & \textcolor{gray}{\textbf{T3}} & NYT 'The Ransom' (2022) & Falsified if the extraction equation fails to predict subsequent debt i\ldots \\
\eqlab{eq:\discretionary{}{}{}13.\discretionary{}{}{}5-\discretionary{}{}{}legitimation-\discretionary{}{}{}constraint-\discretionary{}{}{}set} & Ch.~13, l.~10039 & \textcolor{gray}{\textbf{T3}} & \textit{(ordinal)} & Falsified if E\_global actors consistently operate outside their stated \ldots \\
\eqlab{eq:\discretionary{}{}{}13.\discretionary{}{}{}6-\discretionary{}{}{}authorization-\discretionary{}{}{}precondition} & Ch.~13, l.~10050 & \textcolor{gray}{\textbf{T3}} & \textit{(ordinal)} & Falsified if E\_global achieves extraction without satisfying authorizat\ldots \\
\eqlab{eq:\discretionary{}{}{}13.\discretionary{}{}{}7-\discretionary{}{}{}coercion-\discretionary{}{}{}legitimacy-\discretionary{}{}{}frontier} & Ch.~13, l.~10056 & \textcolor{gray}{\textbf{T3}} & \textit{(ordinal)} & Falsified if observed coercion levels in 'voluntary' agreements consist\ldots \\
\eqlab{eq:\discretionary{}{}{}13.\discretionary{}{}{}8-\discretionary{}{}{}firmin-\discretionary{}{}{}extraction-\discretionary{}{}{}threshold} & Ch.~13, l.~10062 & \textcolor{gray}{\textbf{T3}} & \textit{(ordinal)} & Falsified if E\_global achieves extraction with deployed force below F* \ldots \\
\eqlab{eq:\discretionary{}{}{}13.\discretionary{}{}{}9-\discretionary{}{}{}infeasible-\discretionary{}{}{}constraint-\discretionary{}{}{}set} & Ch.~13, l.~10075 & \textcolor{gray}{\textbf{T3}} & \textit{(ordinal)} & Falsified if E\_global achieves both extraction success and legitimation\ldots \\
\eqlab{eq:\discretionary{}{}{}13.\discretionary{}{}{}10-\discretionary{}{}{}asymmetric-\discretionary{}{}{}leverage-\discretionary{}{}{}ratio} & Ch.~13, l.~10100 & \textcolor{blue!70!black}{\textbf{T2}} & Firmin (1885) --- De l'égalité des races humaines & Falsified if a weaker actor's invocation of E\_global's legitimation arc\ldots \\
\eqlab{eq:\discretionary{}{}{}13.\discretionary{}{}{}11-\discretionary{}{}{}legitimation-\discretionary{}{}{}abandonment-\discretionary{}{}{}condition} & Ch.~13, l.~10107 & \textcolor{gray}{\textbf{T3}} & \textit{(ordinal)} & Falsified if E\_global maintains legitimation constraints even when extr\ldots \\
\eqlab{eq:\discretionary{}{}{}13.\discretionary{}{}{}12-\discretionary{}{}{}global-\discretionary{}{}{}containment-\discretionary{}{}{}strategy} & Ch.~13, l.~10123 & \textcolor{gray}{\textbf{T3}} & \textit{(ordinal)} & Falsified if global containment strategies differ structurally from the\ldots \\
\eqlab{eq:\discretionary{}{}{}13.\discretionary{}{}{}13-\discretionary{}{}{}un-\discretionary{}{}{}vote-\discretionary{}{}{}distribution} & Ch.~13, l.~10162 & \textcolor{gray}{\textbf{T3}} & UN Digital Library --- UNSC vote records & Falsified if UN Security Council voting distribution fails to map onto \ldots \\
\textbf{\eqlab{eq:\discretionary{}{}{}13.\discretionary{}{}{}14-\discretionary{}{}{}imperial-\discretionary{}{}{}core-\discretionary{}{}{}collapse}} & Ch.~13, l.~9208 & \textcolor{green!60!black}{\textbf{T1}} & World Bank Development Indicators --- GDP trajectory \ldots & Falsified if documented rising powers (China, OPEC, Asian Tigers) achie\ldots \\
\textbf{\eqlab{eq:\discretionary{}{}{}13.\discretionary{}{}{}15-\discretionary{}{}{}interface-\discretionary{}{}{}swap-\discretionary{}{}{}trigger}} & Ch.~13, l.~10280 & \textcolor{green!60!black}{\textbf{T1}} & Federal Reserve H.4.1: Historical Gold Stock Data (\ldots & Falsified if the gold-cover ratio was still above 80\% at the time of th\ldots \\
\eqlab{eq:\discretionary{}{}{}13.\discretionary{}{}{}16-\discretionary{}{}{}polymorphic-\discretionary{}{}{}reboot-\discretionary{}{}{}operator} & Ch.~13, l.~10310 & \textcolor{blue!70!black}{\textbf{T2}} & Gowa 1983 (Closing the Gold Window) & Falsified if Assets(E, A\_new) < Assets(E, A\_old) across any of the thre\ldots \\
\textbf{\eqlab{eq:\discretionary{}{}{}13.\discretionary{}{}{}17-\discretionary{}{}{}elite-\discretionary{}{}{}asset-\discretionary{}{}{}continuity-\discretionary{}{}{}invariant}} & Ch.~13, l.~10320 & \textcolor{green!60!black}{\textbf{T1}} & WID.world: Top-0.1\% US Wealth Share 1913-present (P\ldots & Falsified if the top-0.1\% US wealth share declines STRUCTURALLY (> 3 pe\ldots \\
\multicolumn{5}{l}{\small\bfseries Chapter 14: The Algorithmic Epoch: Real-Time Subjugation and the Necessity of the Counter-Virus} \\
\eqlab{eq:\discretionary{}{}{}14.\discretionary{}{}{}1-\discretionary{}{}{}algo-\discretionary{}{}{}prior-\discretionary{}{}{}inheritance} & Ch.~14, l.~10757 & \textcolor{blue!70!black}{\textbf{T2}} & Angwin et al. (2016) --- Machine Bias (ProPublica) & Falsified if ML systems trained on pre-2024 US data show no racial disp\ldots \\
\eqlab{eq:\discretionary{}{}{}14.\discretionary{}{}{}2-\discretionary{}{}{}realtime-\discretionary{}{}{}proximity-\discretionary{}{}{}estimate} & Ch.~14, l.~10777 & \textcolor{blue!70!black}{\textbf{T2}} & Twitter/X political affinity graph data (academic d\ldots & Falsified if real-time social-affinity graph G(t) shows no correlation \ldots \\
\eqlab{eq:\discretionary{}{}{}14.\discretionary{}{}{}3-\discretionary{}{}{}orthogonal-\discretionary{}{}{}vector-\discretionary{}{}{}injection} & Ch.~14, l.~10793 & \textcolor{gray}{\textbf{T3}} & \textit{(ordinal)} & Falsified if documented interference engine operations increase class-a\ldots \\
\eqlab{eq:\discretionary{}{}{}14.\discretionary{}{}{}4-\discretionary{}{}{}correction-\discretionary{}{}{}signal-\discretionary{}{}{}damping} & Ch.~14, l.~10807 & \textcolor{gray}{\textbf{T3}} & \textit{(ordinal)} & Falsified if scale-accurate counter-signals are shown to reduce identit\ldots \\
\eqlab{eq:\discretionary{}{}{}14.\discretionary{}{}{}5-\discretionary{}{}{}optical-\discretionary{}{}{}enclosure} & Ch.~14, l.~10820 & \textcolor{gray}{\textbf{T3}} & Pew Research --- Police-community relations by race & Falsified if documented perception studies show O's population maintain\ldots \\
\eqlab{eq:\discretionary{}{}{}14.\discretionary{}{}{}6-\discretionary{}{}{}decoy-\discretionary{}{}{}transfer-\discretionary{}{}{}coefficient} & Ch.~14, l.~10834 & \textcolor{blue!70!black}{\textbf{T2}} & \textit{(pending)} & Falsified if documented cases show kinetic class energy passing to E at\ldots \\
\eqlab{eq:\discretionary{}{}{}14.\discretionary{}{}{}7-\discretionary{}{}{}ddos-\discretionary{}{}{}bandwidth-\discretionary{}{}{}ceiling} & Ch.~14, l.~10858 & \textcolor{blue!70!black}{\textbf{T2}} & Armed Conflict Location \& Event Data (ACLED) --- US p\ldots & Falsified if F\_enforce simultaneously suppresses n > $\Gamma$ nodes in any doc\ldots \\
\eqlab{eq:\discretionary{}{}{}14.\discretionary{}{}{}8-\discretionary{}{}{}armed-\discretionary{}{}{}swarm-\discretionary{}{}{}bandwidth} & Ch.~14, l.~10869 & \textcolor{blue!70!black}{\textbf{T2}} & \textit{(pending)} & Falsified if F\_enforce simultaneously neutralizes armed nodes with forc\ldots \\
\eqlab{eq:\discretionary{}{}{}14.\discretionary{}{}{}9-\discretionary{}{}{}kinetic-\discretionary{}{}{}capital-\discretionary{}{}{}asymmetry} & Ch.~14, l.~10887 & \textcolor{blue!70!black}{\textbf{T2}} & Pew Research Center --- Gun Ownership Survey 2017 & Falsified if gun-ownership data shows no racial asymmetry in kinetic ca\ldots \\
\eqlab{eq:\discretionary{}{}{}14.\discretionary{}{}{}10-\discretionary{}{}{}sustained-\discretionary{}{}{}siege-\discretionary{}{}{}cost} & Ch.~14, l.~10909 & \textcolor{blue!70!black}{\textbf{T2}} & \textit{(pending)} & Falsified if F\_enforce besieges multiple simultaneous armed nodes with \ldots \\
\eqlab{eq:\discretionary{}{}{}14.\discretionary{}{}{}11-\discretionary{}{}{}escalation-\discretionary{}{}{}ceiling-\discretionary{}{}{}tiers} & Ch.~14, l.~10919 & \textcolor{gray}{\textbf{T3}} & \textit{(ordinal)} & Falsified if F\_enforce escalation tiers do not follow the predicted ban\ldots \\
\eqlab{eq:\discretionary{}{}{}14.\discretionary{}{}{}12-\discretionary{}{}{}inclusion-\discretionary{}{}{}predicate} & Ch.~14, l.~10948 & \textcolor{blue!70!black}{\textbf{T2}} & Bureau of Labor Statistics contingent worker survey & Falsified if data shows V(x,t) $\leq$ R(x,t) for majority members of I\_buffe\ldots \\
\eqlab{eq:\discretionary{}{}{}14.\discretionary{}{}{}13-\discretionary{}{}{}terminal-\discretionary{}{}{}interface-\discretionary{}{}{}swap} & Ch.~14, l.~10959 & \textcolor{gray}{\textbf{T3}} & \textit{(ordinal)} & Falsified if documented late-stage extraction shows I\_buffer benefits m\ldots \\
\eqlab{eq:\discretionary{}{}{}14.\discretionary{}{}{}14-\discretionary{}{}{}liability-\discretionary{}{}{}phase-\discretionary{}{}{}optimization} & Ch.~14, l.~10969 & \textcolor{gray}{\textbf{T3}} & \textit{(ordinal)} & Falsified if documented terminal-phase extraction shows E maintaining I\ldots \\
\eqlab{eq:\discretionary{}{}{}14.\discretionary{}{}{}15-\discretionary{}{}{}noise-\discretionary{}{}{}spectrum-\discretionary{}{}{}index} & Ch.~14, l.~11017 & \textcolor{blue!70!black}{\textbf{T2}} & ACLED --- US social movements 2011--2020 & Falsified if noise-spectrum index of a documented swarm fails to predic\ldots \\
\eqlab{eq:\discretionary{}{}{}14.\discretionary{}{}{}16-\discretionary{}{}{}agnostic-\discretionary{}{}{}swarm-\discretionary{}{}{}payload} & Ch.~14, l.~11036 & \textcolor{gray}{\textbf{T3}} & \textit{(ordinal)} & Falsified if documented multi-node mobilizations with ideological diver\ldots \\
\eqlab{eq:\discretionary{}{}{}14.\discretionary{}{}{}17-\discretionary{}{}{}botnet-\discretionary{}{}{}load-\discretionary{}{}{}theorem} & Ch.~14, l.~11051 & \textcolor{blue!70!black}{\textbf{T2}} & ACLED --- 2020 US protest dataset & Falsified if enforcement grid contains a swarm with n > $\Gamma$\_armed nodes a\ldots \\
\eqlab{eq:\discretionary{}{}{}14.\discretionary{}{}{}18-\discretionary{}{}{}ideological-\discretionary{}{}{}routing-\discretionary{}{}{}operator} & Ch.~14, l.~11072 & \textcolor{gray}{\textbf{T3}} & \textit{(ordinal)} & Falsified if nodes in documented distributed mobilizations fail to rout\ldots \\
\eqlab{eq:\discretionary{}{}{}14.\discretionary{}{}{}19-\discretionary{}{}{}zugzwang-\discretionary{}{}{}payoff-\discretionary{}{}{}function} & Ch.~14, l.~11122 & \textcolor{blue!70!black}{\textbf{T2}} & \textit{(pending)} & Falsified if Elite actors systematically choose attack over stand-down \ldots \\
\eqlab{eq:\discretionary{}{}{}14.\discretionary{}{}{}20-\discretionary{}{}{}defection-\discretionary{}{}{}cascade} & Ch.~14, l.~11164 & \textcolor{blue!70!black}{\textbf{T2}} & James (1938) --- The Black Jacobins (Polish Legion de\ldots & Falsified if documented enforcement nodes fail to defect when community\ldots \\
\eqlab{eq:\discretionary{}{}{}14.\discretionary{}{}{}21-\discretionary{}{}{}empathy-\discretionary{}{}{}permeability} & Ch.~14, l.~11196 & \textcolor{blue!70!black}{\textbf{T2}} & James (1938) --- The Black Jacobins & Falsified if trauma overlap between enforcement and out-group nodes fai\ldots \\
\eqlab{eq:\discretionary{}{}{}14.\discretionary{}{}{}22-\discretionary{}{}{}semantic-\discretionary{}{}{}overwrite} & Ch.~14, l.~11222 & \textcolor{gray}{\textbf{T3}} & Haitian Constitution of 1805 (primary source) & Falsified if 1805 Haitian Constitution's ideological redefinition faile\ldots \\
\multicolumn{5}{l}{\small\bfseries Chapter 17: Conclusion} \\
\eqlab{eq:\discretionary{}{}{}17.\discretionary{}{}{}1-\discretionary{}{}{}racism-\discretionary{}{}{}as-\discretionary{}{}{}institutional-\discretionary{}{}{}vector} & Ch.~17, l.~11531 & \textcolor{gray}{\textbf{T3}} & \textit{(ordinal)} & Falsified if documented cases of marginalized-group prejudice achieve s\ldots \\
\eqlab{eq:\discretionary{}{}{}17.\discretionary{}{}{}2-\discretionary{}{}{}o-\discretionary{}{}{}terminal-\discretionary{}{}{}expansion} & Ch.~17, l.~11586 & \textcolor{blue!70!black}{\textbf{T2}} & US Census Bureau 2020 demographic projections & Falsified if documented demographic trends show O contracting rather th\ldots \\
\multicolumn{5}{l}{\small\bfseries Chapter 19: Equation Registry and Era-Level Calibration} \\
\eqlab{eq:\discretionary{}{}{}19.\discretionary{}{}{}1-\discretionary{}{}{}kernel-\discretionary{}{}{}objective-\discretionary{}{}{}registry-\discretionary{}{}{}entry} & Ch.~19, l.~11717 & \textcolor{gray}{\textbf{T3}} & \textit{(ordinal)} & Redundant with eq:5 and eq:56; falsified by the same conditions as thos\ldots \\
\eqlab{eq:\discretionary{}{}{}19.\discretionary{}{}{}2-\discretionary{}{}{}suppression-\discretionary{}{}{}envelope-\discretionary{}{}{}registry-\discretionary{}{}{}entry} & Ch.~19, l.~11721 & \textcolor{gray}{\textbf{T3}} & \textit{(ordinal)} & Redundant with eq:8 and eq:58; falsified by the same conditions as thos\ldots \\
\eqlab{eq:\discretionary{}{}{}19.\discretionary{}{}{}3-\discretionary{}{}{}crash-\discretionary{}{}{}condition-\discretionary{}{}{}registry-\discretionary{}{}{}entry} & Ch.~19, l.~11725 & \textcolor{gray}{\textbf{T3}} & \textit{(ordinal)} & Redundant with eq:9 and eq:59; falsified by the same conditions as thos\ldots \\
\eqlab{eq:\discretionary{}{}{}19.\discretionary{}{}{}4-\discretionary{}{}{}tier-\discretionary{}{}{}ordering-\discretionary{}{}{}registry-\discretionary{}{}{}entry} & Ch.~19, l.~11729 & \textcolor{gray}{\textbf{T3}} & \textit{(ordinal)} & Redundant with eq:60; falsified if any tier ordering is reversed in com\ldots \\
\end{longtable}
\endgroup

%
% Requires: longtable, booktabs, xcolor, array (loaded in main preamble)

\begingroup
\small
\setlength{\LTpre}{6pt}
\setlength{\LTpost}{6pt}
% Equation labels are long unbreakable \texttt strings (up to 47 chars). Without
% explicit break opportunities they overrun the label column and print on top of
% the case-study column. \eqlab emits the label with zero-width discretionary
% breaks after ':', '-' and '.', inserted by the generator.
\providecommand{\eqlab}[1]{\texttt{#1}}
\begin{longtable}{@{}%
  >{\raggedright\arraybackslash}p{3.2cm}
  >{\raggedright\arraybackslash}p{2.5cm}
  >{\centering\arraybackslash}p{1.0cm}
  >{\raggedright\arraybackslash}p{3.3cm}
  >{\raggedright\arraybackslash}p{5.2cm}@{}}
\textbf{Equation label} &
\textbf{Case study location} &
\textbf{Tier} &
\textbf{Primary data source} &
\textbf{Falsification (truncated)} \\
\endfirsthead
\textbf{Equation label} &
\textbf{Case study location} &
\textbf{Tier} &
\textbf{Primary data source} &
\textbf{Falsification (truncated)} \\
\endhead
\multicolumn{5}{r}{\small\itshape continued on next page} \\
\endfoot
\endlastfoot
\multicolumn{5}{l}{\small\bfseries Chapter 1: Redefining Racism} \\
\eqlab{eq:\discretionary{}{}{}1.\discretionary{}{}{}1-\discretionary{}{}{}enclosure-\discretionary{}{}{}score} & Ch.~1, l.~241 & \textcolor{gray}{\textbf{T3}} & \textit{(ordinal)} & Falsified if a population with all three outlets blocked (e\_i=1) shows \ldots \\
\eqlab{eq:\discretionary{}{}{}1.\discretionary{}{}{}2-\discretionary{}{}{}total-\discretionary{}{}{}enclosure} & Ch.~1, l.~257 & \textcolor{gray}{\textbf{T3}} & \textit{(ordinal)} & Falsified if a fully blocked population (e\_1=e\_2=e\_3=1) achieves coordi\ldots \\
\eqlab{eq:\discretionary{}{}{}1.\discretionary{}{}{}3-\discretionary{}{}{}conventional-\discretionary{}{}{}causal-\discretionary{}{}{}arrow} & Ch.~1, l.~291 & \textcolor{gray}{\textbf{T3}} & \textit{(ordinal)} & Falsified if documented evidence shows individual prejudice predating E\ldots \\
\eqlab{eq:\discretionary{}{}{}1.\discretionary{}{}{}4-\discretionary{}{}{}elite-\discretionary{}{}{}causal-\discretionary{}{}{}arrow} & Ch.~1, l.~297 & \textcolor{gray}{\textbf{T3}} & Zurara & Falsified if racialization systems are shown to emerge without Elite ec\ldots \\
\textbf{\eqlab{eq:\discretionary{}{}{}1.\discretionary{}{}{}5-\discretionary{}{}{}kernel-\discretionary{}{}{}optimization}} & Ch.~1, l.~495 & \textcolor{green!60!black}{\textbf{T1}} & Piketty & Falsified if a documented period shows sustained decline in Elite extra\ldots \\
\eqlab{eq:\discretionary{}{}{}1.\discretionary{}{}{}6-\discretionary{}{}{}interface-\discretionary{}{}{}strategy-\discretionary{}{}{}set} & Ch.~1, l.~576 & \textcolor{gray}{\textbf{T3}} & \textit{(ordinal)} & Falsified if a documented Elite extraction system used a strategy outsi\ldots \\
\eqlab{eq:\discretionary{}{}{}1.\discretionary{}{}{}7-\discretionary{}{}{}interface-\discretionary{}{}{}optimizer} & Ch.~1, l.~608 & \textcolor{blue!70!black}{\textbf{T2}} & Dudziak (2000) --- Cold War Civil Rights & Falsified if a historical interface transition increased net cost C\_tot\ldots \\
\textbf{\eqlab{eq:\discretionary{}{}{}1.\discretionary{}{}{}8-\discretionary{}{}{}suppression-\discretionary{}{}{}envelope}} & Ch.~1, l.~714 & \textcolor{blue!70!black}{\textbf{T2}} & BLS Union Membership Historical Series & Falsified if M(t) consistently exceeds $\Sigma$\_sup(t) for extended periods wi\ldots \\
\eqlab{eq:\discretionary{}{}{}1.\discretionary{}{}{}9-\discretionary{}{}{}crash-\discretionary{}{}{}condition-\discretionary{}{}{}derivative} & Ch.~1, l.~710 & \textcolor{gray}{\textbf{T3}} & \textit{(ordinal)} & Falsified if a crash event (Bacon's Rebellion, Haitian Revolution) occu\ldots \\
\eqlab{eq:\discretionary{}{}{}1.\discretionary{}{}{}10-\discretionary{}{}{}effective-\discretionary{}{}{}coherence} & Ch.~1, l.~716 & \textcolor{blue!70!black}{\textbf{T2}} & ANES cross-racial coalition data 1948--2020 & Falsified if a crash event occurs when M\_eff(t) < $\tau$, or if no crash occ\ldots \\
\eqlab{eq:\discretionary{}{}{}1.\discretionary{}{}{}11-\discretionary{}{}{}phase-\discretionary{}{}{}loading} & Ch.~1, l.~838 & \textcolor{gray}{\textbf{T3}} & ANES Time Series Study 1948--2020 & Falsified if survey-measured cross-group solidarity remains high when $\Phi$\ldots \\
\multicolumn{5}{l}{\small\bfseries Chapter 2: Version 1.0: Initializing the Vector (15th-Century Portugal)} \\
\eqlab{eq:\discretionary{}{}{}2.\discretionary{}{}{}1-\discretionary{}{}{}racism-\discretionary{}{}{}vector} & Ch.~2, l.~1091 & \textcolor{gray}{\textbf{T3}} & \textit{(ordinal)} & Falsified if documented racism operates without a state-power direction\ldots \\
\eqlab{eq:\discretionary{}{}{}2.\discretionary{}{}{}2-\discretionary{}{}{}status-\discretionary{}{}{}suppression-\discretionary{}{}{}allocation} & Ch.~2, l.~1163 & \textcolor{gray}{\textbf{T3}} & Du Bois (1935) --- Black Reconstruction in America & Falsified if Buffer Class loyalty is maintained without a non-material \ldots \\
\eqlab{eq:\discretionary{}{}{}2.\discretionary{}{}{}3-\discretionary{}{}{}psi-\discretionary{}{}{}null-\discretionary{}{}{}in-\discretionary{}{}{}survival} & Ch.~2, l.~1167 & \textcolor{gray}{\textbf{T3}} & \textit{(ordinal)} & Falsified if material wages to Buffer Class are shown to be independent\ldots \\
\eqlab{eq:\discretionary{}{}{}2.\discretionary{}{}{}4-\discretionary{}{}{}roman-\discretionary{}{}{}membership-\discretionary{}{}{}function} & Ch.~2, l.~1208 & \textcolor{gray}{\textbf{T3}} & Bradley (1994) --- Slavery and Society at Rome & Falsified if Roman enslaved population is shown to be phenotypically ho\ldots \\
\eqlab{eq:\discretionary{}{}{}2.\discretionary{}{}{}5-\discretionary{}{}{}american-\discretionary{}{}{}racial-\discretionary{}{}{}function} & Ch.~2, l.~1217 & \textcolor{gray}{\textbf{T3}} & Hening's Statutes at Large (Virginia) Vol. 2 & Falsified if documented American slavery assigned status based on non-p\ldots \\
\eqlab{eq:\discretionary{}{}{}2.\discretionary{}{}{}6-\discretionary{}{}{}roman-\discretionary{}{}{}american-\discretionary{}{}{}structural-\discretionary{}{}{}parallel} & Ch.~2, l.~1235 & \textcolor{gray}{\textbf{T3}} & \textit{(ordinal)} & Falsified if the structural mapping between Roman and American slavery \ldots \\
\eqlab{eq:\discretionary{}{}{}2.\discretionary{}{}{}7-\discretionary{}{}{}three-\discretionary{}{}{}tier-\discretionary{}{}{}portugal-\discretionary{}{}{}inequality} & Ch.~2, l.~1280 & \textcolor{gray}{\textbf{T3}} & \textit{(ordinal)} & Falsified if measured material outcomes show Benefit(I\_buffer) $\geq$ Benefi\ldots \\
\eqlab{eq:\discretionary{}{}{}2.\discretionary{}{}{}8-\discretionary{}{}{}church-\discretionary{}{}{}science-\discretionary{}{}{}feedback-\discretionary{}{}{}loop} & Ch.~2, l.~1335 & \textcolor{gray}{\textbf{T3}} & Cartwright (1851) --- Diseases and Peculiarities of t\ldots & Falsified if the religious-scientific legitimation loop is disrupted fr\ldots \\
\multicolumn{5}{l}{\small\bfseries Chapter 3: The Application: Bacon's Rebellion, the Buffer Class, and the Constitutional Patch (1676--1787)} \\
\textbf{\eqlab{eq:\discretionary{}{}{}3.\discretionary{}{}{}1-\discretionary{}{}{}bacon-\discretionary{}{}{}solidarity-\discretionary{}{}{}condition}} & Ch.~3, l.~1580 & \textcolor{green!60!black}{\textbf{T1}} & Morgan (1975) --- American Slavery & Falsified if an antebellum cross-racial coalition reached M(t) > $\tau$ with\ldots \\
\eqlab{eq:\discretionary{}{}{}3.\discretionary{}{}{}2-\discretionary{}{}{}boundary-\discretionary{}{}{}enforcement-\discretionary{}{}{}1691} & Ch.~3, l.~1507 & \textcolor{gray}{\textbf{T3}} & Hening's Statutes at Large (Virginia) Vol. 3 & Falsified if racial boundary enforcement persisted without legal penalt\ldots \\
\eqlab{eq:\discretionary{}{}{}3.\discretionary{}{}{}3-\discretionary{}{}{}race-\discretionary{}{}{}destroys-\discretionary{}{}{}class} & Ch.~3, l.~1539 & \textcolor{gray}{\textbf{T3}} & Morgan (1975) --- American Slavery & Falsified if colonial-era racial coding produced class solidarity rathe\ldots \\
\eqlab{eq:\discretionary{}{}{}3.\discretionary{}{}{}4-\discretionary{}{}{}three-\discretionary{}{}{}tier-\discretionary{}{}{}post-\discretionary{}{}{}bacon} & Ch.~3, l.~1560 & \textcolor{gray}{\textbf{T3}} & \textit{(ordinal)} & Falsified if the post-Bacon ordering shows Buffer Class benefits exceed\ldots \\
\eqlab{eq:\discretionary{}{}{}3.\discretionary{}{}{}5-\discretionary{}{}{}kinetic-\discretionary{}{}{}necessary-\discretionary{}{}{}condition} & Ch.~3, l.~1595 & \textcolor{green!60!black}{\textbf{T1}} & Morgan (1975) --- American Slavery & Falsified if documented cases show Elite extraction maintained when K\_E\ldots \\
\multicolumn{5}{l}{\small\bfseries Chapter 5: The Gendered Axis: Coverture, Eugenics, and the Reproductive Extraction Kernel} \\
\eqlab{eq:\discretionary{}{}{}5.\discretionary{}{}{}1-\discretionary{}{}{}compliance-\discretionary{}{}{}differential} & Ch.~5, l.~2272 & \textcolor{blue!70!black}{\textbf{T2}} & Baptist (2014) --- The Half Has Never Been Told & Falsified if V\_c(x) $\leq$ V\_r(x) for subjects in documented extraction regi\ldots \\
\multicolumn{5}{l}{\small\bfseries Chapter 6: The Enforcement Engine: Slave Patrols, the 13th Amendment, and the Compounding Model (1704--1865)} \\
\eqlab{eq:\discretionary{}{}{}6.\discretionary{}{}{}1-\discretionary{}{}{}lethal-\discretionary{}{}{}autonomy-\discretionary{}{}{}gradient} & Ch.~6, l.~2889 & \textcolor{blue!70!black}{\textbf{T2}} & LEOSA --- 18 U.S.C. §926B statutory text & Falsified if Mapping Police Violence data shows parity in police killin\ldots \\
\textbf{\eqlab{eq:\discretionary{}{}{}6.\discretionary{}{}{}2-\discretionary{}{}{}four-\discretionary{}{}{}tier-\discretionary{}{}{}benefit-\discretionary{}{}{}ordering}} & Ch.~6, l.~2994 & \textcolor{gray}{\textbf{T3}} & Piketty-Saez-Zucman Distributional National Accounts & Falsified if comprehensive wealth, health, or carceral data shows Benef\ldots \\
\eqlab{eq:\discretionary{}{}{}6.\discretionary{}{}{}3-\discretionary{}{}{}thirteenth-\discretionary{}{}{}amendment-\discretionary{}{}{}exception} & Ch.~6, l.~3080 & \textcolor{gray}{\textbf{T3}} & Blackmon (2008) --- Slavery by Another Name & Falsified if 13th Amendment exception clause is shown to have been appl\ldots \\
\eqlab{eq:\discretionary{}{}{}6.\discretionary{}{}{}4-\discretionary{}{}{}additive-\discretionary{}{}{}policy-\discretionary{}{}{}model} & Ch.~6, l.~3144 & \textcolor{gray}{\textbf{T3}} & \textit{(ordinal)} & Falsified if additive accumulation predicts compound policy outcomes wi\ldots \\
\eqlab{eq:\discretionary{}{}{}6.\discretionary{}{}{}5-\discretionary{}{}{}compounding-\discretionary{}{}{}temporal-\discretionary{}{}{}model} & Ch.~6, l.~3154 & \textcolor{blue!70!black}{\textbf{T2}} & Hamilton \& Darity (2017) --- The Political Economy of\ldots & Falsified if policy shock effects on O\textasciicircum{}capacity are shown to be indepen\ldots \\
\eqlab{eq:\discretionary{}{}{}6.\discretionary{}{}{}6-\discretionary{}{}{}asymmetric-\discretionary{}{}{}enforcement-\discretionary{}{}{}multiplier} & Ch.~6, l.~3172 & \textcolor{green!60!black}{\textbf{T1}} & ACLU (2020) --- A Tale of Two Countries: Racially Tar\ldots & Falsified if ACLU cannabis arrest data shows no racial disparity in pol\ldots \\
\eqlab{eq:\discretionary{}{}{}6.\discretionary{}{}{}7-\discretionary{}{}{}policy-\discretionary{}{}{}sequence-\discretionary{}{}{}noncommutative} & Ch.~6, l.~3177 & \textcolor{gray}{\textbf{T3}} & \textit{(ordinal)} & Falsified if order of policy application (P\_1 then P\_2 vs. P\_2 then P\_1\ldots \\
\eqlab{eq:\discretionary{}{}{}6.\discretionary{}{}{}8-\discretionary{}{}{}capacity-\discretionary{}{}{}chain-\discretionary{}{}{}1619} & Ch.~6, l.~3193 & \textcolor{blue!70!black}{\textbf{T2}} & Darity \& Mullen (2020) --- From Here to Equality (rep\ldots & Falsified if O\_1619 capacity is non-significantly different from O\_1450\ldots \\
\eqlab{eq:\discretionary{}{}{}6.\discretionary{}{}{}9-\discretionary{}{}{}capacity-\discretionary{}{}{}chain-\discretionary{}{}{}1865} & Ch.~6, l.~3196 & \textcolor{blue!70!black}{\textbf{T2}} & Blackmon (2008) --- Slavery by Another Name & Falsified if O\_1865 capacity is not reduced from O\_1619 as measured by \ldots \\
\eqlab{eq:\discretionary{}{}{}6.\discretionary{}{}{}10-\discretionary{}{}{}capacity-\discretionary{}{}{}chain-\discretionary{}{}{}1934} & Ch.~6, l.~3199 & \textcolor{green!60!black}{\textbf{T1}} & Mapping Inequality --- HOLC Redlining Maps & Falsified if redlining-exposed tracts show no differential wealth accum\ldots \\
\eqlab{eq:\discretionary{}{}{}6.\discretionary{}{}{}11-\discretionary{}{}{}capacity-\discretionary{}{}{}chain-\discretionary{}{}{}1971} & Ch.~6, l.~3230 & \textcolor{green!60!black}{\textbf{T1}} & ACLU (2020) --- Cannabis Arrests Report & Falsified if War on Drugs enforcement shows no differential impact on O\ldots \\
\textbf{\eqlab{eq:\discretionary{}{}{}6.\discretionary{}{}{}12-\discretionary{}{}{}capacity-\discretionary{}{}{}compounding-\discretionary{}{}{}full}} & Ch.~6, l.~3169 & \textcolor{green!60!black}{\textbf{T1}} & ACLU (2020) --- A Tale of Two Countries & Falsified if any one policy shock ($\alpha$, $\beta$, $\gamma$, $\delta$) is shown to have zero ma\ldots \\
\multicolumn{5}{l}{\small\bfseries Chapter 7: The Containment: Pullman, Redlining, and the Wages of Whiteness (1894--1965)} \\
\eqlab{eq:\discretionary{}{}{}7.\discretionary{}{}{}1-\discretionary{}{}{}pullman-\discretionary{}{}{}corollary} & Ch.~7, l.~3624 & \textcolor{gray}{\textbf{T3}} & Roediger (1991) --- The Wages of Whiteness & Falsified if a documented case shows I\_buffer achieving collective gain\ldots \\
\multicolumn{5}{l}{\small\bfseries Chapter 8: Tweedism and the Puppet Class: The Algorithmic Filter on Democracy} \\
\eqlab{eq:\discretionary{}{}{}8.\discretionary{}{}{}1-\discretionary{}{}{}vertex-\discretionary{}{}{}partition-\discretionary{}{}{}graph} & Ch.~8, l.~4020 & \textcolor{gray}{\textbf{T3}} & \textit{(ordinal)} & Falsified if leader/follower partition fails to predict convergence in \ldots \\
\eqlab{eq:\discretionary{}{}{}8.\discretionary{}{}{}2-\discretionary{}{}{}tier-\discretionary{}{}{}to-\discretionary{}{}{}graph-\discretionary{}{}{}mapping} & Ch.~8, l.~4024 & \textcolor{gray}{\textbf{T3}} & \textit{(ordinal)} & Falsified if the tier-to-graph mapping produces inconsistent prediction\ldots \\
\eqlab{eq:\discretionary{}{}{}8.\discretionary{}{}{}3-\discretionary{}{}{}fractional-\discretionary{}{}{}agent-\discretionary{}{}{}dynamics} & Ch.~8, l.~4032 & \textcolor{gray}{\textbf{T3}} & Li et al. --- Mittag-Leffler stability results (mathe\ldots & Falsified if integer-order dynamics ($\alpha$=1) predict social-mobility conve\ldots \\
\eqlab{eq:\discretionary{}{}{}8.\discretionary{}{}{}4-\discretionary{}{}{}compounding-\discretionary{}{}{}chain-\discretionary{}{}{}formal} & Ch.~8, l.~3187 & \textcolor{blue!70!black}{\textbf{T2}} & Federal Reserve SCF racial wealth data & Falsified if multiplicative chain prediction diverges significantly fro\ldots \\
\eqlab{eq:\discretionary{}{}{}8.\discretionary{}{}{}5-\discretionary{}{}{}stationary-\discretionary{}{}{}leader-\discretionary{}{}{}condition} & Ch.~8, l.~4046 & \textcolor{gray}{\textbf{T3}} & \textit{(ordinal)} & Falsified if documented Elite actors show systematic response to u\_i (u\ldots \\
\eqlab{eq:\discretionary{}{}{}8.\discretionary{}{}{}6-\discretionary{}{}{}proximity-\discretionary{}{}{}connectivity-\discretionary{}{}{}condition} & Ch.~8, l.~4128 & \textcolor{gray}{\textbf{T3}} & \textit{(ordinal)} & Falsified if social-affinity graph connectivity is shown to be independ\ldots \\
\eqlab{eq:\discretionary{}{}{}8.\discretionary{}{}{}7-\discretionary{}{}{}navigation-\discretionary{}{}{}components} & Ch.~8, l.~4220 & \textcolor{gray}{\textbf{T3}} & \textit{(ordinal)} & Falsified if goal and constraint functions do not bound follower trajec\ldots \\
\eqlab{eq:\discretionary{}{}{}8.\discretionary{}{}{}8-\discretionary{}{}{}navigation-\discretionary{}{}{}function} & Ch.~8, l.~4223 & \textcolor{gray}{\textbf{T3}} & \textit{(ordinal)} & Falsified if $\phi$\_i fails to produce convergence in any graph satisfying t\ldots \\
\eqlab{eq:\discretionary{}{}{}8.\discretionary{}{}{}9-\discretionary{}{}{}gradient-\discretionary{}{}{}control-\discretionary{}{}{}law} & Ch.~8, l.~4228 & \textcolor{gray}{\textbf{T3}} & \textit{(ordinal)} & Falsified if negative gradient of $\phi$\_i fails to drive convergence in any\ldots \\
\eqlab{eq:\discretionary{}{}{}8.\discretionary{}{}{}10-\discretionary{}{}{}convex-\discretionary{}{}{}hull-\discretionary{}{}{}convergence} & Ch.~8, l.~4316 & \textcolor{gray}{\textbf{T3}} & \textit{(ordinal)} & Falsified if followers fail to converge to Elite convex hull in any gra\ldots \\
\eqlab{eq:\discretionary{}{}{}8.\discretionary{}{}{}11-\discretionary{}{}{}mittag-\discretionary{}{}{}leffler-\discretionary{}{}{}stability} & Ch.~8, l.~4402 & \textcolor{gray}{\textbf{T3}} & Chetty et al. (2018) --- Race and Economic Opportunit\ldots & Falsified if social-mobility data for $\alpha$ $\in$ (0,1) shows exponential rathe\ldots \\
\textbf{\eqlab{eq:\discretionary{}{}{}8.\discretionary{}{}{}12-\discretionary{}{}{}class-\discretionary{}{}{}alignment-\discretionary{}{}{}base-\discretionary{}{}{}waves}} & Ch.~8, l.~4968 & \textcolor{green!60!black}{\textbf{T1}} & Google Trends & Falsified if class alignment signal is not measurable as a distinct fre\ldots \\
\eqlab{eq:\discretionary{}{}{}8.\discretionary{}{}{}13-\discretionary{}{}{}subgroup-\discretionary{}{}{}compound-\discretionary{}{}{}phase} & Ch.~8, l.~4968 & \textcolor{blue!70!black}{\textbf{T2}} & Google Trends & Falsified if compound phase $\Phi$\_j fails to predict subgroup political ali\ldots \\
\eqlab{eq:\discretionary{}{}{}8.\discretionary{}{}{}14-\discretionary{}{}{}net-\discretionary{}{}{}solidarity-\discretionary{}{}{}signal} & Ch.~8, l.~4968 & \textcolor{green!60!black}{\textbf{T1}} & Google Trends & Falsified if net solidarity signal cannot be measured as distinct from \ldots \\
\eqlab{eq:\discretionary{}{}{}8.\discretionary{}{}{}15-\discretionary{}{}{}solidarity-\discretionary{}{}{}collapse-\discretionary{}{}{}condition} & Ch.~8, l.~4968 & \textcolor{green!60!black}{\textbf{T1}} & Google Trends & Falsified if M(t) remains above $\tau$ even when S\_class approaches zero acr\ldots \\
\eqlab{eq:\discretionary{}{}{}8.\discretionary{}{}{}16-\discretionary{}{}{}interference-\discretionary{}{}{}control-\discretionary{}{}{}objective} & Ch.~8, l.~4968 & \textcolor{green!60!black}{\textbf{T1}} & Google Trends & Falsified if observed interference engine output increases S\_class rath\ldots \\
\eqlab{eq:\discretionary{}{}{}8.\discretionary{}{}{}17-\discretionary{}{}{}circular-\discretionary{}{}{}dispersion-\discretionary{}{}{}operator} & Ch.~8, l.~4968 & \textcolor{blue!70!black}{\textbf{T2}} & ANES Time Series --- cross-group solidarity items & Falsified if circular dispersion of phase values fails to predict cross\ldots \\
\textbf{\eqlab{eq:\discretionary{}{}{}8.\discretionary{}{}{}18-\discretionary{}{}{}tweedism-\discretionary{}{}{}agenda-\discretionary{}{}{}path}} & Ch.~8, l.~4844 & \textcolor{green!60!black}{\textbf{T1}} & Gilens \& Page (2014) --- Testing Theories of American\ldots & Falsified if Gilens-Page analysis shows median voter preferences predic\ldots \\
\multicolumn{5}{l}{\small\bfseries Chapter 9: The Recompile: COINTELPRO, the Variable Swap, and the War on Drugs (1968--1994)} \\
\textbf{\eqlab{eq:\discretionary{}{}{}9.\discretionary{}{}{}1-\discretionary{}{}{}lead-\discretionary{}{}{}crime-\discretionary{}{}{}compounding}} & Ch.~9, l.~6344 & \textcolor{green!60!black}{\textbf{T1}} & Reyes (2007) --- Environmental Policy as Social Policy & Falsified if Reyes (2007) blood-lead time-series fails to predict crime\ldots \\
\eqlab{eq:\discretionary{}{}{}9.\discretionary{}{}{}2-\discretionary{}{}{}epistemic-\discretionary{}{}{}suppression-\discretionary{}{}{}variable} & Ch.~9, l.~6344 & \textcolor{blue!70!black}{\textbf{T2}} & Kitman (2000) --- The Secret History of Lead (The Nat\ldots & Falsified if documented lead industry cover-up timeline shows no correl\ldots \\
\eqlab{eq:\discretionary{}{}{}9.\discretionary{}{}{}3-\discretionary{}{}{}multi-\discretionary{}{}{}vector-\discretionary{}{}{}lead-\discretionary{}{}{}exposure} & Ch.~9, l.~6344 & \textcolor{green!60!black}{\textbf{T1}} & EPA Air Quality Monitoring Data & Falsified if cumulative lead exposure from all three vectors is not sig\ldots \\
\eqlab{eq:\discretionary{}{}{}9.\discretionary{}{}{}4-\discretionary{}{}{}redlining-\discretionary{}{}{}containment-\discretionary{}{}{}field} & Ch.~9, l.~6344 & \textcolor{green!60!black}{\textbf{T1}} & Mapping Inequality HOLC maps & Falsified if HOLC map grades show no statistically significant correlat\ldots \\
\eqlab{eq:\discretionary{}{}{}9.\discretionary{}{}{}5-\discretionary{}{}{}school-\discretionary{}{}{}funding-\discretionary{}{}{}property-\discretionary{}{}{}value} & Ch.~9, l.~6337 & \textcolor{green!60!black}{\textbf{T1}} & NCES Public School Finance Survey & Falsified if school funding is shown to be independent of property valu\ldots \\
\eqlab{eq:\discretionary{}{}{}9.\discretionary{}{}{}6-\discretionary{}{}{}infrastructure-\discretionary{}{}{}quality-\discretionary{}{}{}funding} & Ch.~9, l.~6337 & \textcolor{blue!70!black}{\textbf{T2}} & NCES School Facility Condition Survey & Falsified if school infrastructure quality is shown to be unrelated to \ldots \\
\eqlab{eq:\discretionary{}{}{}9.\discretionary{}{}{}7-\discretionary{}{}{}school-\discretionary{}{}{}lead-\discretionary{}{}{}exposure-\discretionary{}{}{}inverse} & Ch.~9, l.~6337 & \textcolor{green!60!black}{\textbf{T1}} & EPA 3Ts for Reducing Lead in Drinking Water in Scho\ldots & Falsified if school infrastructure quality is shown to be unrelated to \ldots \\
\eqlab{eq:\discretionary{}{}{}9.\discretionary{}{}{}8-\discretionary{}{}{}community-\discretionary{}{}{}capacity-\discretionary{}{}{}lead-\discretionary{}{}{}reduction} & Ch.~9, l.~6337 & \textcolor{blue!70!black}{\textbf{T2}} & Aizer \& Currie (2019) --- Lead and Juvenile Delinquen\ldots & Falsified if community economic capacity is shown to be independent of \ldots \\
\eqlab{eq:\discretionary{}{}{}9.\discretionary{}{}{}9-\discretionary{}{}{}property-\discretionary{}{}{}value-\discretionary{}{}{}capacity-\discretionary{}{}{}feedback} & Ch.~9, l.~6337 & \textcolor{green!60!black}{\textbf{T1}} & Mapping Inequality HOLC maps + ACS property value d\ldots & Falsified if property values in subsequent periods are shown to be inde\ldots \\
\eqlab{eq:\discretionary{}{}{}9.\discretionary{}{}{}10-\discretionary{}{}{}highway-\discretionary{}{}{}lead-\discretionary{}{}{}spatial-\discretionary{}{}{}concentration} & Ch.~9, l.~7178 & \textcolor{green!60!black}{\textbf{T1}} & Mapping Inequality DSL (Nelson et al.) --- HOLC grade\ldots & Falsified if HOLC-D tracts show no excess firearm or lead burden relati\ldots \\
\multicolumn{5}{l}{\small\bfseries Chapter 10: The Full Algorithm: Demographic Paradox, Cannibalization, and the 5-Tier Reveal (1994--Present)} \\
\eqlab{eq:\discretionary{}{}{}10.\discretionary{}{}{}1-\discretionary{}{}{}outgroup-\discretionary{}{}{}racialized-\discretionary{}{}{}subset} & Ch.~10, l.~7051 & \textcolor{gray}{\textbf{T3}} & US Census Bureau demographic data & Falsified if O\_racialized is shown to be coextensive with O rather than\ldots \\
\eqlab{eq:\discretionary{}{}{}10.\discretionary{}{}{}2-\discretionary{}{}{}buffer-\discretionary{}{}{}class-\discretionary{}{}{}shrinkage} & Ch.~10, l.~7055 & \textcolor{blue!70!black}{\textbf{T2}} & Pew Research Center --- America's Shrinking Middle Cl\ldots & Falsified if Buffer Class size is shown to be expanding rather than con\ldots \\
\eqlab{eq:\discretionary{}{}{}10.\discretionary{}{}{}3-\discretionary{}{}{}cannibalization-\discretionary{}{}{}equation} & Ch.~10, l.~7153 & \textcolor{gray}{\textbf{T3}} & Saez-Zucman wealth share series & Falsified if documented Elite extraction rate increases without corresp\ldots \\
\eqlab{eq:\discretionary{}{}{}10.\discretionary{}{}{}4-\discretionary{}{}{}predatory-\discretionary{}{}{}minmax-\discretionary{}{}{}full-\discretionary{}{}{}definition} & Ch.~10, l.~7303 & \textcolor{gray}{\textbf{T3}} & \textit{(ordinal)} & Falsified if the five-tier hierarchy fails to map onto any contemporary\ldots \\
\eqlab{eq:\discretionary{}{}{}10.\discretionary{}{}{}5-\discretionary{}{}{}kernel-\discretionary{}{}{}objective-\discretionary{}{}{}canonical} & Ch.~10, l.~7477 & \textcolor{green!60!black}{\textbf{T1}} & Piketty-Saez-Zucman Distributional National Accounts & Falsified if a documented period shows sustained Elite wealth-share dec\ldots \\
\eqlab{eq:\discretionary{}{}{}10.\discretionary{}{}{}6-\discretionary{}{}{}effective-\discretionary{}{}{}resistance-\discretionary{}{}{}variable} & Ch.~10, l.~7309 & \textcolor{gray}{\textbf{T3}} & \textit{(ordinal)} & Falsified if M\_eff fails to capture the difference between raw class co\ldots \\
\eqlab{eq:\discretionary{}{}{}10.\discretionary{}{}{}7-\discretionary{}{}{}suppression-\discretionary{}{}{}envelope-\discretionary{}{}{}canonical} & Ch.~10, l.~7345 & \textcolor{blue!70!black}{\textbf{T2}} & \textit{(pending)} & Falsified if documented suppression tools outside this set constitute t\ldots \\
\eqlab{eq:\discretionary{}{}{}10.\discretionary{}{}{}8-\discretionary{}{}{}crash-\discretionary{}{}{}condition-\discretionary{}{}{}canonical} & Ch.~10, l.~7319 & \textcolor{gray}{\textbf{T3}} & \textit{(ordinal)} & Falsified if crash events occur in documented cases without M\_eff(t) > \ldots \\
\eqlab{eq:\discretionary{}{}{}10.\discretionary{}{}{}9-\discretionary{}{}{}tier-\discretionary{}{}{}benefit-\discretionary{}{}{}ordering-\discretionary{}{}{}canonical} & Ch.~10, l.~7324 & \textcolor{gray}{\textbf{T3}} & \textit{(ordinal)} & Falsified if any two-tier comparison fails to show stated ordering in c\ldots \\
\eqlab{eq:\discretionary{}{}{}10.\discretionary{}{}{}10-\discretionary{}{}{}system-\discretionary{}{}{}objective-\discretionary{}{}{}extraction-\discretionary{}{}{}ratio} & Ch.~10, l.~7405 & \textcolor{gray}{\textbf{T3}} & \textit{(ordinal)} & Falsified if documented extraction systems show this ratio declining wi\ldots \\
\eqlab{eq:\discretionary{}{}{}10.\discretionary{}{}{}11-\discretionary{}{}{}system-\discretionary{}{}{}stability-\discretionary{}{}{}algorithm} & Ch.~10, l.~7405 & \textcolor{gray}{\textbf{T3}} & \textit{(ordinal)} & Falsified if system stability is maintained when extraction/resistance \ldots \\
\textbf{\eqlab{eq:\discretionary{}{}{}10.\discretionary{}{}{}12-\discretionary{}{}{}o-\discretionary{}{}{}final-\discretionary{}{}{}construction}} & Ch.~10, l.~6828 & \textcolor{green!60!black}{\textbf{T1}} & Bureau of Justice Statistics incarceration by race \ldots & Falsified if modern mass-incarceration demographics fail to reproduce t\ldots \\
\eqlab{eq:\discretionary{}{}{}10.\discretionary{}{}{}13-\discretionary{}{}{}reclassification-\discretionary{}{}{}operator} & Ch.~10, l.~7734 & \textcolor{gray}{\textbf{T3}} & \textit{(ordinal)} & Falsified if the operator fails to predict observed ideological sorting\ldots \\
\eqlab{eq:\discretionary{}{}{}10.\discretionary{}{}{}14-\discretionary{}{}{}kinetic-\discretionary{}{}{}power-\discretionary{}{}{}distribution} & Ch.~10, l.~7768 & \textcolor{blue!70!black}{\textbf{T2}} & Pew Research Center --- Gun Ownership Survey 2017 & Falsified if kinetic power is not concentrated in F\_enforce as a propor\ldots \\
\eqlab{eq:\discretionary{}{}{}10.\discretionary{}{}{}15-\discretionary{}{}{}temporal-\discretionary{}{}{}enclosure-\discretionary{}{}{}definition} & Ch.~10, l.~7820 & \textcolor{gray}{\textbf{T3}} & Reinhart-Sbrancia 2015 (Economic Policy) & Falsified if sovereign debt accumulation is demonstrably unrelated to p\ldots \\
\textbf{\eqlab{eq:\discretionary{}{}{}10.\discretionary{}{}{}16-\discretionary{}{}{}financial-\discretionary{}{}{}repression-\discretionary{}{}{}condition}} & Ch.~10, l.~7825 & \textcolor{green!60!black}{\textbf{T1}} & FRED 10Y Treasury Constant Maturity Rate (GS10) & Falsified if: (1) real interest rates remain positive throughout the 19\ldots \\
\eqlab{eq:\discretionary{}{}{}10.\discretionary{}{}{}17-\discretionary{}{}{}temporal-\discretionary{}{}{}extraction-\discretionary{}{}{}rate} & Ch.~10, l.~7825 & \textcolor{green!60!black}{\textbf{T1}} & FRED: Federal Debt Total Public Debt as \% of GDP (G\ldots & Falsified if the indicator function 1[r < g] is inactive during documen\ldots \\
\textbf{\eqlab{eq:\discretionary{}{}{}10.\discretionary{}{}{}18-\discretionary{}{}{}psi-\discretionary{}{}{}degradation-\discretionary{}{}{}function}} & Ch.~10, l.~7905 & \textcolor{green!60!black}{\textbf{T1}} & FRED: Real Median Hourly Compensation (B4701C0A052N\ldots & Falsified if: (1) Chow test at 1971 shows no statistically significant \ldots \\
\eqlab{eq:\discretionary{}{}{}10.\discretionary{}{}{}19-\discretionary{}{}{}saeculum-\discretionary{}{}{}phase-\discretionary{}{}{}map} & Ch.~10, l.~7900 & \textcolor{blue!70!black}{\textbf{T2}} & Strauss \& Howe, The Fourth Turning (1997) & Falsified if: (1) the generational archetype sequence does not follow t\ldots \\
\textbf{\eqlab{eq:\discretionary{}{}{}10.\discretionary{}{}{}20-\discretionary{}{}{}dalio-\discretionary{}{}{}extraction-\discretionary{}{}{}bound}} & Ch.~10, l.~7950 & \textcolor{blue!70!black}{\textbf{T2}} & Maddison Project Database 2020 (GDP, trade output) & Falsified if: (1) the reconstructed US Empire Index does not show a pea\ldots \\
\multicolumn{5}{l}{\small\bfseries Chapter 11: The Kinetic Guarantee: Arms Asymmetry, the Second Amendment, and the Disarmament Timeline} \\
\textbf{\eqlab{eq:\discretionary{}{}{}11.\discretionary{}{}{}1-\discretionary{}{}{}extraction-\discretionary{}{}{}precondition}} & Ch.~11, l.~7992 & \textcolor{green!60!black}{\textbf{T1}} & Primary legislation texts (sullivanlaw, nfa1934, mu\ldots & Falsified if a documented extraction system maintained stability with L\ldots \\
\eqlab{eq:\discretionary{}{}{}11.\discretionary{}{}{}2-\discretionary{}{}{}disarmament-\discretionary{}{}{}agenda-\discretionary{}{}{}path} & Ch.~11, l.~7992 & \textcolor{green!60!black}{\textbf{T1}} & Primary legislation and ruling texts (11 events) & Falsified if documented gun-control legislation passed under high $\Phi$\_loa\ldots \\
\eqlab{eq:\discretionary{}{}{}11.\discretionary{}{}{}3-\discretionary{}{}{}restriction-\discretionary{}{}{}precedent-\discretionary{}{}{}accumulation} & Ch.~11, l.~7992 & \textcolor{green!60!black}{\textbf{T1}} & Primary legislation and ruling texts (11 events) & Falsified if any documented legislative or ruling event reduced cumulat\ldots \\
\eqlab{eq:\discretionary{}{}{}11.\discretionary{}{}{}4-\discretionary{}{}{}restriction-\discretionary{}{}{}scope-\discretionary{}{}{}expansion} & Ch.~11, l.~7992 & \textcolor{green!60!black}{\textbf{T1}} & Primary legislation texts and enforcement records & Falsified if the historical record showed restriction scope declining o\ldots \\
\multicolumn{5}{l}{\small\bfseries Chapter 12: The Contradiction: Why Reform Serves the Algorithm} \\
\eqlab{eq:\discretionary{}{}{}12.\discretionary{}{}{}1-\discretionary{}{}{}reform-\discretionary{}{}{}absorption-\discretionary{}{}{}mechanism} & Ch.~12, l.~8859 & \textcolor{gray}{\textbf{T3}} & Piketty-Saez top income share data & Falsified if a documented reform reduces Elite extraction share (max) r\ldots \\
\eqlab{eq:\discretionary{}{}{}12.\discretionary{}{}{}2-\discretionary{}{}{}concession-\discretionary{}{}{}theorem} & Ch.~12, l.~8920 & \textcolor{blue!70!black}{\textbf{T2}} & Piven \& Cloward (1977) --- Poor People's Movements & Falsified if Elite-initiated reforms consistently exceed the minimum re\ldots \\
\eqlab{eq:\discretionary{}{}{}12.\discretionary{}{}{}3-\discretionary{}{}{}judiciary-\discretionary{}{}{}discrimination-\discretionary{}{}{}detection} & Ch.~12, l.~9071 & \textcolor{gray}{\textbf{T3}} & Washington v. Davis & Falsified if judicial detection function produces false positives for p\ldots \\
\eqlab{eq:\discretionary{}{}{}12.\discretionary{}{}{}4-\discretionary{}{}{}proxy-\discretionary{}{}{}discrimination-\discretionary{}{}{}equivalence} & Ch.~12, l.~9111 & \textcolor{gray}{\textbf{T3}} & Alexander (2010) --- The New Jim Crow & Falsified if P\_proxy achieves racial separation at rates significantly \ldots \\
\eqlab{eq:\discretionary{}{}{}12.\discretionary{}{}{}4a-\discretionary{}{}{}judicial-\discretionary{}{}{}double-\discretionary{}{}{}agent} & Ch.~12, l.~10000 & \textcolor{blue!70!black}{\textbf{T2}} & Local SCOTUS markdown corpus & Falsified if the SCOTUS corpus does not separate civil-rights/voting ca\ldots \\
\textbf{\eqlab{eq:\discretionary{}{}{}12.\discretionary{}{}{}5-\discretionary{}{}{}haitian-\discretionary{}{}{}theorem-\discretionary{}{}{}nonkinetic}} & Ch.~12, l.~8507 & \textcolor{green!60!black}{\textbf{T1}} & James (1938) The Black Jacobins & Falsified if any documented non-kinetic reform achieves $\Delta$max < 0 (susta\ldots \\
\textbf{\eqlab{eq:\discretionary{}{}{}12.\discretionary{}{}{}6-\discretionary{}{}{}haitian-\discretionary{}{}{}theorem-\discretionary{}{}{}kinetic}} & Ch.~12, l.~8507 & \textcolor{green!60!black}{\textbf{T1}} & James (1938) The Black Jacobins & Falsified if any kinetic event fails to produce at least temporary loca\ldots \\
\eqlab{eq:\discretionary{}{}{}12.\discretionary{}{}{}7-\discretionary{}{}{}gendered-\discretionary{}{}{}outgroup-\discretionary{}{}{}set} & Ch.~12, l.~9597 & \textcolor{gray}{\textbf{T3}} & \textit{(ordinal)} & Falsified if O\_gendered is shown to be structured differently from O\_ra\ldots \\
\eqlab{eq:\discretionary{}{}{}12.\discretionary{}{}{}8-\discretionary{}{}{}intersection-\discretionary{}{}{}double-\discretionary{}{}{}partition} & Ch.~12, l.~9602 & \textcolor{gray}{\textbf{T3}} & Crenshaw (1989) --- Demarginalizing the Intersection \ldots & Falsified if intersection O\_racialized $\cap$ O\_gendered does not show compo\ldots \\
\eqlab{eq:\discretionary{}{}{}12.\discretionary{}{}{}9-\discretionary{}{}{}multiplicative-\discretionary{}{}{}intersection-\discretionary{}{}{}compounding} & Ch.~12, l.~9607 & \textcolor{blue!70!black}{\textbf{T2}} & AAUW pay-gap data by race and gender & Falsified if Black women's wealth gap from white men is not multiplicat\ldots \\
\eqlab{eq:\discretionary{}{}{}12.\discretionary{}{}{}10-\discretionary{}{}{}nonviolence-\discretionary{}{}{}mandate-\discretionary{}{}{}function} & Ch.~12, l.~9709 & \textcolor{gray}{\textbf{T3}} & Carson (1981) --- In Struggle: SNCC and the Black Awa\ldots & Falsified if documented liberation movements adopting nonviolence achie\ldots \\
\multicolumn{5}{l}{\small\bfseries Chapter 13: The Global Containment Field: Scaling the Algorithm} \\
\eqlab{eq:\discretionary{}{}{}13.\discretionary{}{}{}1-\discretionary{}{}{}global-\discretionary{}{}{}compounding-\discretionary{}{}{}capacity} & Ch.~13, l.~9925 & \textcolor{blue!70!black}{\textbf{T2}} & World Bank development indicators --- GDP per capita \ldots & Falsified if post-colonial nation's capacity trajectory shows exponenti\ldots \\
\eqlab{eq:\discretionary{}{}{}13.\discretionary{}{}{}2-\discretionary{}{}{}global-\discretionary{}{}{}predatory-\discretionary{}{}{}minmax} & Ch.~13, l.~9973 & \textcolor{gray}{\textbf{T3}} & Klein (2007) --- The Shock Doctrine & Falsified if documented international economic relationships fail to ma\ldots \\
\eqlab{eq:\discretionary{}{}{}13.\discretionary{}{}{}3-\discretionary{}{}{}externalized-\discretionary{}{}{}enforcement-\discretionary{}{}{}envelope} & Ch.~13, l.~9993 & \textcolor{gray}{\textbf{T3}} & UN MINUSTAH mission records & Falsified if Haiti's domestic coercive capacity is shown to be structur\ldots \\
\eqlab{eq:\discretionary{}{}{}13.\discretionary{}{}{}4-\discretionary{}{}{}killick-\discretionary{}{}{}extraction-\discretionary{}{}{}equation} & Ch.~13, l.~10021 & \textcolor{gray}{\textbf{T3}} & NYT 'The Ransom' (2022) & Falsified if the extraction equation fails to predict subsequent debt i\ldots \\
\eqlab{eq:\discretionary{}{}{}13.\discretionary{}{}{}5-\discretionary{}{}{}legitimation-\discretionary{}{}{}constraint-\discretionary{}{}{}set} & Ch.~13, l.~10039 & \textcolor{gray}{\textbf{T3}} & \textit{(ordinal)} & Falsified if E\_global actors consistently operate outside their stated \ldots \\
\eqlab{eq:\discretionary{}{}{}13.\discretionary{}{}{}6-\discretionary{}{}{}authorization-\discretionary{}{}{}precondition} & Ch.~13, l.~10050 & \textcolor{gray}{\textbf{T3}} & \textit{(ordinal)} & Falsified if E\_global achieves extraction without satisfying authorizat\ldots \\
\eqlab{eq:\discretionary{}{}{}13.\discretionary{}{}{}7-\discretionary{}{}{}coercion-\discretionary{}{}{}legitimacy-\discretionary{}{}{}frontier} & Ch.~13, l.~10056 & \textcolor{gray}{\textbf{T3}} & \textit{(ordinal)} & Falsified if observed coercion levels in 'voluntary' agreements consist\ldots \\
\eqlab{eq:\discretionary{}{}{}13.\discretionary{}{}{}8-\discretionary{}{}{}firmin-\discretionary{}{}{}extraction-\discretionary{}{}{}threshold} & Ch.~13, l.~10062 & \textcolor{gray}{\textbf{T3}} & \textit{(ordinal)} & Falsified if E\_global achieves extraction with deployed force below F* \ldots \\
\eqlab{eq:\discretionary{}{}{}13.\discretionary{}{}{}9-\discretionary{}{}{}infeasible-\discretionary{}{}{}constraint-\discretionary{}{}{}set} & Ch.~13, l.~10075 & \textcolor{gray}{\textbf{T3}} & \textit{(ordinal)} & Falsified if E\_global achieves both extraction success and legitimation\ldots \\
\eqlab{eq:\discretionary{}{}{}13.\discretionary{}{}{}10-\discretionary{}{}{}asymmetric-\discretionary{}{}{}leverage-\discretionary{}{}{}ratio} & Ch.~13, l.~10100 & \textcolor{blue!70!black}{\textbf{T2}} & Firmin (1885) --- De l'égalité des races humaines & Falsified if a weaker actor's invocation of E\_global's legitimation arc\ldots \\
\eqlab{eq:\discretionary{}{}{}13.\discretionary{}{}{}11-\discretionary{}{}{}legitimation-\discretionary{}{}{}abandonment-\discretionary{}{}{}condition} & Ch.~13, l.~10107 & \textcolor{gray}{\textbf{T3}} & \textit{(ordinal)} & Falsified if E\_global maintains legitimation constraints even when extr\ldots \\
\eqlab{eq:\discretionary{}{}{}13.\discretionary{}{}{}12-\discretionary{}{}{}global-\discretionary{}{}{}containment-\discretionary{}{}{}strategy} & Ch.~13, l.~10123 & \textcolor{gray}{\textbf{T3}} & \textit{(ordinal)} & Falsified if global containment strategies differ structurally from the\ldots \\
\eqlab{eq:\discretionary{}{}{}13.\discretionary{}{}{}13-\discretionary{}{}{}un-\discretionary{}{}{}vote-\discretionary{}{}{}distribution} & Ch.~13, l.~10162 & \textcolor{gray}{\textbf{T3}} & UN Digital Library --- UNSC vote records & Falsified if UN Security Council voting distribution fails to map onto \ldots \\
\textbf{\eqlab{eq:\discretionary{}{}{}13.\discretionary{}{}{}14-\discretionary{}{}{}imperial-\discretionary{}{}{}core-\discretionary{}{}{}collapse}} & Ch.~13, l.~9208 & \textcolor{green!60!black}{\textbf{T1}} & World Bank Development Indicators --- GDP trajectory \ldots & Falsified if documented rising powers (China, OPEC, Asian Tigers) achie\ldots \\
\textbf{\eqlab{eq:\discretionary{}{}{}13.\discretionary{}{}{}15-\discretionary{}{}{}interface-\discretionary{}{}{}swap-\discretionary{}{}{}trigger}} & Ch.~13, l.~10280 & \textcolor{green!60!black}{\textbf{T1}} & Federal Reserve H.4.1: Historical Gold Stock Data (\ldots & Falsified if the gold-cover ratio was still above 80\% at the time of th\ldots \\
\eqlab{eq:\discretionary{}{}{}13.\discretionary{}{}{}16-\discretionary{}{}{}polymorphic-\discretionary{}{}{}reboot-\discretionary{}{}{}operator} & Ch.~13, l.~10310 & \textcolor{blue!70!black}{\textbf{T2}} & Gowa 1983 (Closing the Gold Window) & Falsified if Assets(E, A\_new) < Assets(E, A\_old) across any of the thre\ldots \\
\textbf{\eqlab{eq:\discretionary{}{}{}13.\discretionary{}{}{}17-\discretionary{}{}{}elite-\discretionary{}{}{}asset-\discretionary{}{}{}continuity-\discretionary{}{}{}invariant}} & Ch.~13, l.~10320 & \textcolor{green!60!black}{\textbf{T1}} & WID.world: Top-0.1\% US Wealth Share 1913-present (P\ldots & Falsified if the top-0.1\% US wealth share declines STRUCTURALLY (> 3 pe\ldots \\
\multicolumn{5}{l}{\small\bfseries Chapter 14: The Algorithmic Epoch: Real-Time Subjugation and the Necessity of the Counter-Virus} \\
\eqlab{eq:\discretionary{}{}{}14.\discretionary{}{}{}1-\discretionary{}{}{}algo-\discretionary{}{}{}prior-\discretionary{}{}{}inheritance} & Ch.~14, l.~10757 & \textcolor{blue!70!black}{\textbf{T2}} & Angwin et al. (2016) --- Machine Bias (ProPublica) & Falsified if ML systems trained on pre-2024 US data show no racial disp\ldots \\
\eqlab{eq:\discretionary{}{}{}14.\discretionary{}{}{}2-\discretionary{}{}{}realtime-\discretionary{}{}{}proximity-\discretionary{}{}{}estimate} & Ch.~14, l.~10777 & \textcolor{blue!70!black}{\textbf{T2}} & Twitter/X political affinity graph data (academic d\ldots & Falsified if real-time social-affinity graph G(t) shows no correlation \ldots \\
\eqlab{eq:\discretionary{}{}{}14.\discretionary{}{}{}3-\discretionary{}{}{}orthogonal-\discretionary{}{}{}vector-\discretionary{}{}{}injection} & Ch.~14, l.~10793 & \textcolor{gray}{\textbf{T3}} & \textit{(ordinal)} & Falsified if documented interference engine operations increase class-a\ldots \\
\eqlab{eq:\discretionary{}{}{}14.\discretionary{}{}{}4-\discretionary{}{}{}correction-\discretionary{}{}{}signal-\discretionary{}{}{}damping} & Ch.~14, l.~10807 & \textcolor{gray}{\textbf{T3}} & \textit{(ordinal)} & Falsified if scale-accurate counter-signals are shown to reduce identit\ldots \\
\eqlab{eq:\discretionary{}{}{}14.\discretionary{}{}{}5-\discretionary{}{}{}optical-\discretionary{}{}{}enclosure} & Ch.~14, l.~10820 & \textcolor{gray}{\textbf{T3}} & Pew Research --- Police-community relations by race & Falsified if documented perception studies show O's population maintain\ldots \\
\eqlab{eq:\discretionary{}{}{}14.\discretionary{}{}{}6-\discretionary{}{}{}decoy-\discretionary{}{}{}transfer-\discretionary{}{}{}coefficient} & Ch.~14, l.~10834 & \textcolor{blue!70!black}{\textbf{T2}} & \textit{(pending)} & Falsified if documented cases show kinetic class energy passing to E at\ldots \\
\eqlab{eq:\discretionary{}{}{}14.\discretionary{}{}{}7-\discretionary{}{}{}ddos-\discretionary{}{}{}bandwidth-\discretionary{}{}{}ceiling} & Ch.~14, l.~10858 & \textcolor{blue!70!black}{\textbf{T2}} & Armed Conflict Location \& Event Data (ACLED) --- US p\ldots & Falsified if F\_enforce simultaneously suppresses n > $\Gamma$ nodes in any doc\ldots \\
\eqlab{eq:\discretionary{}{}{}14.\discretionary{}{}{}8-\discretionary{}{}{}armed-\discretionary{}{}{}swarm-\discretionary{}{}{}bandwidth} & Ch.~14, l.~10869 & \textcolor{blue!70!black}{\textbf{T2}} & \textit{(pending)} & Falsified if F\_enforce simultaneously neutralizes armed nodes with forc\ldots \\
\eqlab{eq:\discretionary{}{}{}14.\discretionary{}{}{}9-\discretionary{}{}{}kinetic-\discretionary{}{}{}capital-\discretionary{}{}{}asymmetry} & Ch.~14, l.~10887 & \textcolor{blue!70!black}{\textbf{T2}} & Pew Research Center --- Gun Ownership Survey 2017 & Falsified if gun-ownership data shows no racial asymmetry in kinetic ca\ldots \\
\eqlab{eq:\discretionary{}{}{}14.\discretionary{}{}{}10-\discretionary{}{}{}sustained-\discretionary{}{}{}siege-\discretionary{}{}{}cost} & Ch.~14, l.~10909 & \textcolor{blue!70!black}{\textbf{T2}} & \textit{(pending)} & Falsified if F\_enforce besieges multiple simultaneous armed nodes with \ldots \\
\eqlab{eq:\discretionary{}{}{}14.\discretionary{}{}{}11-\discretionary{}{}{}escalation-\discretionary{}{}{}ceiling-\discretionary{}{}{}tiers} & Ch.~14, l.~10919 & \textcolor{gray}{\textbf{T3}} & \textit{(ordinal)} & Falsified if F\_enforce escalation tiers do not follow the predicted ban\ldots \\
\eqlab{eq:\discretionary{}{}{}14.\discretionary{}{}{}12-\discretionary{}{}{}inclusion-\discretionary{}{}{}predicate} & Ch.~14, l.~10948 & \textcolor{blue!70!black}{\textbf{T2}} & Bureau of Labor Statistics contingent worker survey & Falsified if data shows V(x,t) $\leq$ R(x,t) for majority members of I\_buffe\ldots \\
\eqlab{eq:\discretionary{}{}{}14.\discretionary{}{}{}13-\discretionary{}{}{}terminal-\discretionary{}{}{}interface-\discretionary{}{}{}swap} & Ch.~14, l.~10959 & \textcolor{gray}{\textbf{T3}} & \textit{(ordinal)} & Falsified if documented late-stage extraction shows I\_buffer benefits m\ldots \\
\eqlab{eq:\discretionary{}{}{}14.\discretionary{}{}{}14-\discretionary{}{}{}liability-\discretionary{}{}{}phase-\discretionary{}{}{}optimization} & Ch.~14, l.~10969 & \textcolor{gray}{\textbf{T3}} & \textit{(ordinal)} & Falsified if documented terminal-phase extraction shows E maintaining I\ldots \\
\eqlab{eq:\discretionary{}{}{}14.\discretionary{}{}{}15-\discretionary{}{}{}noise-\discretionary{}{}{}spectrum-\discretionary{}{}{}index} & Ch.~14, l.~11017 & \textcolor{blue!70!black}{\textbf{T2}} & ACLED --- US social movements 2011--2020 & Falsified if noise-spectrum index of a documented swarm fails to predic\ldots \\
\eqlab{eq:\discretionary{}{}{}14.\discretionary{}{}{}16-\discretionary{}{}{}agnostic-\discretionary{}{}{}swarm-\discretionary{}{}{}payload} & Ch.~14, l.~11036 & \textcolor{gray}{\textbf{T3}} & \textit{(ordinal)} & Falsified if documented multi-node mobilizations with ideological diver\ldots \\
\eqlab{eq:\discretionary{}{}{}14.\discretionary{}{}{}17-\discretionary{}{}{}botnet-\discretionary{}{}{}load-\discretionary{}{}{}theorem} & Ch.~14, l.~11051 & \textcolor{blue!70!black}{\textbf{T2}} & ACLED --- 2020 US protest dataset & Falsified if enforcement grid contains a swarm with n > $\Gamma$\_armed nodes a\ldots \\
\eqlab{eq:\discretionary{}{}{}14.\discretionary{}{}{}18-\discretionary{}{}{}ideological-\discretionary{}{}{}routing-\discretionary{}{}{}operator} & Ch.~14, l.~11072 & \textcolor{gray}{\textbf{T3}} & \textit{(ordinal)} & Falsified if nodes in documented distributed mobilizations fail to rout\ldots \\
\eqlab{eq:\discretionary{}{}{}14.\discretionary{}{}{}19-\discretionary{}{}{}zugzwang-\discretionary{}{}{}payoff-\discretionary{}{}{}function} & Ch.~14, l.~11122 & \textcolor{blue!70!black}{\textbf{T2}} & \textit{(pending)} & Falsified if Elite actors systematically choose attack over stand-down \ldots \\
\eqlab{eq:\discretionary{}{}{}14.\discretionary{}{}{}20-\discretionary{}{}{}defection-\discretionary{}{}{}cascade} & Ch.~14, l.~11164 & \textcolor{blue!70!black}{\textbf{T2}} & James (1938) --- The Black Jacobins (Polish Legion de\ldots & Falsified if documented enforcement nodes fail to defect when community\ldots \\
\eqlab{eq:\discretionary{}{}{}14.\discretionary{}{}{}21-\discretionary{}{}{}empathy-\discretionary{}{}{}permeability} & Ch.~14, l.~11196 & \textcolor{blue!70!black}{\textbf{T2}} & James (1938) --- The Black Jacobins & Falsified if trauma overlap between enforcement and out-group nodes fai\ldots \\
\eqlab{eq:\discretionary{}{}{}14.\discretionary{}{}{}22-\discretionary{}{}{}semantic-\discretionary{}{}{}overwrite} & Ch.~14, l.~11222 & \textcolor{gray}{\textbf{T3}} & Haitian Constitution of 1805 (primary source) & Falsified if 1805 Haitian Constitution's ideological redefinition faile\ldots \\
\multicolumn{5}{l}{\small\bfseries Chapter 17: Conclusion} \\
\eqlab{eq:\discretionary{}{}{}17.\discretionary{}{}{}1-\discretionary{}{}{}racism-\discretionary{}{}{}as-\discretionary{}{}{}institutional-\discretionary{}{}{}vector} & Ch.~17, l.~11531 & \textcolor{gray}{\textbf{T3}} & \textit{(ordinal)} & Falsified if documented cases of marginalized-group prejudice achieve s\ldots \\
\eqlab{eq:\discretionary{}{}{}17.\discretionary{}{}{}2-\discretionary{}{}{}o-\discretionary{}{}{}terminal-\discretionary{}{}{}expansion} & Ch.~17, l.~11586 & \textcolor{blue!70!black}{\textbf{T2}} & US Census Bureau 2020 demographic projections & Falsified if documented demographic trends show O contracting rather th\ldots \\
\multicolumn{5}{l}{\small\bfseries Chapter 19: Equation Registry and Era-Level Calibration} \\
\eqlab{eq:\discretionary{}{}{}19.\discretionary{}{}{}1-\discretionary{}{}{}kernel-\discretionary{}{}{}objective-\discretionary{}{}{}registry-\discretionary{}{}{}entry} & Ch.~19, l.~11717 & \textcolor{gray}{\textbf{T3}} & \textit{(ordinal)} & Redundant with eq:5 and eq:56; falsified by the same conditions as thos\ldots \\
\eqlab{eq:\discretionary{}{}{}19.\discretionary{}{}{}2-\discretionary{}{}{}suppression-\discretionary{}{}{}envelope-\discretionary{}{}{}registry-\discretionary{}{}{}entry} & Ch.~19, l.~11721 & \textcolor{gray}{\textbf{T3}} & \textit{(ordinal)} & Redundant with eq:8 and eq:58; falsified by the same conditions as thos\ldots \\
\eqlab{eq:\discretionary{}{}{}19.\discretionary{}{}{}3-\discretionary{}{}{}crash-\discretionary{}{}{}condition-\discretionary{}{}{}registry-\discretionary{}{}{}entry} & Ch.~19, l.~11725 & \textcolor{gray}{\textbf{T3}} & \textit{(ordinal)} & Redundant with eq:9 and eq:59; falsified by the same conditions as thos\ldots \\
\eqlab{eq:\discretionary{}{}{}19.\discretionary{}{}{}4-\discretionary{}{}{}tier-\discretionary{}{}{}ordering-\discretionary{}{}{}registry-\discretionary{}{}{}entry} & Ch.~19, l.~11729 & \textcolor{gray}{\textbf{T3}} & \textit{(ordinal)} & Redundant with eq:60; falsified if any tier ordering is reversed in com\ldots \\
\end{longtable}
\endgroup


\chapter{The Photon Model of Polarizing Information}
\label{app:photon_model}

The electrodynamic framework developed in Chapters~\ref{ch:redefining}--\ref{ch:algorithmic_epoch} treats the Interference Engine as a continuous field: a superposition of cultural magnetic fields $B_k$ with time-varying demographic weights $\rho_k(t)$. This appendix extends the field model into the quantized regime by treating discrete information packets as photon-like carriers of polarization.

The extension is speculative. It is presented as a formal analogy with predictive value, not as a claim that social media posts are literal electromagnetic quanta.

\section{From Continuous Field to Discrete Quanta}

In the continuous model, the Engine injects power into the $k$-th identity axis at rate:
\begin{equation}\label{eq:app-photon-field-power}
P_k = \frac{1}{2\mu_0} |B_k|^2 \, V
\end{equation}
where $V$ is the effective volume of the population exposed to the field. In the analog regime, $V$ was bounded by physical reach; in the digital regime, $V$ approaches the entire network.

The continuous model predicts total power; the photon model quantizes the distribution of that power across individual exposures. Define a \textbf{polarizing photon} as a discrete information packet (tweet, meme, short-form video, algorithmic recommendation) that carries a specific polarization state along one or more identity axes.

\section{The Photon Ansatz}

A polarizing photon has the following properties:

\begin{enumerate}
    \item \textbf{Energy}: $E_{\gamma} = h \nu$, where $\nu$ is the refresh/repetition frequency of the axis and $h$ is the attention-economy constant. High-frequency axes (e.g., transphobia cycles in 2023) have higher photon energy than low-frequency axes (e.g., racial housing covenants, which operated on decadal cycles).
    \item \textbf{Polarization}: Each photon carries a spin state $\sigma_k \in \{-1, 0, +1\}$ along axis $k$. The spin determines whether the packet reinforces Buffer Class alignment ($+1$), Out-group extraction ($-1$), or neutral transmission ($0$).
    \item \textbf{Propagation velocity}: $c_{\text{effective}}$, the network latency. In practice, this is milliseconds for domestic propagation and seconds for global propagation---effectively instantaneous at human perception scales.
    \item \textbf{No rest mass}: The photon exists only in motion. Once shared, it propagates until absorbed (viewed, processed, emotionally reacted to) or decohered (drowned by competing signals).
\end{enumerate}

\section{Spectral Density and the Small-Demographic Anomaly}

The photon model explains the small-demographic targeting puzzle (Section~VI of the electrodynamic synthesis) with precision. The total power $P_k$ injected into axis $k$ can be written as:
\begin{equation}\label{eq:app-photon-spectral-density}
P_k = N_{\gamma} \cdot E_{\gamma} = N_{\gamma} \cdot h \nu_k
\end{equation}
where $N_{\gamma}$ is the photon number (packet count) and $\nu_k$ is the axis frequency.

For a major axis like race ($B_1$), the population base is large. To avoid saturating the Buffer Class (burnout / legitimacy collapse), the Engine must distribute the same power across many low-energy photons: high $N_{\gamma}$, low $\nu_k$. The field is spatially diffuse.

For a minor axis like transphobia ($B_3$), the population base is small ($\sim$0.6\%). The Engine can concentrate the same power into fewer photons of higher individual energy: low $N_{\gamma}$, high $\nu_k$. The Buffer Class has been pre-conditioned by patriarchal $B_{\text{gender}}$ to exhibit high magnetic susceptibility $\chi_m$ on this axis, so the high-energy photons achieve the same deflection at lower total packet count.

The result is a \textbf{spectral optimization}: the Engine selects the axis frequency $\nu_k$ and photon number $N_{\gamma}$ that maximize deflection per unit cost:
\begin{equation}\label{eq:app-photon-optimization}
\max_{k, N_{\gamma}} \; \chi_{m,k} \cdot N_{\gamma} \cdot h \nu_k \quad \text{subject to} \quad C_{\text{engine}}^{\text{digital}}(N_{\gamma}, \nu_k) \leq C_{\text{budget}}
\end{equation}

This is why a demographic of 0.6\% can command a field energy comparable to a demographic of 13\%. The energy is distributed according to the spectral efficiency of deflection, decoupling field share from population proportion.

\section{Implications for Counter-Engineering}

If the interference is carried by photon-like packets, then resistance can be modeled as \textbf{polarization filtering}. The Counter-AI (Chapter~\ref{ch:post_kinetic}) must function as a polarizing filter that selectively attenuates high-$\nu_k$ photons while amplifying class-band photons ($f_{\text{class}}$).

The analogy suggests specific engineering requirements:
\begin{itemize}
    \item \textbf{Frequency-selective attenuation}: Detect and suppress packets whose $\nu_k$ exceeds the natural frequency of organic political discourse. Machine-generated high-frequency polarization (botnets, synthetic media) should be filtered before human exposure.
    \item \textbf{Polarization rotation}: Convert orthogonal deflections back into class-aligned momentum. When a packet with $\sigma_k = +1$ (racial resentment) is detected, the Counter-AI publishes the corresponding class-axis packet showing how $E$ benefits from the same resentment.
    \item \textbf{Coherent amplification}: Class-band photons must be amplified constructively. The Counter-AI must increase $N_{\gamma}$ at $f_{\text{class}}$ until the class-band power exceeds the identity-band power---reversing the spectral redistribution documented in Chapter~\ref{ch:spectral_carrier}.
\end{itemize}

\chapter{Universality and the Finite Topology of Power}
\label{app:universality}

The electrodynamic isomorphism developed in this book maps social oppression onto circuit theory, thermodynamics, and control systems. The question raised in the engineering conversation is whether this isomorphism is a useful analogy or whether it reflects a deeper structural constraint: \textbf{are there only finitely many ways to organize complex systems of power and control?}

This appendix frames that question as a falsifiable conjecture.

\section{The Universality Class Conjecture}

In statistical mechanics, a \textbf{universality class} is a set of physical systems that share the same critical exponents and scaling behavior despite differing in microscopic detail. The Ising model, a lattice gas, and a binary fluid all belong to the same universality class because they share the same symmetry and dimensionality.

The conjecture is that the Predatory Min-Max Function belongs to a universality class of \textbf{hierarchical far-from-equilibrium extractors}: systems that concentrate energy, matter, or information against a gradient while maintaining stability against internal dissipation.

\begin{conjecture}[Universality of the Extraction Kernel]
\label{conj:universality}
Any system satisfying the following three conditions will converge toward the five-node extraction topology $(E, P_{\text{uppet}}, F_{\text{enforce}}, I_{\text{buffer}}, O_{\text{racialized}})$ and the phase-loading control law $\rho_k(t)$:
\begin{enumerate}
    \item \textbf{Bounded resources}: The system operates under a resource constraint that prevents uniform distribution to all nodes.
    \item \textbf{Positive feedback}: A subset of nodes can amplify its own resource share by altering the rules of distribution.
    \item \textbf{Dissipation threat}: The excluded nodes possess sufficient coherence capacity to threaten system stability if they mobilize.
\end{enumerate}
\end{conjecture}

The conjecture claims that the \textit{topology} of the solution is invariant: a small controlling node, an enforcement intermediary, a co-opted buffer, and an excluded sink. The specific identity axes ($B_k$) are culturally contingent; the circuit topology is fixed across contexts.

\section{Evidence from Control Theory}

The five-node topology is the minimal control structure capable of solving the constrained optimization problem:
\begin{equation}\label{eq:app-universality-objective}
\max_{u} \; \mathcal{E}(t) \quad \text{subject to} \quad M(t) < \tau, \; \dot{q} = f(q, u), \; q(0) = q_0
\end{equation}
where $u$ is the control input and $q$ is the state vector. Control theorists have proven that the optimal solution to such problems under resource constraints and dissipation threats takes the form of a \textbf{hierarchical feedback structure}: a leader node ($E$), a controller ($P_{\text{uppet}}$), an actuator ($F_{\text{enforce}}$), and a partitioned follower set ($I_{\text{buffer}} \cup O_{\text{racialized}}$) \cite{slotine_li, khalil_nonlinear}.

The mathematics requires a \textit{dissipation channel}: a path through which the coherence threat $M(t)$ is converted into harmless thermal noise. The racial partition is one historically contingent instantiation of this channel. The specific choice of dissipation channel---race, caste, religion, nationality---is historically contingent. The requirement of a dissipation channel is invariant across instantiations.

\section{Falsification Criteria}

The conjecture is falsifiable if any of the following empirical conditions are met:
\begin{enumerate}
    \item A complex society with bounded resources, positive feedback, and dissipation threat maintains stable egalitarian distribution for $>5$ generations without hierarchical partition.
    \item A hierarchical extractor society collapses into egalitarian equilibrium without kinetic intervention from the excluded population.
    \item A system with the three conditions converges on a non-hierarchical topology (e.g., pure anarchy, rotating leadership, or stochastic resource allocation) and maintains it stably.
\end{enumerate}

To date, no society in the historical dataset satisfies any of these conditions. This leaves the conjecture consistent with the dataset; it is not a proof.

\section{Implications}

If the conjecture holds, the framework's prescriptions are \textbf{engineering requirements}. The Open-Source Republic must be designed as a \textit{different topology entirely}: one that removes the positive-feedback condition (via hard-caps on extraction) and the dissipation-threat condition (via universal baseline provisioning), thereby changing the universality class of the system.

\backmatter

\printbibliography

\end{document}
